%% sockctrl.tex
%%
%% Copyright (C) 2004-2018 Simone Piccardi.  Permission is granted to
%% copy, distribute and/or modify this document under the terms of the GNU Free
%% Documentation License, Version 1.1 or any later version published by the
%% Free Software Foundation; with the Invariant Sections being "Prefazione",
%% with no Front-Cover Texts, and with no Back-Cover Texts.  A copy of the
%% license is included in the section entitled "GNU Free Documentation
%% License".
%%

\chapter{La gestione dei socket}
\label{cha:sock_generic_management}

Esamineremo in questo capitolo una serie di funzionalità aggiuntive relative
alla gestione dei socket, come la gestione della risoluzione di nomi e
indirizzi, le impostazioni delle varie proprietà ed opzioni relative ai
socket, e le funzioni di controllo che permettono di modificarne il
comportamento.


\section{La risoluzione dei nomi}
\label{sec:sock_name_resolution}

Negli esempi dei capitoli precedenti abbiamo sempre identificato le singole
macchine attraverso indirizzi numerici, sfruttando al più le funzioni di
conversione elementare illustrate in sez.~\ref{sec:sock_addr_func} che
permettono di passare da un indirizzo espresso in forma \textit{dotted
  decimal} ad un numero. Vedremo in questa sezione le funzioni utilizzate per
poter utilizzare dei nomi simbolici al posto dei valori numerici, e viceversa
quelle che permettono di ottenere i nomi simbolici associati ad indirizzi,
porte o altre proprietà del sistema.


\subsection{La struttura del \textit{resolver}}
\label{sec:sock_resolver}

\itindbeg{resolver} La risoluzione dei nomi è associata tradizionalmente al
servizio del \itindex{Domain~Name~Service~(DNS)} \textit{Domain Name Service}
che permette di identificare le macchine su internet invece che per numero IP
attraverso il relativo \textsl{nome a dominio}.\footnote{non staremo ad
  entrare nei dettagli della definizione di cosa è un nome a dominio, dandolo
  per noto, una introduzione alla problematica si trova in \cite{AGL} (cap.~9)
  mentre per una trattazione approfondita di tutte le problematiche relative
  al DNS si può fare riferimento a \cite{DNSbind}.} In realtà per DNS si
intendono spesso i server che forniscono su internet questo servizio, mentre
nel nostro caso affronteremo la problematica dal lato client, di un qualunque
programma che necessita di compiere questa operazione.

\begin{figure}[!htb]
  \centering \includegraphics[width=11cm]{img/resolver}
  \caption{Schema di funzionamento delle funzioni del \textit{resolver}.}
  \label{fig:sock_resolver_schema}
\end{figure}

Inoltre quella fra nomi a dominio e indirizzi IP non è l'unica corrispondenza
possibile fra nomi simbolici e valori numerici, come abbiamo visto anche in
sez.~\ref{sec:sys_user_group} per le corrispondenze fra nomi di utenti e
gruppi e relativi identificatori numerici; per quanto riguarda però tutti i
nomi associati a identificativi o servizi relativi alla rete il servizio di
risoluzione è gestito in maniera unificata da un insieme di funzioni fornite
con le librerie del C, detto appunto \textit{resolver}.

Lo schema di funzionamento del \textit{resolver} è illustrato in
fig.~\ref{fig:sock_resolver_schema}; in sostanza i programmi hanno a
disposizione un insieme di funzioni di libreria con cui chiamano il
\textit{resolver}, indicate con le frecce nere. Ricevuta la richiesta è
quest'ultimo che, sulla base della sua configurazione, esegue le operazioni
necessarie a fornire la risposta, che possono essere la lettura delle
informazioni mantenute nei relativi dei file statici presenti sulla macchina,
una interrogazione ad un DNS (che a sua volta, per il funzionamento del
protocollo, può interrogarne altri) o la richiesta ad altri server per i quali
sia fornito il supporto, come LDAP.\footnote{la sigla LDAP fa riferimento ad
  un protocollo, il \textit{Lightweight Directory Access Protocol}, che
  prevede un meccanismo per la gestione di \textsl{elenchi} di informazioni
  via rete; il contenuto di un elenco può essere assolutamente generico, e
  questo permette il mantenimento dei più vari tipi di informazioni su una
  infrastruttura di questo tipo.}

La configurazione del \textit{resolver} attiene più alla amministrazione di
sistema che alla programmazione, ciò non di meno, prima di trattare le varie
funzioni di librerie utilizzate dai programmi, vale la pena fare una
panoramica generale.  Originariamente la configurazione del \textit{resolver}
riguardava esclusivamente le questioni relative alla gestione dei nomi a
dominio, e prevedeva solo l'utilizzo del DNS e del file statico
\conffile{/etc/hosts}.

Per questo aspetto il file di configurazione principale del sistema è
\conffile{/etc/resolv.conf} che contiene in sostanza l'elenco degli indirizzi
IP dei server DNS da contattare; a questo si affiancava (fino alla \acr{glibc}
2.4) il file \conffile{/etc/host.conf} il cui scopo principale era indicare
l'ordine in cui eseguire la risoluzione dei nomi (se usare prima i valori di
\conffile{/etc/hosts} o quelli del DNS). Tralasciamo i dettagli relativi alle
varie direttive che possono essere usate in questi file, che si trovano nelle
rispettive pagine di manuale.

Con il tempo però è divenuto possibile fornire diversi sostituti per
l'utilizzo delle associazione statiche in \conffile{/etc/hosts}, inoltre oltre
alla risoluzione dei nomi a dominio ci sono anche altri nomi da risolvere,
come quelli che possono essere associati ad una rete (invece che ad una
singola macchina) o ai gruppi di macchine definiti dal servizio
NIS,\footnote{il \textit{Network Information Service} è un servizio, creato da
  Sun, e poi diffuso su tutte le piattaforme unix-like, che permette di
  raggruppare all'interno di una rete (in quelli che appunto vengono chiamati
  \textit{netgroup}) varie macchine, centralizzando i servizi di definizione
  di utenti e gruppi e di autenticazione, oggi è sempre più spesso sostituito
  da LDAP.} o come quelli dei protocolli e dei servizi che sono mantenuti nei
file statici \conffile{/etc/protocols} e \conffile{/etc/services}.  

Molte di queste informazioni non si trovano su un DNS, ma in una rete locale
può essere molto utile centralizzare il mantenimento di alcune di esse su
opportuni server.  Inoltre l'uso di diversi supporti possibili per le stesse
informazioni (ad esempio il nome delle macchine può essere mantenuto sia
tramite \conffile{/etc/hosts}, che con il DNS, che con NIS) comporta il
problema dell'ordine in cui questi vengono interrogati. Con le implementazioni
classiche i vari supporti erano introdotti modificando direttamente le
funzioni di libreria, prevedendo un ordine di interrogazione predefinito e non
modificabile (a meno di una ricompilazione delle librerie stesse).

\itindbeg{Name~Service~Switch~(NSS)} 

Per risolvere questa serie di problemi la risoluzione dei nomi a dominio
eseguità dal \textit{resolver} è stata inclusa all'interno di un meccanismo
generico per la risoluzione di corrispondenze fra nomi ed informazioni ad essi
associate chiamato \textit{Name Service Switch}, cui abbiamo accennato anche in
sez.~\ref{sec:sys_user_group} per quanto riguarda la gestione dei dati
associati a utenti e gruppi. Il sistema è stato introdotto la prima volta
nella libreria standard di Solaris e la \acr{glibc} ha ripreso lo stesso
schema; si tenga presente che questo sistema non esiste per altre librerie
standard come la \acr{libc5} o la \acr{uClib}.

Il \textit{Name Service Switch} (cui spesso si fa riferimento con l'acronimo
NSS) è un sistema di librerie dinamiche che permette di definire in maniera
generica sia i supporti su cui mantenere i dati di corrispondenza fra nomi e
valori numerici, sia l'ordine in cui effettuare le ricerche sui vari supporti
disponibili. Il sistema prevede una serie di possibili classi di
corrispondenza, quelle attualmente definite sono riportate in
tab.~\ref{tab:sys_NSS_classes}.

\begin{table}[htb]
  \footnotesize
  \centering
  \begin{tabular}[c]{|l|p{8cm}|}
    \hline
    \textbf{Classe} & \textbf{Tipo di corrispondenza}\\
    \hline
    \hline
    \texttt{passwd}   & Corrispondenze fra nome dell'utente e relative
                        proprietà (\ids{UID}, gruppo principale, ecc.).\\  
    \texttt{shadow}   & Corrispondenze fra username e password dell'utente
                        (e altre informazioni relative alle password).\\  
    \texttt{group}    & Corrispondenze fra nome del gruppo e proprietà dello 
                        stesso.\\  
    \texttt{aliases}  & Alias per la posta elettronica.\\ 
    \texttt{ethers}   & Corrispondenze fra numero IP e MAC address della
                        scheda di rete.\\ 
    \texttt{hosts}    & Corrispondenze fra nome a dominio e numero IP.\\ 
    \texttt{netgroup} & Corrispondenze fra gruppo di rete e macchine che lo
                        compongono.\\  
    \texttt{networks} & Corrispondenze fra nome di una rete e suo indirizzo
                        IP.\\  
    \texttt{protocols}& Corrispondenze fra nome di un protocollo e relativo
                        numero identificativo.\\ 
    \texttt{rpc}      & Corrispondenze fra nome di un servizio RPC e relativo 
                        numero identificativo.\\ 
    \texttt{publickey}& Chiavi pubbliche e private usate per gli RFC sicuri,
                        utilizzate da NFS e NIS+. \\ 
    \texttt{services} & Corrispondenze fra nome di un servizio e numero di
                        porta. \\ 
    \hline
  \end{tabular}
  \caption{Le diverse classi di corrispondenze definite
    all'interno del \textit{Name Service Switch}.} 
  \label{tab:sys_NSS_classes}
\end{table}

% TODO rivedere meglio la tabella

Il sistema del \textit{Name Service Switch} è controllato dal contenuto del
file \conffile{/etc/nsswitch.conf}; questo contiene una riga di configurazione
per ciascuna di queste classi, che viene inizia col nome di
tab.~\ref{tab:sys_NSS_classes} seguito da un carattere ``\texttt{:}'' e
prosegue con la lista dei \textsl{servizi} su cui le relative informazioni
sono raggiungibili, scritti nell'ordine in cui si vuole siano interrogati.
Pertanto nelle versioni recenti delle librerie è questo file e non
\conffile{/etc/host.conf}  a indicare l'ordine con cui si esegue la
risoluzione dei nomi.

Ogni servizio è specificato a sua volta da un nome, come \texttt{file},
\texttt{dns}, \texttt{db}, ecc.  che identifica la libreria dinamica che
realizza l'interfaccia con esso. Per ciascun servizio se \texttt{NAME} è il
nome utilizzato dentro \conffile{/etc/nsswitch.conf}, dovrà essere presente
(usualmente in \file{/lib}) una libreria \texttt{libnss\_NAME} che ne
implementa le funzioni.

In ogni caso, qualunque sia la modalità con cui ricevono i dati o il supporto
su cui vengono mantenuti, e che si usino o meno funzionalità aggiuntive
fornite dal sistema del \textit{Name Service Switch}, dal punto di vista di un
programma che deve effettuare la risoluzione di un nome a dominio, tutto
quello che conta sono le funzioni classiche che il \textit{resolver} mette a
disposizione (è cura della \acr{glibc} tenere conto della presenza del
\textit{Name Service Switch}) e sono queste quelle che tratteremo nelle
sezioni successive.  

\itindend{Name~Service~Switch~(NSS)}


\subsection{Le funzioni di interrogazione del DNS}
\label{sec:sock_resolver_functions}

Prima di trattare le funzioni usate normalmente nella risoluzione dei nomi a
dominio conviene trattare in maniera più dettagliata il servizio DNS. Come
accennato questo, benché esso in teoria sia solo uno dei possibili supporti su
cui mantenere le informazioni, in pratica costituisce il meccanismo principale
con cui vengono risolti i nomi a dominio. Inolte esso può fornire anche
ulteriori informazioni oltre relative alla risoluzione dei nomi a dominio.
Per questo motivo esistono una serie di funzioni di libreria che servono
specificamente ad eseguire delle interrogazioni verso un server DNS, funzioni
che poi vengono utilizzate anche per realizzare le funzioni generiche di
libreria usate dal sistema del \textit{resolver}.

Il sistema del DNS è in sostanza di un database distribuito organizzato in
maniera gerarchica, i dati vengono mantenuti in tanti server distinti ciascuno
dei quali si occupa della risoluzione del proprio \textsl{dominio}; i nomi a
dominio sono organizzati in una struttura ad albero analoga a quella
dell'albero dei file, con domini di primo livello (come i \texttt{.org}),
secondo livello (come \texttt{.truelite.it}), ecc.  In questo caso le
separazioni sono fra i vari livelli sono definite dal carattere ``\texttt{.}''
ed i nomi devono essere risolti da destra verso sinistra.\footnote{per chi si
  stia chiedendo quale sia la radice di questo albero, cioè l'equivalente di
  ``\texttt{/}'', la risposta è il dominio speciale ``\texttt{.}'', che in
  genere non viene mai scritto esplicitamente, ma che, come chiunque abbia
  configurato un server DNS sa bene, esiste ed è gestito dai cosiddetti
  \textit{root DNS} che risolvono i domini di primo livello.} Il meccanismo
funziona con il criterio della \textsl{delegazione}, un server responsabile
per un dominio di primo livello può delegare la risoluzione degli indirizzi
per un suo dominio di secondo livello ad un altro server, il quale a sua volta
potrà delegare la risoluzione di un eventuale sotto-dominio di terzo livello ad
un altro server ancora.

Come accennato un server DNS è in grado di fare molto altro rispetto alla
risoluzione di un nome a dominio in un indirizzo IP: ciascuna voce nel
database viene chiamata \textit{resource record}, e può contenere diverse
informazioni. In genere i \textit{resource record} vengono classificati per la
\textsl{classe di indirizzi} cui i dati contenuti fanno riferimento, e per il
\textsl{tipo} di questi ultimi (ritroveremo classi di indirizzi e tipi di
record più avanti in tab.~\ref{tab:DNS_address_class} e
tab.~\ref{tab:DNS_record_type}).  Oggigiorno i dati mantenuti nei server DNS
sono quasi esclusivamente relativi ad indirizzi internet, per cui in pratica
viene utilizzata soltanto una classe di indirizzi; invece le corrispondenze
fra un nome a dominio ed un indirizzo IP sono solo uno fra i vari tipi di
informazione che un server DNS fornisce normalmente.

L'esistenza di vari tipi di informazioni è un altro dei motivi per cui il
\textit{resolver} prevede, oltre a quelle relative alla semplice risoluzione
dei nomi, un insieme di funzioni specifiche dedicate all'interrogazione di un
server DNS, tutte nella forma \texttt{res\_}\textsl{\texttt{nome}}. La prima
di queste funzioni è \funcd{res\_init}, il cui prototipo è:

\begin{funcproto}{
\fhead{netinet/in.h} 
\fhead{arpa/nameser.h} 
\fhead{resolv.h}
\fdecl{int res\_init(void)}
\fdesc{Inizializza il sistema del \textit{resolver}.} 
}
{La funzione ritorna $0$ in caso di successo e $-1$ per un errore.
}
\end{funcproto}

La funzione legge il contenuto dei file di configurazione per impostare il
dominio di default, gli indirizzi dei server DNS da contattare e l'ordine
delle ricerche; se non sono specificati server verrà utilizzato l'indirizzo
locale, e se non è definito un dominio di default sarà usato quello associato
con l'indirizzo locale (ma questo può essere sovrascritto con l'uso della
variabile di ambiente \envvar{LOCALDOMAIN}). In genere non è necessario
eseguire questa funzione esplicitamente, in quanto viene automaticamente
chiamata la prima volta che si esegue una qualunque delle altre.

Le impostazioni e lo stato del \textit{resolver} inizializzati da
\func{res\_init} vengono mantenuti in una serie di variabili raggruppate nei
campi di una apposita struttura. Questa struttura viene definita in
\headfiled{resolv.h} e mantenuta nella variabile globale \var{\_res}, che
viene utilizzata internamente da tutte le funzioni dell'interfaccia. Questo
consente anche di accedere direttamente al contenuto della variabile
all'interno di un qualunque programma, una volta che la sia opportunamente
dichiarata con:
\includecodesnip{listati/resolv_option.c}

Dato che l'uso di una variabile globale rende tutte le funzioni
dell'interfaccia classica non rientranti, queste sono state deprecate in
favore di una nuova interfaccia in cui esse sono state sostituite da
altrettante nuove funzioni, il cui nome è ottenuto apponendo una
``\texttt{n}'' al nome di quella tradizionale (cioè nella forma
\texttt{res\_n\textsl{nome}}). Tutte le nuove funzioni sono identiche alle
precedenti, ma hanno un primo argomento aggiuntivo, \param{statep}, puntatore
ad una struttura dello stesso tipo di \var{\_res}. Questo consente di usare
una variabile locale per mantenere lo stato del \textit{resolver}, rendendo le
nuove funzioni rientranti.  In questo caso per poter utilizzare il nuovo
argomento occorrerà una opportuna dichiarazione del relativo tipo di dato con:
\includecodesnip{listati/resolv_newoption.c}

Così la nuova funzione utilizzata per inizializzare il \textit{resolver} (che
come la precedente viene chiamata automaticamente da tutte altre funzioni) è
\funcd{res\_ninit}, ed il suo prototipo è:

\begin{funcproto}{
\fhead{netinet/in.h} 
\fhead{arpa/nameser.h} 
\fhead{resolv.h}
\fdecl{int res\_ninit(res\_state statep)}
\fdesc{Inizializza il sistema del \textit{resolver}.} 
}
{La funzione ritorna $0$ in caso di successo e $-1$ per un errore.
}
\end{funcproto}

Indipendentemente da quale versione delle funzioni si usino, tutti i campi
della struttura (\var{\_res} o la variabile puntata da \param{statep}) sono ad
uso interno, e vengono usualmente inizializzate da \func{res\_init} o
\func{res\_ninit} in base al contenuto dei file di configurazione e ad una
serie di valori di default. L'unico campo che può essere utile modificare è
\var{\_res.options} (o l'equivalente della variabile puntata
da\param{statep}), una maschera binaria che contiene una serie di bit che
esprimono le opzioni che permettono di controllare il comportamento del
\textit{resolver}.

\begin{table}[htb]
  \centering
  \footnotesize
  \begin{tabular}[c]{|l|p{8cm}|}
    \hline
    \textbf{Costante} & \textbf{Significato} \\
    \hline
    \hline
    \constd{RES\_INIT}       & Viene attivato se è stata chiamata
                               \func{res\_init}. \\
    \constd{RES\_DEBUG}      & Stampa dei messaggi di debug.\\
    \constd{RES\_AAONLY}     & Accetta solo risposte autoritative.\\
    \constd{RES\_USEVC}      & Usa connessioni TCP per contattare i server 
                               invece che l'usuale UDP.\\
    \constd{RES\_PRIMARY}    & Interroga soltanto server DNS primari.\\
    \constd{RES\_IGNTC}      & Ignora gli errori di troncamento, non ritenta la
                               richiesta con una connessione TCP.\\
    \constd{RES\_RECURSE}    & Imposta il bit che indica che si desidera
                               eseguire una interrogazione ricorsiva.\\
    \constd{RES\_DEFNAMES}   & Se attivo \func{res\_search} aggiunge il nome
                               del dominio di default ai nomi singoli (che non
                               contengono cioè un ``\texttt{.}'').\\
    \constd{RES\_STAYOPEN}   & Usato con \const{RES\_USEVC} per mantenere
                               aperte le connessioni TCP fra interrogazioni
                               diverse.\\
    \constd{RES\_DNSRCH}     & Se attivo \func{res\_search} esegue le ricerche
                               di nomi di macchine nel dominio corrente o nei
                               domini ad esso sovrastanti.\\
    \constd{RES\_INSECURE1}  & Blocca i controlli di sicurezza di tipo 1.\\
    \constd{RES\_INSECURE2}  & Blocca i controlli di sicurezza di tipo 2.\\
    \constd{RES\_NOALIASES}  & Blocca l'uso della variabile di ambiente
                               \envvar{HOSTALIASES}.\\ 
    \constd{RES\_USE\_INET6} & Restituisce indirizzi IPv6 con
                               \func{gethostbyname}. \\
    \constd{RES\_ROTATE}     & Ruota la lista dei server DNS dopo ogni
                               interrogazione.\\
    \constd{RES\_NOCHECKNAME}& Non controlla i nomi per verificarne la
                               correttezza sintattica. \\
    \constd{RES\_KEEPTSIG}   & Non elimina i record di tipo \texttt{TSIG}.\\
    \constd{RES\_BLAST}      & Effettua un ``\textit{blast}'' inviando
                               simultaneamente le richieste a tutti i server;
                               non ancora implementata. \\
    \constd{RES\_DEFAULT}    & Combinazione di \const{RES\_RECURSE},
                               \const{RES\_DEFNAMES} e \const{RES\_DNSRCH}.\\
    \hline
  \end{tabular}
  \caption{Costanti utilizzabili come valori per \var{\_res.options}.}
  \label{tab:resolver_option}
\end{table}

Per utilizzare questa funzionalità per modificare le impostazioni direttamente
da programma occorrerà impostare un opportuno valore per questo campo ed
invocare esplicitamente \func{res\_init} o \func{res\_ninit}, dopo di che le
altre funzioni prenderanno le nuove impostazioni. Le costanti che definiscono
i vari bit di questo campo, ed il relativo significato sono illustrate in
tab.~\ref{tab:resolver_option}; trattandosi di una maschera binaria un valore
deve essere espresso con un opportuno OR aritmetico di dette costanti; ad
esempio il valore di default delle opzioni, espresso dalla costante
\const{RES\_DEFAULT}, è definito come:
\includecodesnip{listati/resolv_option_def.c}

Non tratteremo il significato degli altri campi non essendovi necessità di
modificarli direttamente; gran parte di essi sono infatti impostati dal
contenuto dei file di configurazione, mentre le funzionalità controllate da
alcuni di esse possono essere modificate con l'uso delle opportune variabili
di ambiente come abbiamo visto per \envvar{LOCALDOMAIN}. In particolare con
\envvar{RES\_RETRY} si soprassiede il valore del campo \var{retry} che
controlla quante volte viene ripetuto il tentativo di connettersi ad un server
DNS prima di dichiarare fallimento; il valore di default è 4, un valore nullo
significa bloccare l'uso del DNS. Infine con \envvar{RES\_TIMEOUT} si
soprassiede il valore del campo \var{retrans} (preimpostato al valore della
omonima costante \const{RES\_TIMEOUT} di \headfile{resolv.h}) che è il valore
preso come base (in numero di secondi) per definire la scadenza di una
richiesta, ciascun tentativo di richiesta fallito viene ripetuto raddoppiando
il tempo di scadenza per il numero massimo di volte stabilito da
\texttt{RES\_RETRY}.

La funzione di interrogazione principale è \funcd{res\_query}
(\funcd{res\_nquery} per la nuova interfaccia), che serve ad eseguire una
richiesta ad un server DNS per un nome a dominio \textsl{completamente
  specificato} (quello che si chiama
\itindex{Fully~Qualified~Domain~Name~(FQDN)} FQDN, \textit{Fully Qualified 
  Domain Name}); il loro prototipo è:

\begin{funcproto}{
\fhead{netinet/in.h} 
\fhead{arpa/nameser.h} 
\fhead{resolv.h}
\fdecl{int res\_query(const char *dname, int class, int type,
              unsigned char *answer, int anslen)}
\fdecl{int res\_nquery(res\_state statep, const char *dname, int class, int
  type, \\
  \phantom{int res\_nquery(}unsigned char *answer, int anslen)}
\fdesc{Esegue una interrogazione al DNS.} 
}
{Le funzioni ritornano un valore positivo pari alla lunghezza dei dati scritti
  nel buffer \param{answer} in caso di successo e $-1$ per un errore.
}
\end{funcproto}

Le funzioni eseguono una interrogazione ad un server DNS relativa al nome da
risolvere passato nella stringa indirizzata da \param{dname}, inoltre deve
essere specificata la classe di indirizzi in cui eseguire la ricerca con
\param{class}, ed il tipo di \textit{resource record} che si vuole ottenere
con \param{type}. Il risultato della ricerca verrà scritto nel buffer di
lunghezza \param{anslen} puntato da \param{answer} che si sarà opportunamente
allocato in precedenza.

Una seconda funzione di ricerca analoga a \func{res\_query}, che prende gli
stessi argomenti ma che esegue l'interrogazione con le funzionalità
addizionali previste dalle due opzioni \const{RES\_DEFNAMES} e
\const{RES\_DNSRCH}, è \funcd{res\_search} (\funcd{res\_nsearch} per la nuova
interfaccia), il cui prototipo è:

\begin{funcproto}{
\fhead{netinet/in.h} 
\fhead{arpa/nameser.h} 
\fhead{resolv.h}
\fdecl{int res\_search(const char *dname, int class, int type,
  unsigned char *answer, \\
  \phantom{int res\_search}int anslen)}
\fdecl{int res\_nsearch(res\_state statep, const char *dname, int class, 
  int type, \\
  \phantom{int res\_nsearch(}unsigned char *answer, int anslen)}
\fdesc{Esegue una interrogazione al DNS.} 
}
{Le funzioni ritornano un valore positivo pari alla lunghezza dei dati scritti
  nel buffer \param{answer} in caso di successo e $-1$ per un errore.
}
\end{funcproto}

In sostanza la funzione ripete una serie di chiamate a \func{res\_query}
(\func{res\_nquery}) aggiungendo al nome contenuto nella stringa \param{dname}
il dominio di default da cercare, fermandosi non appena trova un risultato.
Il risultato di entrambe le funzioni viene scritto nel formato opportuno (che
sarà diverso a seconda del tipo di record richiesto) nel buffer di ritorno;
sarà compito del programma (o di altre funzioni) estrarre i relativi dati,
esistono una serie di funzioni interne usate per la scansione di questi dati,
per chi fosse interessato una trattazione dettagliata è riportata nel
quattordicesimo capitolo di \cite{DNSbind}.

Le classi di indirizzi supportate da un server DNS sono tre, ma di queste in
pratica oggi viene utilizzata soltanto quella degli indirizzi internet; le
costanti che identificano dette classi, da usare come valore per l'argomento
\param{class} delle precedenti funzioni, sono riportate in
tab.~\ref{tab:DNS_address_class} (esisteva in realtà anche una classe
\constd{C\_CSNET} per la omonima rete, ma è stata dichiarata obsoleta).

\begin{table}[htb]
  \centering
  \footnotesize
  \begin{tabular}[c]{|l|p{8cm}|}
    \hline
    \textbf{Costante} & \textbf{Significato} \\
    \hline
    \hline
    \constd{C\_IN}   & Indirizzi internet, in pratica i soli utilizzati oggi.\\
    \constd{C\_HS}   & Indirizzi \textit{Hesiod}, utilizzati solo al MIT, oggi
                       completamente estinti. \\
    \constd{C\_CHAOS}& Indirizzi per la rete \textit{Chaosnet}, un'altra rete
                       sperimentale nata al MIT. \\
    \constd{C\_ANY}  & Indica un indirizzo di classe qualunque.\\
    \hline
  \end{tabular}
  \caption{Costanti identificative delle classi di indirizzi per l'argomento
    \param{class} di \func{res\_query}.}
  \label{tab:DNS_address_class}
\end{table}

Come accennato le tipologie di dati che sono mantenibili su un server DNS sono
diverse, ed a ciascuna di essa corrisponde un diverso tipo di \textit{resource
  record}. L'elenco delle costanti, ripreso dai file di dichiarazione
\headfiled{arpa/nameser.h} e \headfiled{arpa/nameser\_compat.h}, che
definiscono i valori che si possono usare per l'argomento \param{type} per
specificare il tipo di \textit{resource record} da richiedere è riportato in
tab.~\ref{tab:DNS_record_type}; le costanti (tolto il \texttt{T\_} iniziale)
hanno gli stessi nomi usati per identificare i record nei file di zona di
BIND,\footnote{BIND, acronimo di \textit{Berkley Internet Name Domain}, è una
  implementazione di un server DNS, ed, essendo utilizzata nella stragrande
  maggioranza dei casi, fa da riferimento; i dati relativi ad un certo dominio
  (cioè i suoi \textit{resource record} vengono mantenuti in quelli che sono
  usualmente chiamati \textsl{file di zona}, e in essi ciascun tipo di dominio
  è identificato da un nome che è appunto identico a quello delle costanti di
  tab.~\ref{tab:DNS_record_type} senza il \texttt{T\_} iniziale.} e che
normalmente sono anche usati come nomi per indicare i record.

\begin{table}[!htb]
  \centering
  \footnotesize
  \begin{tabular}[c]{|l|l|}
    \hline
    \textbf{Costante} & \textbf{Significato} \\
    \hline
    \hline
    \constd{T\_A}     & Indirizzo di una stazione.\\
    \constd{T\_NS}    & Server DNS autoritativo per il dominio richiesto.\\
    \constd{T\_MD}    & Destinazione per la posta elettronica.\\
    \constd{T\_MF}    & Redistributore per la posta elettronica.\\
    \constd{T\_CNAME} & Nome canonico.\\
    \constd{T\_SOA}   & Inizio di una zona di autorità.\\
    \constd{T\_MB}    & Nome a dominio di una casella di posta.\\
    \constd{T\_MG}    & Nome di un membro di un gruppo di posta.\\
    \constd{T\_MR}    & Nome di un cambiamento di nome per la posta.\\
    \constd{T\_NULL}  & Record nullo.\\
    \constd{T\_WKS}   & Servizio noto.\\
    \constd{T\_PTR}   & Risoluzione inversa di un indirizzo numerico.\\
    \constd{T\_HINFO} & Informazione sulla stazione.\\
    \constd{T\_MINFO} & Informazione sulla casella di posta.\\
    \constd{T\_MX}    & Server cui instradare la posta per il dominio.\\
    \constd{T\_TXT}   & Stringhe di testo (libere).\\
    \constd{T\_RP}    & Nome di un responsabile (\textit{responsible person}).\\
    \constd{T\_AFSDB} & Database per una cella AFS.\\
    \constd{T\_X25}   & Indirizzo di chiamata per X.25.\\
    \constd{T\_ISDN}  & Indirizzo di chiamata per ISDN.\\
    \constd{T\_RT}    & Router.\\
    \constd{T\_NSAP}  & Indirizzo NSAP.\\
    \constd{T\_NSAP\_PTR}& Risoluzione inversa per NSAP (deprecato).\\
    \constd{T\_SIG}   & Firma digitale di sicurezza.\\
    \constd{T\_KEY}   & Chiave per firma.\\
    \constd{T\_PX}    & Corrispondenza per la posta X.400.\\
    \constd{T\_GPOS}  & Posizione geografica.\\
    \constd{T\_AAAA}  & Indirizzo IPv6.\\
    \constd{T\_LOC}   & Informazione di collocazione.\\
    \constd{T\_NXT}   & Dominio successivo.\\
    \constd{T\_EID}   & Identificatore di punto conclusivo.\\
    \constd{T\_NIMLOC}& Posizionatore \textit{nimrod}.\\
    \constd{T\_SRV}   & Servizio.\\
    \constd{T\_ATMA}  & Indirizzo ATM.\\
    \constd{T\_NAPTR} & Puntatore ad una \textit{naming authority}.\\
    \constd{T\_TSIG}  & Firma di transazione.\\
    \constd{T\_IXFR}  & Trasferimento di zona incrementale.\\
    \constd{T\_AXFR}  & Trasferimento di zona di autorità.\\
    \constd{T\_MAILB} & Trasferimento di record di caselle di posta.\\
    \constd{T\_MAILA} & Trasferimento di record di server di posta.\\
    \constd{T\_ANY}   & Valore generico.\\
    \hline
  \end{tabular}
  \caption{Costanti identificative del tipo di record per l'argomento
    \param{type} di \func{res\_query}.}
  \label{tab:DNS_record_type}
\end{table}


L'elenco di tab.~\ref{tab:DNS_record_type} è quello di \textsl{tutti} i
\textit{resource record} definiti, con una breve descrizione del relativo
significato.  Di tutti questi però viene impiegato correntemente solo un
piccolo sottoinsieme, alcuni sono obsoleti ed altri fanno riferimento a dati
applicativi che non ci interessano non avendo nulla a che fare con la
risoluzione degli indirizzi IP, pertanto non entreremo nei dettagli del
significato di tutti i \textit{resource record}, ma solo di quelli usati dalle
funzioni del \textit{resolver}. Questi sono sostanzialmente i seguenti (per
indicarli si è usata la notazione dei file di zona di BIND):
\begin{basedescript}{\desclabelwidth{1.2cm}\desclabelstyle{\nextlinelabel}}
\item[\texttt{A}] viene usato per indicare la corrispondenza fra un nome a
  dominio ed un indirizzo IPv4; ad esempio la corrispondenza fra
  \texttt{jojo.truelite.it} e l'indirizzo IP \texttt{62.48.34.25}.
\item[\texttt{AAAA}] viene usato per indicare la corrispondenza fra un nome a
  dominio ed un indirizzo IPv6; è chiamato in questo modo dato che la
  dimensione di un indirizzo IPv6 è quattro volte quella di un indirizzo IPv4.
\item[\texttt{PTR}] per fornire la corrispondenza inversa fra un indirizzo IP
  ed un nome a dominio ad esso associato si utilizza questo tipo di record (il
  cui nome sta per \textit{pointer}).
\item[\texttt{CNAME}] qualora si abbiamo più nomi che corrispondono allo
  stesso indirizzo (come ad esempio \texttt{www.truelite.it} e
  \texttt{sources.truelite.it}, che fanno entrambi riferimento alla stessa
  macchina (nel caso \texttt{dodds.truelite.it}) si può usare questo tipo di
  record per creare degli \textit{alias} in modo da associare un qualunque
  altro nome al \textsl{nome canonico} della macchina (si chiama così quello
  associato al record \texttt{A}).
\end{basedescript}

Come accennato in caso di successo le due funzioni di richiesta restituiscono
il risultato della interrogazione al server, in caso di insuccesso l'errore
invece viene segnalato da un valore di ritorno pari a $-1$, ma in questo caso,
non può essere utilizzata la variabile \var{errno} per riportare un codice di
errore, in quanto questo viene impostato per ciascuna delle chiamate al
sistema utilizzate dalle funzioni del \textit{resolver}, non avrà alcun
significato nell'indicare quale parte del procedimento di risoluzione è
fallita.

Per questo motivo è stata definita una variabile di errore separata,
\var{h\_errno}, che viene utilizzata dalle funzioni del \textit{resolver} per
indicare quale problema ha causato il fallimento della risoluzione del nome.
Ad essa si può accedere una volta che la si dichiara con:
\includecodesnip{listati/herrno.c} 
ed i valori che può assumere, con il relativo significato, sono riportati in
tab.~\ref{tab:h_errno_values}.

\begin{table}[!htb]
  \centering
  \footnotesize
  \begin{tabular}[c]{|l|p{9cm}|}
    \hline
    \textbf{Costante} & \textbf{Significato} \\
    \hline
    \hline
    \constd{HOST\_NOT\_FOUND}& L'indirizzo richiesto non è valido e la
                               macchina indicata è sconosciuta.\\
    \constd{NO\_ADDRESS}     & Il nome a dominio richiesto è valido, ma non ha
                               un indirizzo associato ad esso
                               (alternativamente può essere indicato come 
                               \constd{NO\_DATA}).\\
    \constd{NO\_RECOVERY}    & Si è avuto un errore non recuperabile
                               nell'interrogazione di un server DNS.\\
    \constd{TRY\_AGAIN}      & Si è avuto un errore temporaneo
                               nell'interrogazione di un server DNS, si può
                               ritentare l'interrogazione in un secondo
                               tempo.\\
    \hline
  \end{tabular}
  \caption{Valori possibili della variabile \var{h\_errno}.}
  \label{tab:h_errno_values}
\end{table}

Insieme alla nuova variabile vengono definite anche delle nuove funzioni per
stampare l'errore a video, analoghe a quelle di sez.~\ref{sec:sys_strerror}
per \var{errno}, ma che usano il valore di \var{h\_errno}; la prima è
\funcd{herror} ed il suo prototipo è:

{\centering
\vspace{3pt}
\begin{funcbox}{
\fhead{netdb.h}
\fdecl{void herror(const char *string)}
\fdesc{Stampa un errore di risoluzione.} 
}
\end{funcbox}}

La funzione è l'analoga di \func{perror} e stampa sullo \textit{standard error} un
messaggio di errore corrispondente al valore corrente di \var{h\_errno}, a cui
viene anteposta la stringa \param{string} passata come argomento.  La seconda
funzione è \funcd{hstrerror} ed il suo prototipo è:

{\centering
\vspace{3pt}
\begin{funcbox}{
\fhead{netdb.h}
\fdecl{const char *hstrerror(int err)}
\fdesc{Restituisce una stringa corrispondente ad un errore di risoluzione.}
}
\end{funcbox}}

\noindent che, come  l'analoga \func{strerror}, restituisce una stringa con un
messaggio di errore già formattato, corrispondente al codice passato come
argomento (che si presume sia dato da \var{h\_errno}).

\itindend{resolver}


\subsection{La vecchia interfaccia per la risoluzione dei nomi a dominio}
\label{sec:sock_name_services}

La principale funzionalità del \textit{resolver} resta quella di risolvere i
nomi a dominio in indirizzi IP, per cui non ci dedicheremo oltre alle funzioni
per effetture delle richieste generiche al DNS ed esamineremo invece le
funzioni del \textit{resolver} dedicate specificamente a questo. Tratteremo in
questa sezione l'interfaccia tradizionale, che ormai è deprecata, mentre
vedremo nella sezione seguente la nuova interfaccia.

La prima funzione dell'interfaccia tradizionale è \funcd{gethostbyname} il cui
scopo è ottenere l'indirizzo di una stazione noto il suo nome a dominio, il
suo prototipo è:

\begin{funcproto}{
\fhead{netdb.h}
\fdecl{struct hostent *gethostbyname(const char *name)}
\fdesc{Determina l'indirizzo associato ad un nome a dominio.} 
}

{La funzione ritorna il puntatore ad una
  struttura di tipo \struct{hostent} contenente i dati associati al nome a
  dominio in caso di successo o un puntatore nullo per un errore.}
\end{funcproto}

La funzione prende come argomento una stringa \param{name} contenente il nome
a dominio che si vuole risolvere, in caso di successo i dati ad esso relativi
vengono memorizzati in una opportuna struttura \struct{hostent} la cui
definizione è riportata in fig.~\ref{fig:sock_hostent_struct}. 

\begin{figure}[!htb]
  \footnotesize \centering
  \begin{minipage}[c]{0.80\textwidth}
    \includestruct{listati/hostent.h}
  \end{minipage}
  \caption{La struttura \structd{hostent} per la risoluzione dei nomi a
    dominio e degli indirizzi IP.}
  \label{fig:sock_hostent_struct}
\end{figure}

Quando un programma chiama \func{gethostbyname} e questa usa il DNS per
effettuare la risoluzione del nome, è con i valori contenuti nei relativi
record che vengono riempite le varie parti della struttura \struct{hostent}.
Il primo campo della struttura, \var{h\_name} contiene sempre il \textsl{nome
  canonico}, che nel caso del DNS è appunto il nome associato ad un record
\texttt{A}. Il secondo campo della struttura, \var{h\_aliases}, invece è un
puntatore ad vettore di puntatori, terminato da un puntatore nullo. Ciascun
puntatore del vettore punta ad una stringa contenente uno degli altri
possibili nomi associati allo stesso \textsl{nome canonico} (quelli che nel
DNS vengono inseriti come record di tipo \texttt{CNAME}).

Il terzo campo della struttura, \var{h\_addrtype}, indica il tipo di indirizzo
che è stato restituito, e può assumere soltanto i valori \const{AF\_INET} o
\const{AF\_INET6}, mentre il quarto campo, \var{h\_length}, indica la
lunghezza dell'indirizzo stesso in byte. 

Infine il campo \var{h\_addr\_list} è il puntatore ad un vettore di puntatori
ai singoli indirizzi; il vettore è terminato da un puntatore nullo.  Inoltre,
come illustrato in fig.~\ref{fig:sock_hostent_struct}, viene definito il campo
\var{h\_addr} come sinonimo di \code{h\_addr\_list[0]}, cioè un riferimento
diretto al primo indirizzo della lista.

Oltre ai normali nomi a dominio la funzione accetta come argomento
\param{name} anche indirizzi numerici, in formato dotted decimal per IPv4 o
con la notazione illustrata in sez.~\ref{sec:IP_ipv6_notation} per IPv6. In
tal caso \func{gethostbyname} non eseguirà nessuna interrogazione remota, ma
si limiterà a copiare la stringa nel campo \var{h\_name} ed a creare la
corrispondente struttura \var{in\_addr} da indirizzare con
\code{h\_addr\_list[0]}.

Con l'uso di \func{gethostbyname} si ottengono solo gli indirizzi IPv4, se si
vogliono ottenere degli indirizzi IPv6 occorrerà prima impostare l'opzione
\const{RES\_USE\_INET6} nel campo \texttt{\_res.options} e poi chiamare
\func{res\_init} (vedi sez.~\ref{sec:sock_resolver_functions}) per modificare
le opzioni del \textit{resolver}; dato che questo non è molto comodo è stata
definita (è una estensione fornita dalla \acr{glibc}, disponibile anche in
altri sistemi unix-like) un'altra funzione, \funcd{gethostbyname2}, il cui
prototipo è:

\begin{funcproto}{
\fhead{netdb.h}
\fhead{sys/socket.h}
\fdecl{struct hostent *gethostbyname2(const char *name, int af)}
\fdesc{Determina l'indirizzo del tipo scelto associato ad un nome a dominio.} 
}

{La funzione ritorna il puntatore ad una
  struttura di tipo \struct{hostent} contenente i dati associati al nome a
  dominio in caso di successo e  un puntatore nullo per un errore.}
\end{funcproto}

In questo caso la funzione prende un secondo argomento \param{af} che indica
quale famiglia di indirizzi è quella che dovrà essere utilizzata per
selezionare i risultati restituiti dalla funzione; i soli valori consentiti
sono \const{AF\_INET} o \const{AF\_INET6} per indicare rispettivamente IPv4 e
IPv6 (per questo è necessario l'uso di \headfile{sys/socket.h}). Per tutto il
resto la funzione è identica a \func{gethostbyname}, ed identici sono i suoi
risultati.

\begin{figure}[!htb]
  \footnotesize \centering
  \begin{minipage}[c]{\codesamplewidth}
    \includecodesample{listati/mygethost.c}
  \end{minipage}
  \normalsize
  \caption{Esempio di codice per la risoluzione di un indirizzo.}
  \label{fig:mygethost_example}
\end{figure}

Vediamo allora un primo esempio dell'uso delle funzioni di risoluzione, in
fig.~\ref{fig:mygethost_example} è riportato un estratto del codice di un
programma che esegue una semplice interrogazione al \textit{resolver} usando
\func{gethostbyname} e poi ne stampa a video i risultati. Al solito il
sorgente completo, che comprende il trattamento delle opzioni ed una funzione
per stampare un messaggio di aiuto, è nel file \texttt{mygethost.c} dei
sorgenti allegati alla guida.

Il programma richiede un solo argomento che specifichi il nome da cercare,
senza il quale (\texttt{\small 15--18}) esce con un errore. Dopo di che
(\texttt{\small 20}) si limita a chiamare \func{gethostbyname}, ricevendo il
risultato nel puntatore \var{data}. Questo (\texttt{\small 21--24}) viene
controllato per rilevare eventuali errori, nel qual caso il programma esce
dopo aver stampato un messaggio con \func{herror}. 

Se invece la risoluzione è andata a buon fine si inizia (\texttt{\small 25})
con lo stampare il nome canonico, dopo di che (\texttt{\small 26--30}) si
stampano eventuali altri nomi. Per questo prima (\texttt{\small 26}) si prende
il puntatore alla cima della lista che contiene i nomi e poi (\texttt{\small
  27--30}) si esegue un ciclo che sarà ripetuto fin tanto che nella lista si
troveranno dei puntatori validi per le stringhe dei nomi (si ricordi che la
lista viene terminata da un puntatore nullo); prima (\texttt{\small 28}) si
stamperà la stringa e poi (\texttt{\small 29}) si provvederà ad incrementare
il puntatore per passare al successivo elemento della lista.

Una volta stampati i nomi si passerà a stampare gli indirizzi, il primo passo
(\texttt{\small 31--38}) è allora quello di riconoscere il tipo di indirizzo
sulla base del valore del campo \var{h\_addrtype}, stampandolo a video. Si è
anche previsto di stampare un errore nel caso (che non dovrebbe mai accadere)
di un indirizzo non valido.

Infine (\texttt{\small 39--44}) si stamperanno i valori degli indirizzi, di
nuovo (\texttt{\small 39}) si inizializzerà un puntatore alla cima della lista
e si eseguirà un ciclo fintanto che questo punterà ad indirizzi validi in
maniera analoga a quanto fatto in precedenza per i nomi a dominio. Si noti
come, essendo il campo \var{h\_addr\_list} un puntatore ad strutture di
indirizzi generiche, questo sia ancora di tipo \texttt{char **} e si possa
riutilizzare lo stesso puntatore usato per i nomi.

Per ciascun indirizzo valido si provvederà (\texttt{\small 41}) a una
conversione con la funzione \func{inet\_ntop} (vedi
sez.~\ref{sec:sock_addr_func}) passandole gli opportuni argomenti, questa
restituirà la stringa da stampare (\texttt{\small 42}) con il valore
dell'indirizzo in \var{buffer}, che si è avuto la cura di dichiarare
inizialmente (\texttt{\small 10}) con dimensioni adeguate; dato che la
funzione è in grado di tenere conto automaticamente del tipo di indirizzo non
ci sono precauzioni particolari da prendere.\footnote{volendo essere pignoli
  si dovrebbe controllarne lo stato di uscita, lo si è tralasciato per non
  appesantire il codice, dato che in caso di indirizzi non validi si sarebbe
  avuto un errore con \func{gethostbyname}, ma si ricordi che la sicurezza non
  è mai troppa.}

Le funzioni illustrate finora hanno un difetto: utilizzano tutte una area di
memoria interna per allocare i contenuti della struttura \struct{hostent} e
per questo non possono essere rientranti. L'uso della memoria interna inoltre
comporta anche che in due successive chiamate i dati potranno essere
sovrascritti.

Si tenga presente poi che copiare il contenuto della sola struttura non è
sufficiente per salvare tutti i dati, in quanto questa contiene puntatori ad
altri dati, che pure possono essere sovrascritti; per questo motivo, se si
vuole salvare il risultato di una chiamata, occorrerà eseguire quella che si
chiama una \itindex{deep~copy} \textit{deep copy}.\footnote{si chiama così
  quella tecnica per cui, quando si deve copiare il contenuto di una struttura
  complessa (con puntatori che puntano ad altri dati, che a loro volta possono
  essere puntatori ad altri dati) si deve copiare non solo il contenuto della
  struttura, ma eseguire una scansione per risolvere anche tutti i puntatori
  contenuti in essa (e così via se vi sono altre sotto-strutture con altri
  puntatori) e copiare anche i dati da questi referenziati.}

Per ovviare a questi problemi nella \acr{glibc} sono definite anche delle
versioni rientranti delle precedenti funzioni, al solito queste sono
caratterizzate dall'avere un suffisso \texttt{\_r}, pertanto avremo le due
funzioni \funcd{gethostbyname\_r} e \funcd{gethostbyname2\_r} i cui prototipi
sono:

\begin{funcproto}{
\fhead{netdb.h}
\fhead{sys/socket.h}
\fdecl{int gethostbyname\_r(const char *name, struct hostent *ret, 
    char *buf, size\_t buflen,\\
\phantom{int gethostbyname\_r(}struct hostent **result, int *h\_errnop)}
\fdecl{int gethostbyname2\_r(const char *name, int af,
         struct hostent *ret, char *buf,\\
\phantom{int gethostbyname2\_r(}size\_t buflen, struct hostent **result, int *h\_errnop)}

\fdesc{Versioni rientranti delle funzioni \func{gethostbyname} e
  \func{gethostbyname2}.} 
}

{Le funzioni ritornano $0$ in caso di successo ed un valore diverso da zero per
  un errore.}
\end{funcproto}

Gli argomenti \param{name} (e \param{af} per \func{gethostbyname2\_r}) hanno
lo stesso significato visto in precedenza. Tutti gli altri argomenti hanno lo
stesso significato per entrambe le funzioni. Per evitare l'uso di variabili
globali si dovrà allocare preventivamente una struttura \struct{hostent} in
cui ricevere il risultato, passandone l'indirizzo alla funzione nell'argomento
\param{ret}.  Inoltre, dato che \struct{hostent} contiene dei puntatori, dovrà
essere allocato anche un buffer in cui le funzioni possano scrivere tutti i
dati del risultato dell'interrogazione da questi puntati; l'indirizzo e la
lunghezza di questo buffer devono essere indicati con gli argomenti
\param{buf} e \param{buflen}.

Gli ultimi due argomenti vengono utilizzati per avere indietro i risultati
come \textit{value result argument}, si deve specificare l'indirizzo della
variabile su cui la funzione dovrà salvare il codice di errore
con \param{h\_errnop} e quello su cui dovrà salvare il puntatore che si userà
per accedere i dati con \param{result}.

In caso di successo entrambe le funzioni restituiscono un valore nullo,
altrimenti restituiscono un valore non nulla e all'indirizzo puntato
da \param{result} sarà salvato un puntatore nullo, mentre a quello puntato da
\param{h\_errnop} sarà salvato il valore del codice di errore, dato che per
essere rientrante la funzione non può la variabile globale \var{h\_errno}. In
questo caso il codice di errore, oltre ai valori di
tab.~\ref{tab:h_errno_values}, può avere anche quello di \errcode{ERANGE}
qualora il buffer allocato su \param{buf} non sia sufficiente a contenere i
dati, in tal caso si dovrà semplicemente ripetere l'esecuzione della funzione
con un buffer di dimensione maggiore.

Una delle caratteristiche delle interrogazioni al servizio DNS è che queste
sono normalmente eseguite con il protocollo UDP, ci sono casi in cui si
preferisce che vengano usate connessioni permanenti con il protocollo TCP. Per
ottenere questo sono previste delle funzioni apposite (si potrebbero impostare
direttamente le opzioni di \var{\_\_res.options}, ma queste funzioni
permettono di semplificare la procedura); la prime sono \funcd{sethostent} e
\func{endhostent}, il cui prototipo è:

\begin{funcproto}{
\fhead{netdb.h}
\fdecl{void sethostent(int stayopen)}
\fdesc{Richiede l'uso di connessioni TCP per le interrogazioni ad un server DNS.} 
\fdecl{void endhostent(void)}
\fdesc{Disattiva l'uso di connessioni TCP per le interrogazioni ad un server DNS.} 
}

{Le funzioni non restituiscono nulla, e non danno errori.}
\end{funcproto}

La funzione \func{sethostent} permette di richiedere l'uso di connessioni TCP
per la richiesta dei dati, e che queste restino aperte per successive
richieste; il valore dell'argomento \param{stayopen} indica se attivare questa
funzionalità, un valore diverso da zero, che indica una condizione vera in C,
attiva la funzionalità.  Per disattivare l'uso delle connessioni TCP si può
invece usare \func{endhostent}, e come si vede la funzione è estremamente
semplice, non richiedendo nessun argomento.

Infine si può richiedere la risoluzione inversa di un indirizzo IP od IPv6,
per ottenerne il nome a dominio ad esso associato, per fare questo si può
usare la funzione \funcd{gethostbyaddr}, il cui prototipo è:

\begin{funcproto}{
\fhead{netdb.h}
\fhead{sys/socket.h} 
\fdecl{struct hostent *gethostbyaddr(const char *addr, int len, int type)}
\fdesc{Richiede la risoluzione inversa di un indirizzo IP.} 
}

{La funzione ritorna l'indirizzo ad una struttura \struct{hostent} in caso di
  successo e \val{NULL} per un errore.}
\end{funcproto}

In questo caso l'argomento \param{addr} dovrà essere il puntatore ad una
appropriata struttura contenente il valore dell'indirizzo IP (o IPv6) che si
vuole risolvere. L'uso del tipo \texttt{char *} per questo argomento è
storico, il dato dovrà essere fornito in una struttura \struct{in\_addr} per
un indirizzo IPv4 ed una struttura \struct{in6\_addr} per un indirizzo IPv6.

Si ricordi inoltre, come illustrato in fig.~\ref{fig:sock_sa_ipv4_struct}, che
mentre \struct{in\_addr} corrisponde in realtà ad un oridinario numero intero
a 32 bit  (da esprimere comunque in \textit{network order}) non altrettanto
avviene per \struct{in6\_addr}, pertanto è sempre opportuno inizializzare
questi indirizzi con \func{inet\_pton} (vedi
sez.~\ref{sec:sock_conv_func_gen}).

Nell'argomento \param{len} se ne dovrà poi specificare la dimensione
(rispettivamente 4 o 16), infine l'argomento
\param{type} deve indicare il tipo di indirizzo, e dovrà essere o
\const{AF\_INET} o \const{AF\_INET6}.

La funzione restituisce, in caso di successo, un puntatore ad una struttura
\struct{hostent}, solo che in questo caso la ricerca viene eseguita
richiedendo al DNS un record di tipo \texttt{PTR} corrispondente all'indirizzo
specificato. In caso di errore al solito viene usata la variabile
\var{h\_errno} per restituire un opportuno codice. In questo caso l'unico
campo del risultato che interessa è \var{h\_name} che conterrà il nome a
dominio, la funzione comunque inizializza anche il primo campo della lista
\var{h\_addr\_list} col valore dell'indirizzo passato come argomento.

Dato che \func{gethostbyaddr} usa un buffer statico, anche di questa funzione
esiste una versione rientrante \funcd{gethostbyaddr\_r} fornita come
estensione dalla \acr{glibc}, il cui prototipo è:

\begin{funcproto}{
\fhead{netdb.h}
\fhead{sys/socket.h} 
\fdecl{struct hostent *gethostbyaddr\_r(const void *addr, socklen\_t len, int type,\\
\phantom{struct hostent *gethostbyaddr\_r(}struct hostent *ret, char *buf, size\_t buflen,\\
\phantom{struct hostent *gethostbyaddr\_r(}struct hostent **result, int *h\_errnop);}
\fdesc{Richiede la risoluzione inversa di un indirizzo IP.} 
}

{La funzione ritorna $0$ in caso di successo e un valore non nullo per un errore.}
\end{funcproto}

La funzione prende per gli argomenti \param{addr}, \param{len} e \param{type}
gli stessi valori di \func{gethostbyaddr} con lo stesso significato, gli
argomenti successivi vengono utilizzati per restituire i dati, sono identici a
quelli già illustrati in per \func{gethostbyname\_r} e
\func{gethostbyname2\_r} e devono essere usati allo stesso modo.

Infine lo standard POSIX prevede la presenza della funzione
\funcd{gethostent}, il cui prototipo è:

\begin{funcproto}{
\fhead{netdb.h}
\fdecl{struct hostent *gethostent(void)}
\fdesc{Ottiene la voce successiva nel database dei nomi a dominio.} 
}

{La funzione ritorna l'indirizzo ad una struttura \struct{hostent} in caso di
  successo e \val{NULL} per un errore.}
\end{funcproto}

La funzione dovrebbe ritornare (come puntatore alla solita struttura
\struct{hostent} allocata internamente) la voce successiva nel database dei
nomi a dominio, ma questo ha un significato soltato quando è relativo alla
lettura dei dati da un file come \conffile{/etc/hosts} e non per i risultati
del DNS. Nel caso della \acr{glibc} questa viene usata allora solo per la
lettura delle voci presenti in quest'ultimo, come avviene anche in altri
sistemi unix-like, ed inoltre ignora le voci relative ad indirizzi IPv6.

Della stessa funzione la \acr{glibc} fornisce anche una versione rientrante
\funcd{gethostent\_r}, il cui prototipo è:

\begin{funcproto}{
\fhead{netdb.h}
\fdecl{struct hostent *gethostent\_r(struct hostent *ret, char *buf, size\_t buflen,\\
\phantom{struct hostent *gethostent\_r(}struct hostent **result, int *h\_errnop);}
\fdesc{Ottiene la voce successiva nel database dei nomi a dominio.} 
}

{La funzione ritorna $0$ in caso di successo e un valore non nullo per un errore.}
\end{funcproto}

La funzione ha lo stesso effetto di \func{gethostent}; gli argomenti servono a
restituire i risultati in maniera rientrante e vanno usati secondo le modalità
già illustrate per \func{gethostbyname\_r} e \func{gethostbyname2\_r}.

Dati i limiti delle funzioni \func{gethostbyname} e \func{gethostbyaddr} con
l'uso di memoria statica che può essere sovrascritta fra due chiamate
successive, e per avere sempre la possibilità di indicare esplicitamente il
tipo di indirizzi voluto (cosa che non è possibile con \func{gethostbyname}),
è stata successivamente proposta,
nell'\href{http://www.ietf.org/rfc/rfc2553.txt}{RFC~2553} un diversa
interfaccia con l'introduzione due nuove funzioni di
risoluzione,\footnote{dette funzioni sono presenti nella \acr{glibc} 2.1.96,
  ma essendo considerate deprecate (vedi
  sez.~\ref{sec:sock_advanced_name_services}) sono state rimosse nelle
  versioni successive.} \funcd{getipnodebyname} e \funcd{getipnodebyaddr}, i
cui prototipi sono:

\begin{funcproto}{
\fhead{netdb.h}
\fhead{sys/types.h}
\fhead{sys/socket.h}
\fdecl{struct hostent *getipnodebyname(const char *name, int af, int
    flags, int *error\_num)} 
\fdesc{Richiede la risoluzione di un nome a dominio.} 
\fdecl{struct hostent *getipnodebyaddr(const void *addr, size\_t len,
    int af, int *error\_num)}
\fdesc{Richiede la risoluzione inversa di un indirizzo IP.} 
}

{Le funzioni ritornano l'indirizzo ad una struttura \struct{hostent} in caso
  di successo e \val{NULL} per un errore.}
\end{funcproto}

Entrambe le funzioni supportano esplicitamente la scelta di una famiglia di
indirizzi con l'argomento \param{af} (che può assumere i valori
\const{AF\_INET} o \const{AF\_INET6}), e restituiscono un codice di errore
(con valori identici a quelli precedentemente illustrati in
tab.~\ref{tab:h_errno_values}) nella variabile puntata da \param{error\_num}.
La funzione \func{getipnodebyaddr} richiede poi che si specifichi l'indirizzo
come per \func{gethostbyaddr} passando anche la lunghezza dello stesso
nell'argomento \param{len}.

La funzione \func{getipnodebyname} prende come primo argomento il nome da
risolvere, inoltre prevede un apposito argomento \param{flags}, da usare come
maschera binaria, che permette di specificarne il comportamento nella
risoluzione dei diversi tipi di indirizzi (IPv4 e IPv6); ciascun bit
dell'argomento esprime una diversa opzione, e queste possono essere specificate
con un OR aritmetico delle costanti riportate in
tab.~\ref{tab:sock_getipnodebyname_flags}.

\begin{table}[!htb]
  \centering
  \footnotesize
  \begin{tabular}[c]{|l|p{8cm}|}
    \hline
    \textbf{Costante} & \textbf{Significato} \\
    \hline
    \hline
    \constd{AI\_V4MAPPED}  & Usato con \const{AF\_INET6} per richiedere una
                             ricerca su un indirizzo IPv4 invece che IPv6; gli
                             eventuali risultati saranno rimappati su indirizzi 
                             IPv6.\\
    \constd{AI\_ALL}       & Usato con \const{AI\_V4MAPPED}; richiede sia
                             indirizzi IPv4 che IPv6, e gli indirizzi IPv4
                             saranno rimappati in IPv6.\\
    \constd{AI\_ADDRCONFIG}& Richiede che una richiesta IPv4 o IPv6 venga
                             eseguita solo se almeno una interfaccia del
                             sistema è associata ad un indirizzo di tale tipo.\\
    \constd{AI\_DEFAULT}   & Il valore di default, è equivalente alla
                             combinazione di \const{AI\_ADDRCONFIG} e di
                             \const{AI\_V4MAPPED}.\\  
    \hline
  \end{tabular}
  \caption{Valori possibili per i bit dell'argomento \param{flags} della
    funzione \func{getipnodebyname}.}
  \label{tab:sock_getipnodebyname_flags}
\end{table}

Entrambe le funzioni restituiscono un puntatore ad una struttura \var{hostent}
che contiene i risultati della ricerca, che viene allocata dinamicamente
insieme a tutto lo spazio necessario a contenere i dati in essa referenziati;
per questo motivo queste funzioni non soffrono dei problemi dovuti all'uso di
una sezione statica di memoria presenti con le precedenti \func{gethostbyname}
e \func{gethostbyaddr}.  L'uso di una allocazione dinamica però comporta anche
la necessità di disallocare esplicitamente la memoria occupata dai risultati
una volta che questi non siano più necessari; a tale scopo viene fornita la
funzione \funcd{freehostent}, il cui prototipo è:

\begin{funcproto}{
\fhead{netdb.h}
\fhead{sys/types.h}
\fhead{sys/socket.h}
\fdecl{void freehostent(struct hostent *ip)}
\fdesc{Disalloca una struttura \var{hostent}.} 
}

{La funzione non ritorna nulla, e non da errori.}
\end{funcproto}

La funzione permette di disallocare una struttura \var{hostent}
precedentemente allocata in una chiamata di \func{getipnodebyname} o
\func{getipnodebyaddr}, e prende come argomento l'indirizzo restituito da una
di queste funzioni.

Infine per concludere la nostra panoramica sulle funzioni di risoluzione dei
nomi dobbiamo citare le funzioni che permettono di interrogare gli altri
servizi di risoluzione dei nomi illustrati in sez.~\ref{sec:sock_resolver}; in
generale infatti ci sono una serie di funzioni nella forma
\texttt{get\textsl{XXX}byname} e \texttt{get\textsl{XXX}byaddr} (dove
\texttt{\textsl{XXX}} indica il servizio) per ciascuna delle informazioni di
rete mantenute dal \textit{Name Service Switch} che permettono rispettivamente
di trovare una corrispondenza cercando per nome o per numero.

L'elenco di queste funzioni è riportato nelle colonne finali di
tab.~\ref{tab:name_resolution_functions}, dove le si sono suddivise rispetto
al tipo di informazione che forniscono (riportato in prima colonna). Nella
tabella si è anche riportato il file su cui vengono ordinariamente mantenute
queste informazioni, che però può essere sostituito da un qualunque supporto
interno al \textit{Name Service Switch} (anche se usualmente questo avviene
solo per la risoluzione degli indirizzi).  Ciascuna funzione fa riferimento ad
una sua apposita struttura che contiene i relativi dati, riportata in terza
colonna.

\begin{table}[!htb]
  \centering
  \footnotesize
  \begin{tabular}[c]{|l|l|l|l|l|}
    \hline
    \textbf{Informazione}&\textbf{File}&\textbf{Struttura}&
    \multicolumn{2}{|c|}{\textbf{Funzioni}}\\
    \hline
    \hline
    indirizzo &\conffile{/etc/hosts}&\struct{hostent}&\func{gethostbyname}&
               \func{gethostbyaddr}\\ 
    servizio  &\conffile{/etc/services}&\struct{servent}&\func{getservbyname}&
               \func{getservbyport}\\ 
    rete      &\conffile{/etc/networks}&\struct{netent}&\funcm{getnetbyname}&
               \funcm{getnetbyaddr}\\ 
    protocollo&\conffile{/etc/protocols}&\struct{protoent}&
               \funcm{getprotobyname}&\funcm{getprotobyaddr}\\ 
    \hline
  \end{tabular}
  \caption{Funzioni di risoluzione dei nomi per i vari servizi del
    \textit{Name Service Switch} riguardanti la rete.}
  \label{tab:name_resolution_functions}
\end{table}

Delle funzioni di tab.~\ref{tab:name_resolution_functions} abbiamo trattato
finora soltanto quelle relative alla risoluzione dei nomi, dato che sono le
più usate, e prevedono praticamente da sempre la necessità di rivolgersi ad
una entità esterna; per le altre invece, estensioni fornite dal \textit{Name
  Service Switch} a parte, si fa sempre riferimento ai dati mantenuti nei
rispettivi file.

Dopo la risoluzione dei nomi a dominio una delle ricerche più comuni è quella
sui nomi dei servizi di rete più comuni (cioè \texttt{http}, \texttt{smtp},
ecc.) da associare alle rispettive porte. Le due funzioni da utilizzare per
questo sono \funcd{getservbyname} e \funcd{getservbyport}, che permettono
rispettivamente di ottenere il numero di porta associato ad un servizio dato
il nome e viceversa; i loro prototipi sono:



\begin{funcproto}{
\fhead{netdb.h} 
\fdecl{struct servent *getservbyname(const char *name, const char *proto)}
\fdecl{struct servent *getservbyport(int port, const char *proto)} 
\fdesc{Risolvono il nome di un servizio nel rispettivo numero di porta e viceversa.} 
}

{Le funzioni ritornano il puntatore ad una struttura \struct{servent} con i
    risultati in caso di successo e \val{NULL} per un errore.}
\end{funcproto}

Entrambe le funzioni prendono come ultimo argomento una stringa \param{proto}
che indica il protocollo per il quale si intende effettuare la ricerca (le
informazioni mantenute in \conffile{/etc/services} infatti sono relative sia
alle porte usate su UDP che su TCP, occorre quindi specificare a quale dei due
protocolli si fa riferimento) che nel caso di IP può avere come valori
possibili solo \texttt{udp} o \texttt{tcp};\footnote{in teoria si potrebbe
  avere un qualunque protocollo fra quelli citati in
  \conffile{/etc/protocols}, posto che lo stesso supporti il concetto di
  \textsl{porta}, in pratica questi due sono gli unici presenti.}  se si
specifica un puntatore nullo la ricerca sarà eseguita su un protocollo
qualsiasi.

Il primo argomento è il nome del servizio per \func{getservbyname},
specificato tramite la stringa \param{name}, mentre \func{getservbyport}
richiede il numero di porta in \param{port}. Entrambe le funzioni eseguono una
ricerca sul file \conffile{/etc/services}\footnote{il \textit{Name Service
    Switch} astrae il concetto a qualunque supporto su cui si possano
  mantenere i suddetti dati.}  ed estraggono i dati dalla prima riga che
corrisponde agli argomenti specificati; se la risoluzione ha successo viene
restituito un puntatore ad una apposita struttura \struct{servent} contenente
tutti i risultati, altrimenti viene restituito un puntatore nullo.  Si tenga
presente che anche in questo caso i dati vengono mantenuti in una area di
memoria statica e che quindi la funzione non è rientrante.

\begin{figure}[!htb]
  \footnotesize \centering
  \begin{minipage}[c]{0.80\textwidth}
    \includestruct{listati/servent.h}
  \end{minipage}
  \caption{La struttura \structd{servent} per la risoluzione dei nomi dei
    servizi e dei numeri di porta.}
  \label{fig:sock_servent_struct}
\end{figure}

La definizione della struttura \struct{servent} è riportata in
fig.~\ref{fig:sock_servent_struct}, il primo campo, \var{s\_name} contiene
sempre il nome canonico del servizio, mentre \var{s\_aliases} è un puntatore
ad un vettore di stringhe contenenti gli eventuali nomi alternativi
utilizzabili per identificare lo stesso servizio. Infine \var{s\_port}
contiene il numero di porta e \var{s\_proto} il nome del protocollo.

Come riportato in tab.~\ref{tab:name_resolution_functions} ci sono analoghe
funzioni per la risoluzione del nome dei protocolli e delle reti; non staremo
a descriverle nei dettagli, in quanto il loro uso è molto limitato, esse
comunque utilizzano una loro struttura dedicata del tutto analoga alle
precedenti: tutti i dettagli relativi al loro funzionamento possono essere
trovati nelle rispettive pagine di manuale.

Oltre alle funzioni di ricerca esistono delle ulteriori funzioni che prevedono
una lettura sequenziale delle informazioni mantenute nel \textit{Name Service
  Switch} (in sostanza permettono di leggere i file contenenti le informazioni
riga per riga), che sono analoghe a quelle elencate in
tab.~\ref{tab:sys_passwd_func} per le informazioni relative ai dati degli
utenti e dei gruppi. Nel caso specifico dei servizi avremo allora le tre
funzioni \funcd{setservent}, \funcd{getservent} e \funcd{endservent} i cui
prototipi sono:

\begin{funcproto}{
\fhead{netdb.h} 
\fdecl{struct servent *getservent(void)}
\fdesc{Legge la voce successiva nel file \conffile{/etc/services}.} 
\fdecl{void setservent(int stayopen)} 
\fdesc{Apre il file \conffile{/etc/services} e si posiziona al suo inizio.} 
\fdecl{void endservent(void)}
\fdesc{Chiude il file \conffile{/etc/services}.} 
}

{Le due funzioni \func{setservent} e \func{endservent} non ritornano nulla,
  \func{getservent} restituisce il puntatore ad una struttura \struct{servent}
  in caso di successo e \val{NULL} per un errore o fine del file.}
\end{funcproto}

La prima funzione, \func{getservent}, legge una singola voce a partire dalla
posizione corrente in \conffile{/etc/services}, pertanto si può eseguire una
lettura sequenziale dello stesso invocandola più volte. Se il file non è
aperto provvede automaticamente ad aprirlo, nel qual caso leggerà la prima
voce. La seconda funzione, \func{setservent}, permette di aprire il file
\conffile{/etc/services} per una successiva lettura, ma se il file è già stato
aperto riporta la posizione di lettura alla prima voce del file, in questo
modo si può far ricominciare da capo una lettura sequenziale. 

\begin{table}[!htb]
  \centering
  \footnotesize
  \begin{tabular}[c]{|l|l|l|l|}
    \hline
    \textbf{Informazione}&\multicolumn{3}{|c|}{\textbf{Funzioni}}\\
    \hline
    \hline
    indirizzo &\func{sethostent} &\func{gethostent} &\func{endhostent} \\
    servizio  &\func{setservent} &\func{getservent} &\func{endservent}\\ 
    rete      &\funcm{setnetent}  &\funcm{getnetent}  &\funcm{endnetent}\\ 
    protocollo&\funcm{setprotoent}&\funcm{getprotoent}&\funcm{endprotoent}\\ 
    \hline
  \end{tabular}
  \caption{Funzioni lettura sequenziale dei dati del
    \textit{Name Service Switch}.} 
  \label{tab:name_sequential_read}
\end{table}

L'argomento \param{stayopen} di \func{setservent}, se diverso da zero, fa sì
che il file resti aperto anche fra diverse chiamate a \func{getservbyname} e
\func{getservbyport}; di default dopo una chiamata a queste funzioni il file
viene chiuso, cosicché una successiva chiamata a \func{getservent} riparte
dall'inizio.  La terza funzione, \func{endservent}, provvede semplicemente a
chiudere il file.

Queste tre funzioni per la lettura sequenziale di nuovo sono presenti per
ciascuno dei vari tipi di informazione relative alle reti di
tab.~\ref{tab:name_resolution_functions}; questo significa che esistono
altrettante funzioni nella forma \texttt{setXXXent}, \texttt{getXXXent} e
\texttt{endXXXent}, analoghe alle precedenti per la risoluzione dei servizi,
che abbiamo riportato in tab.~\ref{tab:name_sequential_read}.  Essendo, a
parte il tipo di informazione che viene trattato, sostanzialmente identiche
nel funzionamento e di scarso utilizzo, non staremo a trattarle una per una,
rimandando alle rispettive pagine di manuale.


\subsection{Le funzioni avanzate per la risoluzione dei nomi}
\label{sec:sock_advanced_name_services}

Quelle illustrate nella sezione precedente sono le funzioni classiche per la
risoluzione di nomi ed indirizzi IP, ma abbiamo già visto come esse soffrano
di vari inconvenienti come il fatto che usano informazioni statiche, e non
prevedono la possibilità di avere diverse classi di indirizzi. Anche se sono
state create delle estensioni o metodi diversi che permettono di risolvere
alcuni di questi inconvenienti, comunque esse non forniscono una interfaccia
sufficientemente generica.\footnote{rimane ad esempio il problema generico che
  si deve sapere in anticipo quale tipo di indirizzi IP (IPv4 o IPv6)
  corrispondono ad un certo nome a dominio.}

Inoltre in genere quando si ha a che fare con i socket non esiste soltanto il
problema della risoluzione del nome che identifica la macchina, ma anche
quello del servizio a cui ci si vuole rivolgere.  Per questo motivo con lo
standard POSIX 1003.1-2001 sono state indicate come deprecate le varie
funzioni \func{gethostbyaddr}, \func{gethostbyname}, \var{getipnodebyname} e
\var{getipnodebyaddr} ed è stata introdotta una interfaccia completamente
nuova.

La prima funzione di questa interfaccia è \funcd{getaddrinfo}, che combina le
funzionalità delle precedenti \func{getipnodebyname}, \func{getipnodebyaddr},
\func{getservbyname} e \func{getservbyport}, consentendo di ottenere
contemporaneamente sia la risoluzione di un indirizzo simbolico che del nome
di un servizio; la funzione è stata introdotta, insieme a \func{getnameinfo}
che vedremo più avanti,
nell'\href{http://www.ietf.org/rfc/rfc2553.txt}{RFC~2553} ed il suo prototipo è:


\begin{funcproto}{
\fhead{netdb.h}
\fhead{sys/socket.h} 
\fdecl{int getaddrinfo(const char *node, const char *service, const
    struct addrinfo *hints, \\
\phantom{int getaddrinfo(}struct addrinfo **res)}
\fdesc{Esegue una risoluzione di un nome a dominio e di un nome di servizio.} 
}

{La funzione ritorna $0$ in caso di successo e un codice di errore diverso da
  zero per un errore.}
\end{funcproto}

La funzione prende come primo argomento il nome della macchina che si vuole
risolvere, specificato tramite la stringa \param{node}. Questo argomento,
oltre ad un comune nome a dominio, può indicare anche un indirizzo numerico in
forma \textit{dotted-decimal} per IPv4 o in formato esadecimale per IPv6.  Si
può anche specificare il nome di una rete invece che di una singola macchina.
Il secondo argomento, \param{service}, specifica invece il nome del servizio
che si intende risolvere. Per uno dei due argomenti si può anche usare il
valore \val{NULL}, nel qual caso la risoluzione verrà effettuata soltanto
sulla base del valore dell'altro.

Il terzo argomento, \param{hints}, deve essere invece un puntatore ad una
struttura \struct{addrinfo} usata per dare dei \textsl{suggerimenti} al
procedimento di risoluzione riguardo al protocollo o del tipo di socket che si
intenderà utilizzare; \func{getaddrinfo} infatti permette di effettuare
ricerche generiche sugli indirizzi, usando sia IPv4 che IPv6, e richiedere
risoluzioni sui nomi dei servizi indipendentemente dal protocollo (ad esempio
TCP o UDP) che questi possono utilizzare.

Come ultimo argomento in \param{res} deve essere passato un puntatore ad una
variabile (di tipo puntatore ad una struttura \struct{addrinfo}) che verrà
utilizzata dalla funzione per riportare (come \textit{value result argument})
i propri risultati. La funzione infatti è rientrante, ed alloca autonomamente
tutta la memoria necessaria in cui verranno riportati i risultati della
risoluzione.  La funzione scriverà all'indirizzo puntato da \param{res} il
puntatore iniziale ad una \textit{linked list} di strutture di tipo
\struct{addrinfo} contenenti tutte le informazioni ottenute.

\begin{figure}[!htb]
  \footnotesize \centering
  \begin{minipage}[c]{0.90\textwidth}
    \includestruct{listati/addrinfo.h}
  \end{minipage}
  \caption{La struttura \structd{addrinfo} usata nella nuova interfaccia POSIX
    per la risoluzione di nomi a dominio e servizi.}
  \label{fig:sock_addrinfo_struct}
\end{figure}

Come illustrato la struttura \struct{addrinfo}, la cui definizione è riportata
in fig.~\ref{fig:sock_addrinfo_struct}, viene usata sia in ingresso, per
passare dei valori di controllo alla funzione, che in uscita, per ricevere i
risultati. La definizione è ripresa direttamente dal file \headfiled{netdb.h}
in cui questa struttura viene dichiarata, la pagina di manuale riporta
\type{size\_t} come tipo di dato per il campo \var{ai\_addrlen}, qui viene
usata quanto previsto dallo standard POSIX, in cui viene utilizzato
\type{socklen\_t}; i due tipi di dati sono comunque equivalenti.

Il primo campo, \var{ai\_flags}, è una maschera binaria di bit che
permettono di controllare le varie modalità di risoluzione degli indirizzi,
che viene usato soltanto in ingresso. I tre campi successivi \var{ai\_family},
\var{ai\_socktype}, e \var{ai\_protocol} contengono rispettivamente la
famiglia di indirizzi, il tipo di socket e il protocollo, in ingresso vengono
usati per impostare una selezione (impostandone il valore nella struttura
puntata da \param{hints}), mentre in uscita indicano il tipo di risultato
contenuto nella struttura.

Tutti i campi seguenti vengono usati soltanto in uscita e devono essere nulli
o \val{NULL} in ingresso; il campo \var{ai\_addrlen} indica la dimensione
della struttura degli indirizzi ottenuta come risultato, il cui contenuto sarà
memorizzato nella struttura \struct{sockaddr} posta all'indirizzo puntato dal
campo \var{ai\_addr}. Il campo \var{ai\_canonname} è un puntatore alla stringa
contenente il nome canonico della macchina, ed infine, quando la funzione
restituisce più di un risultato, \var{ai\_next} è un puntatore alla successiva
struttura \struct{addrinfo} della lista.

Ovviamente non è necessario dare dei suggerimenti in ingresso, ed usando
\val{NULL} come valore per l'argomento \param{hints} si possono compiere
ricerche generiche.  Se però si specifica un valore non nullo questo deve
puntare ad una struttura \struct{addrinfo} precedentemente allocata nella
quale siano stati opportunamente impostati i valori dei campi
\var{ai\_family}, \var{ai\_socktype}, \var{ai\_protocol} ed \var{ai\_flags}.

I due campi \var{ai\_family} e \var{ai\_socktype} prendono gli stessi valori
degli analoghi argomenti della funzione \func{socket}; in particolare per
\var{ai\_family} si possono usare i valori di tab.~\ref{tab:net_pf_names} ma
sono presi in considerazione solo \const{AF\_INET} e \const{AF\_INET6}, mentre
se non si vuole specificare nessuna famiglia di indirizzi si può usare il
valore \const{AF\_UNSPEC}.  Allo stesso modo per \var{ai\_socktype} si possono
usare i valori illustrati in sez.~\ref{sec:sock_type} per indicare per quale
tipo di socket si vuole risolvere il servizio indicato, anche se i soli
significativi sono \const{SOCK\_STREAM} e \const{SOCK\_DGRAM}; in questo caso,
se non si vuole effettuare nessuna risoluzione specifica, si potrà usare un
valore nullo.

Il campo \var{ai\_protocol} permette invece di effettuare la selezione dei
risultati per il nome del servizio usando il numero identificativo del
rispettivo protocollo di trasporto (i cui valori possibili sono riportati in
\conffile{/etc/protocols}); di nuovo i due soli valori utilizzabili sono quelli
relativi a UDP e TCP, o il valore nullo che indica di ignorare questo campo
nella selezione.

\begin{table}[!htb]
  \centering
  \footnotesize
  \begin{tabular}[c]{|l|p{8cm}|}
    \hline
    \textbf{Costante} & \textbf{Significato} \\
    \hline
    \hline
    \const{AI\_ADDRCONFIG} & Stesso significato dell'analoga di
                             tab.~\ref{tab:sock_getipnodebyname_flags}.\\ 
    \const{AI\_ALL}        & Stesso significato dell'analoga di
                             tab.~\ref{tab:sock_getipnodebyname_flags}.\\ 
    \constd{AI\_CANONNAME}  & Richiede la restituzione del nome canonico della
                              macchina, che verrà salvato in una stringa il cui
                              indirizzo sarà restituito nel campo
                              \var{ai\_canonname} della prima struttura
                              \struct{addrinfo} dei risultati. Se il nome
                              canonico non è disponibile al suo posto
                              viene restituita una copia di \param{node}. \\ 
    \constd{AI\_NUMERICHOST}& Se impostato il nome della macchina specificato
                              con \param{node} deve essere espresso in forma
                              numerica, altrimenti sarà restituito un errore
                              \const{EAI\_NONAME} (vedi
                              tab.~\ref{tab:addrinfo_error_code}), in questo
                              modo si evita ogni chiamata alle funzioni di
                              risoluzione.\\ 
    \constd{AI\_NUMERICSERVICE}& Analogo di \const{AI\_NUMERICHOST} per la
                                 risoluzione di un servizio, con 
                                 \param{service} che deve essere espresso in forma
                                 numerica.\\ 
    \constd{AI\_PASSIVE}    & Viene utilizzato per ottenere un indirizzo in
                              formato adatto per una successiva chiamata a
                              \func{bind}. Se specificato quando si è usato 
                              \val{NULL} come valore per \param{node} gli
                              indirizzi restituiti saranno inizializzati al
                              valore generico (\const{INADDR\_ANY} per IPv4 e
                              \const{IN6ADDR\_ANY\_INIT} per IPv6), altrimenti
                              verrà usato l'indirizzo dell'interfaccia di
                              \textit{loopback}. Se invece non è impostato gli
                              indirizzi verranno restituiti in formato adatto ad
                              una chiamata a \func{connect} o \func{sendto}.\\
    \const{AI\_V4MAPPED}   & Stesso significato dell'analoga di
                             tab.~\ref{tab:sock_getipnodebyname_flags}.\\  
    \hline
    \const{AI\_CANONIDN}   & Se il nome canonico richiesto con
                             \const{AI\_CANONNAME} è codificato con questo
                             flag la codifica viene convertita in forma
                             leggibile nella localizzazione corrente.\\ 
    \const{AI\_IDN}        & Se specificato il nome viene convertito, se
                             necessario, nella codifica IDN, usando la
                             localizzazione corrente.\\ 
    \const{AI\_IDN\_ALLOW\_UNASSIGNED} & attiva il controllo 
                                         \texttt{IDNA\_ALLOW\_UNASSIGNED}.\\ 
    \const{AI\_AI\_IDN\_USE\_STD3\_ASCII\_RULES} & attiva il controllo
                                            \texttt{IDNA\_USE\_STD3\_ASCII\_RULES}\\ 
    \hline
  \end{tabular}
  \caption{Costanti associate ai bit del campo \var{ai\_flags} della struttura 
    \struct{addrinfo}.} 
  \label{tab:ai_flags_values}
\end{table}

% TODO mettere riferimento a IDNA_ALLOW_UNASSIGNED e IDNA_USE_STD3_ASCII_RULES

Infine gli ultimi dettagli si controllano con il campo \var{ai\_flags}; che
deve essere impostato come una maschera binaria; i bit di questa variabile
infatti vengono usati per dare delle indicazioni sul tipo di risoluzione
voluta, ed hanno valori analoghi a quelli visti in
sez.~\ref{sec:sock_name_services} per \func{getipnodebyname}; il valore di
\var{ai\_flags} può essere impostata con un OR aritmetico delle costanti di
tab.~\ref{tab:ai_flags_values}, ciascuna delle quali identifica un bit della
maschera. 

Nella seconda parte della tabella si sono riportati i valori delle costanti
aggiunte a partire dalla \acr{glibc} 2.3.4 per gestire la
internazionalizazione dei nomi a dominio (IDN o \textit{Internationalized
  Domain Names}) secondo quanto specificato
nell'\href{http://www.ietf.org/rfc/rfc3490.txt}{RFC~3490} (potendo cioè usare
codifiche di caratteri che consentono l'espressione di nomi a dominio in
qualunque lingua).

Come accennato passando un valore \val{NULL} per l'argomento \param{hints} si
effettua una risuluzione generica, equivalente ad aver impostato un valore
nullo per \var{ai\_family} e \var{ai\_socktype}, un valore \const{AF\_UNSPEC}
per \var{ai\_family} e il valore \code{(AI\_V4MAPPED|AI\_ADDRCONFIG)} per
\var{ai\_flags}. 

La funzione restituisce un valore nullo in caso di successo, o un codice in
caso di errore. I valori usati come codice di errore sono riportati in
tab.~\ref{tab:addrinfo_error_code}; dato che la funzione utilizza altre
funzioni e chiamate al sistema per ottenere il suo risultato in generale il
valore di \var{errno} non è significativo, eccetto il caso in cui si sia
ricevuto un errore di \const{EAI\_SYSTEM}, nel qual caso l'errore
corrispondente è riportato tramite \var{errno}.

\begin{table}[!htb]
  \centering
  \footnotesize
  \begin{tabular}[c]{|l|p{10cm}|}
    \hline
    \textbf{Costante} & \textbf{Significato} \\
    \hline
    \hline
    \constd{EAI\_ADDRFAMILY}& La richiesta non ha nessun indirizzo di rete
                              per la famiglia di indirizzi specificata. \\
    \constd{EAI\_AGAIN}   & Il DNS ha restituito un errore di risoluzione  
                            temporaneo, si può ritentare in seguito. \\
    \constd{EAI\_BADFLAGS}& Il campo \var{ai\_flags} contiene dei valori non
                            validi per i flag o si è richiesto
                            \const{AI\_CANONNAME} con \param{name} nullo. \\
    \constd{EAI\_FAIL}    & Il DNS ha restituito un errore di risoluzione  
                            permanente. \\
    \constd{EAI\_FAMILY}  & La famiglia di indirizzi richiesta non è
                            supportata. \\ 
    \constd{EAI\_MEMORY}  & È stato impossibile allocare la memoria necessaria
                            alle operazioni. \\
    \constd{EAI\_NODATA}  & La macchina specificata esiste, ma non ha nessun
                            indirizzo di rete definito. \\
    \constd{EAI\_NONAME}  & Il nome a dominio o il servizio non sono noti,
                            viene usato questo errore anche quando si specifica
                            il valore \val{NULL} per entrambi gli argomenti
                            \param{node} e \param{service}. \\
    \constd{EAI\_SERVICE} & Il servizio richiesto non è disponibile per il tipo
                            di socket richiesto, anche se può esistere per
                            altri tipi di socket. \\
    \constd{EAI\_SOCKTYPE}& Il tipo di socket richiesto non è supportato. \\
    \constd{EAI\_SYSTEM}  & C'è stato un errore di sistema, si può controllare
                            \var{errno} per i dettagli. \\
%    \hline
% TODO estensioni GNU, trovarne la documentazione
%    \constd{EAI\_INPROGRESS}& Richiesta in corso. \\
%    \constd{EAI\_CANCELED}& La richiesta è stata cancellata.\\
%    \constd{EAI\_NOTCANCELED}& La richiesta non è stata cancellata. \\
%    \constd{EAI\_ALLDONE} & Tutte le richieste sono complete. \\
%    \constd{EAI\_INTR}    & Richiesta interrotta. \\
    \hline
  \end{tabular}
  \caption{Costanti associate ai valori dei codici di errore della funzione
    \func{getaddrinfo}.} 
  \label{tab:addrinfo_error_code}
\end{table}

Come per i codici di errore di \func{gethostbyname} anche in questo caso è
fornita una apposita funzione, simile a \func{strerror}, che consente di
utilizzare direttamente il codice restituito dalla funzione per stampare a
video un messaggio esplicativo; la funzione è \funcd{gai\_strerror} ed il suo
prototipo è:

\begin{funcproto}{
\fhead{netdb.h} 
\fdecl{const char *gai\_strerror(int errcode)}
\fdesc{Fornisce il messaggio corrispondente ad un errore di \func{getaddrinfo}.} 
}

{La funzione ritorna il puntatore alla stringa contenente il messaggio di
  errore.}
\end{funcproto}

La funzione restituisce un puntatore alla stringa contenente il messaggio
corrispondente dal codice di errore \param{errcode} ottenuto come valore di
ritorno di \func{getaddrinfo}.  La stringa è allocata staticamente, ma essendo
costante ed accessibile in sola lettura, la funzione è rientrante.

Dato che ad un certo nome a dominio possono corrispondere più indirizzi IP
(sia IPv4 che IPv6), e che un certo servizio può essere fornito su protocolli
e tipi di socket diversi, in generale, a meno di non aver eseguito una
selezione specifica attraverso l'uso di \param{hints}, si otterrà una diversa
struttura \struct{addrinfo} per ciascuna possibilità.  

Ad esempio se si richiede la risoluzione del servizio \textit{echo} per
l'indirizzo \texttt{www.truelite.it}, e si imposta \const{AI\_CANONNAME} per
avere anche la risoluzione del nome canonico, si avrà come risposta della
funzione la lista illustrata in fig.~\ref{fig:sock_addrinfo_list}.

\begin{figure}[!htb]
  \centering
  \includegraphics[width=10cm]{img/addrinfo_list}
  \caption{La \textit{linked list} delle strutture \struct{addrinfo}
    restituite da \func{getaddrinfo}.}
  \label{fig:sock_addrinfo_list}
\end{figure}

Come primo esempio di uso di \func{getaddrinfo} vediamo un programma
elementare di interrogazione del \textit{resolver} basato questa funzione, il
cui corpo principale è riportato in fig.~\ref{fig:mygetaddr_example}. Il
codice completo del programma, compresa la gestione delle opzioni in cui è
gestita l'eventuale inizializzazione dell'argomento \var{hints} per
restringere le ricerche su protocolli, tipi di socket o famiglie di indirizzi,
è disponibile nel file \texttt{mygetaddr.c} dei sorgenti allegati alla guida.

\begin{figure}[!htb]
  \footnotesize \centering
  \begin{minipage}[c]{\codesamplewidth}
    \includecodesample{listati/mygetaddr.c}
  \end{minipage}
  \normalsize
  \caption{Esempio di codice per la risoluzione di un indirizzo.}
  \label{fig:mygetaddr_example}
\end{figure}

Il corpo principale inizia controllando (\texttt{\small 1--5}) il numero di
argomenti passati, che devono essere sempre due, e corrispondere
rispettivamente all'indirizzo ed al nome del servizio da risolvere. A questo
segue la chiamata (\texttt{\small 7}) alla funzione \func{getaddrinfo}, ed il
successivo controllo (\texttt{\small 8--11}) del suo corretto funzionamento,
senza il quale si esce immediatamente stampando il relativo codice di errore.

Se la funzione ha restituito un valore nullo il programma prosegue
inizializzando (\texttt{\small 12}) il puntatore \var{ptr} che sarà usato nel
successivo ciclo (\texttt{\small 14--35}) di scansione della lista delle
strutture \struct{addrinfo} restituite dalla funzione. Prima di eseguire
questa scansione (\texttt{\small 12}) viene stampato il valore del nome
canonico che è presente solo nella prima struttura.

La scansione viene ripetuta (\texttt{\small 14}) fintanto che si ha un
puntatore valido. La selezione principale è fatta sul campo \var{ai\_family},
che stabilisce a quale famiglia di indirizzi fa riferimento la struttura in
esame. Le possibilità sono due, un indirizzo IPv4 o IPv6, se nessuna delle due
si verifica si provvede (\texttt{\small 27--30}) a stampare un messaggio di
errore ed uscire (questa eventualità non dovrebbe comunque mai verificarsi,
almeno fintanto che la funzione \func{getaddrinfo} lavora correttamente).

Per ciascuno delle due possibili famiglie di indirizzi si estraggono le
informazioni che poi verranno stampate alla fine del ciclo (\texttt{\small
  31--34}). Il primo caso esaminato (\texttt{\small 15--21}) è quello degli
indirizzi IPv4, nel qual caso prima se ne stampa l'identificazione
(\texttt{\small 16}) poi si provvede a ricavare la struttura degli indirizzi
(\texttt{\small 17}) indirizzata dal campo \var{ai\_addr}, eseguendo un
opportuno casting del puntatore per poter estrarre da questa la porta
(\texttt{\small 18}) e poi l'indirizzo (\texttt{\small 19}) che verrà
convertito con una chiamata ad \func{inet\_ntop}.

La stessa operazione (\texttt{\small 21--27}) viene ripetuta per gli indirizzi
IPv6, usando la rispettiva struttura degli indirizzi. Si noti anche come in
entrambi i casi per la chiamata a \func{inet\_ntop} si sia dovuto passare il
puntatore al campo contenente l'indirizzo IP nella struttura puntata dal campo
\var{ai\_addr}.\footnote{il meccanismo è complesso a causa del fatto che al
  contrario di IPv4, in cui l'indirizzo IP può essere espresso con un semplice
  numero intero, in IPv6 questo deve essere necessariamente fornito come
  struttura, e pertanto anche se nella struttura puntata da \var{ai\_addr}
  sono presenti direttamente i valori finali, per l'uso con \func{inet\_ntop}
  occorre comunque passare un puntatore agli stessi (ed il costrutto
  \code{\&addr6->sin6\_addr} è corretto in quanto l'operatore \texttt{->} ha
  in questo caso precedenza su \texttt{\&}).}

Una volta estratte dalla struttura \struct{addrinfo} tutte le informazioni
relative alla risoluzione richiesta e stampati i relativi valori, l'ultimo
passo (\texttt{\small 34}) è di estrarre da \var{ai\_next} l'indirizzo della
eventuale successiva struttura presente nella lista e ripetere il ciclo, fin
tanto che, completata la scansione, questo avrà un valore nullo e si potrà
terminare (\texttt{\small 36}) il programma.

Si tenga presente che \func{getaddrinfo} non garantisce nessun particolare
ordinamento della lista delle strutture \struct{addrinfo} restituite, anche se
usualmente i vari indirizzi IP (se ne è presente più di uno) sono forniti
nello stesso ordine in cui vengono inviati dal server DNS. In particolare
nulla garantisce che vengano forniti prima i dati relativi ai servizi di un
determinato protocollo o tipo di socket, se ne sono presenti di diversi.  Se
allora utilizziamo il nostro programma potremo verificare il risultato:
\begin{Console}
[piccardi@gont sources]$ \textbf{./mygetaddr -c  gapil.truelite.it echo}
Canonical name sources2.truelite.it
IPv4 address:
        Indirizzo 62.48.34.25
        Protocollo 6
        Porta 7
IPv4 address:
        Indirizzo 62.48.34.25
        Protocollo 17
        Porta 7
\end{Console}
%$

Una volta estratti i risultati dalla \textit{linked list} puntata
da \param{res} se questa non viene più utilizzata si dovrà avere cura di
disallocare opportunamente tutta la memoria, per questo viene fornita
l'apposita funzione \funcd{freeaddrinfo}, il cui prototipo è:


\begin{funcproto}{
\fhead{netdb.h}
\fdecl{void freeaddrinfo(struct addrinfo *res)}
\fdesc{Libera la memoria allocata da una precedente chiamata a \func{getaddrinfo}.} 
}

{La funzione non restituisce nessun codice di errore.}
\end{funcproto}

La funzione prende come unico argomento il puntatore \param{res}, ottenuto da
una precedente chiamata a \func{getaddrinfo}, e scandisce la lista delle
strutture per liberare tutta la memoria allocata. Dato che la funzione non ha
valori di ritorno deve essere posta molta cura nel passare un valore valido
per \param{res} ed usare un indirizzo non valido o già liberato può avere
conseguenze non prevedibili.

Si tenga presente infine che se si copiano i risultati da una delle strutture
\struct{addrinfo} restituite nella lista indicizzata da \param{res}, occorre
avere cura di eseguire una \textit{deep copy} in cui si copiano anche tutti i
dati presenti agli indirizzi contenuti nella struttura \struct{addrinfo},
perché una volta disallocati i dati con \func{freeaddrinfo} questi non
sarebbero più disponibili.

Anche la nuova interfaccia definita da POSIX prevede una nuova funzione per
eseguire la risoluzione inversa e determinare nomi di servizi e di dominio
dati i rispettivi valori numerici. La funzione che sostituisce le varie
\func{gethostbyname}, \func{getipnodebyname} e \func{getservbyname} è
\funcd{getnameinfo}, ed il suo prototipo è:

\begin{funcproto}{
\fhead{netdb.h}
\fhead{sys/socket.h}
\fdecl{int getnameinfo(const struct sockaddr *sa, socklen\_t salen, \\
\phantom{int getnameinfo(}char *host, size\_t hostlen, char *serv, size\_t
servlen, int flags)} 
\fdesc{Effettua una risoluzione di un indirizzo di rete in maniera
  indipendente dal protocollo.} 
}

{La funzione ritorna $0$ in caso di successo e un codice di errore diverso da
  zero per un errore.}
\end{funcproto}

La principale caratteristica di \func{getnameinfo} è che la funzione è in
grado di eseguire una risoluzione inversa in maniera indipendente dal
protocollo; il suo primo argomento \param{sa} infatti è il puntatore ad una
struttura degli indirizzi generica, che può contenere sia indirizzi IPv4 che
IPv6, la cui dimensione deve comunque essere specificata con l'argomento
\param{salen}. 

I risultati della funzione saranno restituiti nelle due stringhe puntate da
\param{host} e \param{serv}, che dovranno essere state precedentemente
allocate per una lunghezza massima che deve essere specificata con gli altri
due argomenti \param{hostlen} e \param{servlen}. Quando non si è
interessati ad uno dei due, si può passare il valore \val{NULL} come argomento,
così che la corrispondente informazione non venga richiesta. Infine l'ultimo
argomento \param{flags} è una maschera binaria i cui bit consentono di
impostare le modalità con cui viene eseguita la ricerca, e deve essere
specificato attraverso l'OR aritmetico dei valori illustrati in
tab.~\ref{tab:getnameinfo_flags}, nella seconda parte della tabella si sono
aggiunti i valori introdotto con la \acr{glibc} 2.3.4 per gestire la
internazionalizzione dei nomi a dominio.

\begin{table}[!htb]
  \centering
  \footnotesize
  \begin{tabular}[c]{|l|p{8cm}|}
    \hline
    \textbf{Costante} & \textbf{Significato} \\
    \hline
    \hline
    \constd{NI\_DGRAM}      & Richiede che venga restituito il nome del
                              servizio su UDP invece che quello su TCP per quei
                              pichi servizi (porte 512-214) che soni diversi
                              nei due protocolli.\\
    \constd{NI\_NOFQDN}     & Richiede che venga restituita solo il nome della
                              macchina all'interno del dominio al posto del
                              nome completo (FQDN).\\
    \constd{NI\_NAMEREQD}   & Richiede la restituzione di un errore se il nome
                              non può essere risolto.\\
    \constd{NI\_NUMERICHOST}& Richiede che venga restituita la forma numerica
                              dell'indirizzo (questo succede sempre se il nome
                              non può essere ottenuto).\\ 
    \constd{NI\_NUMERICSERV}& Richiede che il servizio venga restituito in
                              forma numerica (attraverso il numero di
                              porta).\\
    \hline
    \const{NI\_IDN}        & Se specificato il nome restituito viene convertito usando la
                             localizzazione corrente, se necessario, nella
                             codifica IDN.\\ 
    \const{NI\_IDN\_ALLOW\_UNASSIGNED} & attiva il controllo 
                                       \texttt{IDNA\_ALLOW\_UNASSIGNED}.\\ 
    \const{NI\_AI\_IDN\_USE\_STD3\_ASCII\_RULES} & attiva il controllo
                                       \texttt{IDNA\_USE\_STD3\_ASCII\_RULES}\\ 
    \hline
  \end{tabular}
  \caption{Costanti associate ai bit dell'argomento \param{flags} della  
    funzione \func{getnameinfo}.} 
  \label{tab:getnameinfo_flags}
\end{table}

La funzione ritorna zero in caso di successo, e scrive i propri risultati agli
indirizzi indicati dagli argomenti \param{host} e \param{serv} come stringhe
terminate dal carattere NUL, a meno che queste non debbano essere troncate
qualora la loro dimensione ecceda quelle specificate dagli argomenti
\param{hostlen} e \param{servlen}. Sono comunque definite le due costanti
\constd{NI\_MAXHOST} e \constd{NI\_MAXSERV}\footnote{in Linux le due costanti
  sono definite in \headfile{netdb.h} ed hanno rispettivamente il valore 1024
  e 12.}  che possono essere utilizzate come limiti massimi.  In caso di
errore viene restituito invece un codice che assume gli stessi valori
illustrati in tab.~\ref{tab:addrinfo_error_code}.

A questo punto possiamo fornire degli esempi di utilizzo della nuova
interfaccia, adottandola per le precedenti implementazioni del client e del
server per il servizio \textit{echo}; dato che l'uso delle funzioni appena
illustrate (in particolare di \func{getaddrinfo}) è piuttosto complesso,
essendo necessaria anche una impostazione diretta dei campi dell'argomento
\param{hints}, provvederemo una interfaccia semplificata per i due casi visti
finora, quello in cui si specifica nel client un indirizzo remoto per la
connessione al server, e quello in cui si specifica nel server un indirizzo
locale su cui porsi in ascolto.

La prima funzione della nostra interfaccia semplificata è \texttt{sockconn}
che permette di ottenere un socket, connesso all'indirizzo ed al servizio
specificati. Il corpo della funzione è riportato in
fig.~\ref{fig:sockconn_code}, il codice completo è nel file \file{SockUtil.c}
dei sorgenti allegati alla guida, che contiene varie funzioni di utilità per
l'uso dei socket.

\begin{figure}[!htb]
  \footnotesize \centering
  \begin{minipage}[c]{\codesamplewidth}
    \includecodesample{listati/sockconn.c}
  \end{minipage}
  \normalsize
  \caption{Il codice della funzione \texttt{sockconn}.}
  \label{fig:sockconn_code}
\end{figure}

La funzione prende quattro argomenti, i primi due sono le stringhe che
indicano il nome della macchina a cui collegarsi ed il relativo servizio su
cui sarà effettuata la risoluzione; seguono il protocollo da usare (da
specificare con il valore numerico di \conffile{/etc/protocols}) ed il tipo di
socket (al solito specificato con i valori illustrati in
sez.~\ref{sec:sock_type}).  La funzione ritorna il valore del file descriptor
associato al socket (un numero positivo) in caso di successo, o $-1$ in caso
di errore.

Per risolvere il problema di non poter passare indietro i valori di ritorno di
\func{getaddrinfo} contenenti i relativi codici di errore si sono stampati i
messaggi d'errore direttamente nella funzione; infatti non si può avere
nessuna certezza che detti valori siano negativi e per cui stampare subito
l'errore diventa necessario per evitare ogni possibile ambiguità nei confronti
del valore di ritorno in caso di successo.

Una volta definite le variabili occorrenti (\texttt{\small 3--5}) la funzione
prima (\texttt{\small 6}) azzera il contenuto della struttura \var{hint} e poi
provvede (\texttt{\small 7--9}) ad inizializzarne i valori necessari per la
chiamata (\texttt{\small 10}) a \func{getaddrinfo}. Di quest'ultima si
controlla (\texttt{\small 12--16}) il codice di ritorno, in modo da stampare
un avviso di errore, azzerare \var{errno} ed uscire in caso di errore.  

Dato che ad una macchina possono corrispondere più indirizzi IP, e di tipo
diverso (sia IPv4 che IPv6), mentre il servizio può essere in ascolto soltanto
su uno solo di questi, si provvede a tentare la connessione per ciascun
indirizzo restituito all'interno di un ciclo (\texttt{\small 18--40}) di
scansione della lista restituita da \func{getaddrinfo}, ma prima
(\texttt{\small 17}) si salva il valore del puntatore per poterlo riutilizzare
alla fine per disallocare la lista.

Il ciclo viene ripetuto (\texttt{\small 18}) fintanto che si hanno indirizzi
validi, ed inizia (\texttt{\small 19}) con l'apertura del socket; se questa
fallisce si controlla (\texttt{\small 20}) se sono disponibili altri
indirizzi, nel qual caso si passa al successivo (\texttt{\small 21}) e si
riprende (\texttt{\small 22}) il ciclo da capo; se non ve ne sono si stampa
l'errore ritornando immediatamente (\texttt{\small 24--27}). 

Quando la creazione del socket ha avuto successo si procede (\texttt{\small
  29}) direttamente con la connessione, di nuovo in caso di fallimento viene
ripetuto (\texttt{\small 30--38}) il controllo se vi sono o no altri indirizzi
da provare nella stessa modalità fatta in precedenza, aggiungendovi però in
entrambi i casi (\texttt{\small 32} e (\texttt{\small 36}) la chiusura del
socket precedentemente aperto, che non è più utilizzabile.

Se la connessione ha avuto successo invece si termina (\texttt{\small 39})
direttamente il ciclo, e prima di ritornare (\texttt{\small 31}) il valore del
file descriptor del socket si provvede (\texttt{\small 30}) a liberare le
strutture \struct{addrinfo} allocate da \func{getaddrinfo} utilizzando il
valore del relativo puntatore precedentemente (\texttt{\small 17}) salvato.
Si noti come per la funzione sia del tutto irrilevante se la struttura
ritornata contiene indirizzi IPv6 o IPv4, in quanto si fa uso direttamente dei
dati relativi alle strutture degli indirizzi di \struct{addrinfo} che sono
opachi rispetto all'uso della funzione \func{connect}.

Per usare questa funzione possiamo allora modificare ulteriormente il nostro
programma client per il servizio \textit{echo}; in questo caso rispetto al
codice usato finora per collegarsi (vedi fig.~\ref{fig:TCP_echo_client_1})
avremo una semplificazione per cui il corpo principale del nostro client
diventerà quello illustrato in fig.~\ref{fig:TCP_echo_fifth}, in cui le
chiamate a \func{socket}, \func{inet\_pton} e \func{connect} sono sostituite
da una singola chiamata a \texttt{sockconn}. Inoltre il nuovo client (il cui
codice completo è nel file \file{TCP\_echo\_fifth.c} dei sorgenti allegati)
consente di utilizzare come argomento del programma un nome a dominio al posto
dell'indirizzo numerico, e può utilizzare sia indirizzi IPv4 che IPv6.

\begin{figure}[!htb]
  \footnotesize \centering
  \begin{minipage}[c]{\codesamplewidth}
    \includecodesample{listati/TCP_echo_fifth.c}
  \end{minipage}
  \normalsize
  \caption{Il nuovo codice per la connessione del client \textit{echo}.}
  \label{fig:TCP_echo_fifth}
\end{figure}

La seconda funzione di ausilio che abbiamo creato è \texttt{sockbind}, il cui
corpo principale è riportato in fig.~\ref{fig:sockbind_code} (al solito il
sorgente completo è nel file \file{sockbind.c} dei sorgenti allegati alla
guida). Come si può notare la funzione è del tutto analoga alla precedente
\texttt{sockconn}, e prende gli stessi argomenti, però invece di eseguire una
connessione con \func{connect} si limita a chiamare \func{bind} per collegare
il socket ad una porta.

\begin{figure}[!htb]
  \footnotesize \centering
  \begin{minipage}[c]{\codesamplewidth}
    \includecodesample{listati/sockbind.c}
  \end{minipage}
  \normalsize
  \caption{Il codice della funzione \texttt{sockbind}.}
  \label{fig:sockbind_code}
\end{figure}

Dato che la funzione è pensata per essere utilizzata da un server ci si può
chiedere a quale scopo mantenere l'argomento \param{host} quando l'indirizzo
di questo è usualmente noto. Si ricordi però quanto detto in
sez.~\ref{sec:TCP_func_bind}, relativamente al significato della scelta di un
indirizzo specifico come argomento di \func{bind}, che consente di porre il
server in ascolto su uno solo dei possibili diversi indirizzi presenti su di
una macchina.  Se non si vuole che la funzione esegua \func{bind} su un
indirizzo specifico, ma utilizzi l'indirizzo generico, occorrerà avere cura di
passare un valore \val{NULL} come valore per l'argomento \var{host}; l'uso
del valore \const{AI\_PASSIVE} serve ad ottenere il valore generico nella
rispettiva struttura degli indirizzi.

Come già detto la funzione è analoga a \texttt{sockconn} ed inizia azzerando
ed inizializzando (\texttt{\small 6--11}) opportunamente la struttura
\var{hint} con i valori ricevuti come argomenti, soltanto che in questo caso
si è usata (\texttt{\small 8}) una impostazione specifica dei flag di
\var{hint} usando \const{AI\_PASSIVE} per indicare che il socket sarà usato
per una apertura passiva. Per il resto la chiamata (\texttt{\small 12--18}) a
\func{getaddrinfo} e ed il ciclo principale (\texttt{\small 20--42}) sono
identici, solo che si è sostituita (\texttt{\small 31}) la chiamata a
\func{connect} con una chiamata a \func{bind}. Anche la conclusione
(\texttt{\small 43--44}) della funzione è identica.

Si noti come anche in questo caso si siano inserite le stampe degli errori
sullo \textit{standard error}, nonostante la funzione possa essere invocata da
un demone. Nel nostro caso questo non è un problema in quanto se la funzione
non ha successo il programma deve uscire immediatamente prima di essere posto
in background, e può quindi scrivere gli errori direttamente sullo
\textit{standard error}.

\begin{figure}[!htbp]
  \footnotesize \centering
  \begin{minipage}[c]{\codesamplewidth}
    \includecodesample{listati/TCP_echod_third.c}
  \end{minipage}
  \normalsize
  \caption{Nuovo codice per l'apertura passiva del server \textit{echo}.}
  \label{fig:TCP_echod_third}
\end{figure}

Con l'uso di questa funzione si può modificare anche il codice del nostro
server \textit{echo}, che rispetto a quanto illustrato nella versione iniziale
di fig.~\ref{fig:TCP_echo_server_first_code} viene modificato nella forma
riportata in fig.~\ref{fig:TCP_echod_third}. In questo caso il socket su cui
porsi in ascolto viene ottenuto (\texttt{\small 15--18}) da \texttt{sockbind}
che si cura anche della eventuale risoluzione di un indirizzo specifico sul
quale si voglia far ascoltare il server.


\section{Le opzioni dei socket}
\label{sec:sock_options}

Benché dal punto di vista del loro uso come canali di trasmissione di dati i
socket vengano trattati allo stesso modo dei file, acceduti tramite i file
descriptor, e gestiti con le ordinarie funzioni di lettura e scrittura dei
file, l'interfaccia standard usata per la gestione dei file generici non è
comunque sufficiente a controllare la moltitudine di caratteristiche
specifiche che li contraddistinguono, considerato tra l'altro che queste
possono essere completamente diverse fra loro a seconda del tipo di socket e
della relativa forma di comunicazione sottostante.

In questa sezione vedremo allora quali sono le funzioni dedicate alla gestione
delle caratteristiche specifiche dei vari tipi di socket, che vengono
raggruppate sotto il nome generico di ``\textit{socket options}'', ma
soprattutto analizzaremo quali sono queste opzioni e quali caretteristiche e
comportamenti dei socket permettono di controllare.


\subsection{Le funzioni di gestione delle opzioni dei socket}
\label{sec:sock_setsockopt}

La modalità principale con cui si possono gestire le caratteristiche dei
socket (ne vedremo delle ulteriori nelle prossime sezioni) è quella che passa
attraverso l'uso di due funzioni di sistema generiche che permettono
rispettivamente di impostarle e di recuperarne il valore corrente. La prima di
queste due funzioni, quella usata per impostare le \textit{socket options}, è
\funcd{setsockopt}, ed il suo prototipo è:

\begin{funcproto}{
\fhead{sys/socket.h}
\fhead{sys/types.h}
\fdecl{int setsockopt(int sock, int level, int optname, const void
    *optval, socklen\_t optlen)}
\fdesc{Imposta le opzioni di un socket.} 
}

{La funzione ritorna $0$ in caso di successo e $-1$ per un errore, nel qual
  caso \var{errno} assumerà uno dei valori: 
  \begin{errlist}
  \item[\errcode{EBADF}]  il file descriptor \param{sock} non è valido.
  \item[\errcode{EFAULT}] l'indirizzo \param{optval} non è valido.
  \item[\errcode{EINVAL}] il valore di \param{optlen} non è valido.
  \item[\errcode{ENOPROTOOPT}] l'opzione scelta non esiste per il livello
    indicato. 
  \item[\errcode{ENOTSOCK}] il file descriptor \param{sock} non corrisponde ad
    un socket.
  \end{errlist}
}
\end{funcproto}

Il primo argomento della funzione, \param{sock}, indica il socket su cui si
intende operare; per indicare l'opzione da impostare si devono usare i due
argomenti successivi, \param{level} e \param{optname}.  Come abbiamo visto in
sez.~\ref{sec:net_protocols} i protocolli di rete sono strutturati su vari
livelli, ed l'interfaccia dei socket può usarne più di uno. Si avranno allora
funzionalità e caratteristiche diverse per ciascun protocollo usato da un
socket, e quindi saranno anche diverse le opzioni che si potranno impostare
per ciascun socket, a seconda del \textsl{livello} (trasporto, rete, ecc.) su
cui si vuole andare ad operare.

Il valore di \param{level} seleziona allora il protocollo su cui vuole
intervenire, mentre \param{optname} permette di scegliere su quale delle
opzioni che sono definite per quel protocollo si vuole operare. In sostanza la
selezione di una specifica opzione viene fatta attraverso una coppia di valori
\param{level} e \param{optname} e chiaramente la funzione avrà successo
soltanto se il protocollo in questione prevede quella opzione ed è utilizzato
dal socket.  Infine \param{level} prevede anche il valore speciale
\const{SOL\_SOCKET} usato per le opzioni generiche che sono disponibili per
qualunque tipo di socket.

\begin{table}[!htb]
  \centering
  \footnotesize
  \begin{tabular}[c]{|l|l|}
    \hline
    \textbf{Livello} & \textbf{Significato} \\
    \hline
    \hline
    \constd{SOL\_SOCKET}& Opzioni generiche dei socket.\\
    \constd{SOL\_IP}    & Opzioni specifiche per i socket che usano IPv4.\\
    \constd{SOL\_TCP}   & Opzioni per i socket che usano TCP.\\
    \constd{SOL\_IPV6}  & Opzioni specifiche per i socket che usano IPv6.\\
    \constd{SOL\_ICMPV6}& Opzioni specifiche per i socket che usano ICMPv6.\\
    \hline
  \end{tabular}
  \caption{Possibili valori dell'argomento \param{level} delle 
    funzioni \func{setsockopt} e \func{getsockopt}.} 
  \label{tab:sock_option_levels}
\end{table}

I valori usati per \param{level}, corrispondenti ad un dato protocollo usato
da un socket, sono quelli corrispondenti al valore numerico che identifica il
suddetto protocollo in \conffile{/etc/protocols}; dato che la leggibilità di un
programma non trarrebbe certo beneficio dall'uso diretto dei valori numerici,
più comunemente si indica il protocollo tramite le apposite costanti
\texttt{SOL\_*} riportate in tab.~\ref{tab:sock_option_levels}, dove si sono
riassunti i valori che possono essere usati per l'argomento
\param{level}. 

La notazione in questo caso è, purtroppo, abbastanza confusa: infatti in Linux
il valore si può impostare sia usando le costanti \texttt{SOL\_*}, che le
analoghe \texttt{IPPROTO\_*} (citate anche da Stevens in \cite{UNP1});
entrambe hanno gli stessi valori che sono equivalenti ai numeri di protocollo
di \conffile{/etc/protocols}, con una eccezione specifica, che è quella del
protocollo ICMP, per la quale non esiste una costante, il che è comprensibile
dato che il suo valore, 1, è quello che viene assegnato a \const{SOL\_SOCKET}.

Il quarto argomento, \param{optval} è un puntatore ad una zona di memoria che
contiene i dati che specificano il valore dell'opzione che si vuole passare al
socket, mentre l'ultimo argomento \param{optlen},\footnote{questo argomento è
  in realtà sempre di tipo \ctyp{int}, come era nelle \acr{libc4} e
  \acr{libc5}; l'uso di \type{socklen\_t} è stato introdotto da POSIX (valgono
  le stesse considerazioni per l'uso di questo tipo di dato fatte in
  sez.~\ref{sec:TCP_func_accept}) ed adottato dalla \acr{glibc}.} è la
dimensione in byte dei dati presenti all'indirizzo indicato da \param{optval}.
Dato che il tipo di dati varia a seconda dell'opzione scelta, occorrerà
individuare qual è quello che deve essere usato, ed utilizzare le opportune
variabili.

La gran parte delle opzioni utilizzano per \param{optval} un valore intero, se
poi l'opzione esprime una condizione logica, il valore è sempre un intero, ma
si dovrà usare un valore non nullo per abilitarla ed un valore nullo per
disabilitarla.  Se invece l'opzione non prevede di dover ricevere nessun tipo
di valore si deve impostare \param{optval} a \val{NULL}. Un piccolo numero
di opzioni però usano dei tipi di dati peculiari, è questo il motivo per cui
\param{optval} è stato definito come puntatore generico.

La seconda funzione usata per controllare le proprietà dei socket è
\funcd{getsockopt}, che serve a leggere i valori delle opzioni dei socket ed a
farsi restituire i dati relativi al loro funzionamento; il suo prototipo è:


\begin{funcproto}{
\fhead{sys/socket.h}
\fhead{sys/types.h}
\fdecl{int getsockopt(int s, int level, int optname, void *optval,
    socklen\_t *optlen)}
\fdesc{Legge le opzioni di un socket.} 
}

{La funzione ritorna $0$ in caso di successo e $-1$ per un errore, nel qual
  caso \var{errno} assumerà uno dei valori: 
  \begin{errlist}
  \item[\errcode{EBADF}] il file descriptor \param{sock} non è valido.
  \item[\errcode{EFAULT}] l'indirizzo \param{optval} o quello di
    \param{optlen} non è valido.
  \item[\errcode{ENOPROTOOPT}] l'opzione scelta non esiste per il livello
    indicato. 
  \item[\errcode{ENOTSOCK}] il file descriptor \param{sock} non corrisponde ad
    un socket.
  \end{errlist}
}
\end{funcproto}

I primi tre argomenti sono identici ed hanno lo stesso significato di quelli
di \func{setsockopt}, anche se non è detto che tutte le opzioni siano definite
per entrambe le funzioni. In questo caso \param{optval} viene usato per
ricevere le informazioni ed indica l'indirizzo a cui andranno scritti i dati
letti dal socket, infine \param{optlen} diventa un puntatore ad una variabile
che viene usata come \textit{value result argument} per indicare, prima della
chiamata della funzione, la lunghezza del buffer allocato per \param{optval} e
per ricevere indietro, dopo la chiamata della funzione, la dimensione
effettiva dei dati scritti su di esso.  Se la dimensione del buffer allocato
per \param{optval} non è sufficiente si avrà un errore.



\subsection{Le opzioni generiche}
\label{sec:sock_generic_options}

Come accennato esiste un insieme generico di opzioni dei socket che possono
applicarsi a qualunque tipo di socket, indipendentemente da quale protocollo
venga poi utilizzato. Una descrizione di queste opzioni è generalmente
disponibile nella settima sezione delle pagine di manuale; nel caso specifico
la si può consultare con \texttt{man 7 socket}. Se si vuole operare su queste
opzioni generiche il livello da utilizzare è \const{SOL\_SOCKET}; si è
riportato un elenco di queste opzioni in tab.~\ref{tab:sock_opt_socklevel}.


\begin{table}[!htb]
  \centering
  \footnotesize
  \begin{tabular}[c]{|l|c|c|c|l|l|}
    \hline
    \textbf{Opzione}&\texttt{get}&\texttt{set}&\textbf{flag}&\textbf{Tipo}&
                    \textbf{Descrizione}\\
    \hline
    \hline
    \const{SO\_ACCEPTCONN}&$\bullet$&        &         &\texttt{int}& 
                          Indica se il socket è in ascolto.\\
    \const{SO\_ATTACH\_FILTER}&     &$\bullet$&         &\texttt{sock\_fprog}& 
                          Aggancia un filtro BPF al socket.\\
    \const{SO\_BINDTODEVICE}&$\bullet$&$\bullet$&         &\texttt{char *}& 
                          Lega il socket ad un dispositivo.\\
    \const{SO\_BROADCAST}&$\bullet$&$\bullet$&$\bullet$&\texttt{int}& 
                          Attiva o disattiva il \textit{broadcast}.\\ 
    \const{SO\_BSDCOMPAT}&$\bullet$&$\bullet$&$\bullet$&\texttt{int}& 
                          Abilita la compatibilità con BSD.\\
    \const{SO\_BUSY\_POLL}&$\bullet$&$\bullet$&$\bullet$&\texttt{int}& 
                          Attiva il ``\textit{busy poll}'' sul socket.\\
    \const{SO\_DEBUG}    &$\bullet$&$\bullet$&$\bullet$&\texttt{int}& 
                          Abilita il debugging sul socket.\\
    \const{SO\_DETACH\_FILTER}&    &$\bullet$&$\bullet$&\texttt{int}& 
                          Rimuove il filtro BPF agganciato al socket.\\
    \const{SO\_DOMAIN}   &$\bullet$&         &         &\texttt{int}& 
                          Legge il tipo di socket.\\
    \const{SO\_DONTROUTE}&$\bullet$&$\bullet$&$\bullet$&\texttt{int}& 
                          Non invia attraverso un gateway.\\
    \const{SO\_ERROR}    &$\bullet$&         &         &\texttt{int}& 
                          Riceve e cancella gli errori pendenti.\\
    \const{SO\_KEEPALIVE}&$\bullet$&$\bullet$&$\bullet$&\texttt{int}& 
                          Controlla l'attività della connessione.\\
    \const{SO\_LINGER}   &$\bullet$&$\bullet$&         &\texttt{linger}&
                          Indugia nella chiusura con dati da spedire.\\
    \const{SO\_LOCK\_FILTER}&    &$\bullet$&$\bullet$&\texttt{int}& 
                          Blocca il filtro BPF agganciato al socket.\\
    \const{SO\_MARK}     &$\bullet$&$\bullet$&         &\texttt{int}&
                          Imposta un ``\textit{firewall mark}'' sul socket.\\
    \const{SO\_OOBINLINE}&$\bullet$&$\bullet$&$\bullet$&\texttt{int}& 
                          Lascia in linea i dati \textit{out-of-band}.\\
    \const{SO\_PASSCRED} &$\bullet$&$\bullet$&$\bullet$&\texttt{int}& 
                          Abilita la ricezione di credenziali.\\
    \const{SO\_PEEK\_OFF} &$\bullet$&$\bullet$&$\bullet$&\texttt{int}& 
                          Imposta il valore del ``\textit{peek offset}''.\\
    \const{SO\_PEERCRED} &$\bullet$&         &         &\texttt{ucred}& 
                          Restituisce le credenziali del processo remoto.\\
    \const{SO\_PRIORITY} &$\bullet$&$\bullet$&         &\texttt{int}& 
                          Imposta la priorità del socket.\\
    \const{SO\_PROTOCOL} &$\bullet$&         &         &\texttt{int}& 
                          Ottiene il protocollo usato dal socket.\\
    \const{SO\_RCVBUF}   &$\bullet$&$\bullet$&         &\texttt{int}& 
                          Imposta dimensione del buffer di ricezione.\\
    \const{SO\_RCVBUFFORCE}&$\bullet$&$\bullet$&         &\texttt{int}& 
                          Forza dimensione del buffer di ricezione.\\
    \const{SO\_RCVLOWAT} &$\bullet$&$\bullet$&$\bullet$&\texttt{int}& 
                          Basso livello sul buffer di ricezione.\\
    \const{SO\_RCVTIMEO} &$\bullet$&$\bullet$&         &\texttt{timeval}& 
                          Timeout in ricezione.\\
    \const{SO\_REUSEADDR}&$\bullet$&$\bullet$&$\bullet$&\texttt{int}& 
                          Consente il riutilizzo di un indirizzo locale.\\
    \const{SO\_REUSEPORT}&$\bullet$&$\bullet$&$\bullet$&\texttt{int}& 
                          Consente il riutilizzo di una porta.\\
    \const{SO\_RXQ\_OVFL} &$\bullet$&$\bullet$&$\bullet$&\texttt{int}& 
                          Richiede messaggio ancillare con pacchetti persi.\\
    \const{SO\_SNDBUF}   &$\bullet$&$\bullet$&         &\texttt{int}& 
                          Imposta dimensione del buffer di trasmissione.\\
    \const{SO\_SNDBUFFORCE}&$\bullet$&$\bullet$&         &\texttt{int}& 
                          Forza dimensione del buffer di trasmissione.\\
    \const{SO\_SNDLOWAT} &$\bullet$&$\bullet$&         &\texttt{int}&
                          Basso livello sul buffer di trasmissione.\\
    \const{SO\_SNDTIMEO} &$\bullet$&$\bullet$&         &\texttt{timeval}& 
                          Timeout in trasmissione.\\
    \const{SO\_TIMESTAMP}&$\bullet$&$\bullet$&$\bullet$&\texttt{int}&
                          Abilita/disabilita la ricezione dei \textit{timestamp}.\\
    \const{SO\_TYPE}     &$\bullet$&         &         &\texttt{int}& 
                          Restituisce il tipo di socket.\\
   \hline
  \end{tabular}
  \caption{Le opzioni disponibili al livello \const{SOL\_SOCKET}.} 
  \label{tab:sock_opt_socklevel}
\end{table}

% TODO aggiungere e documentare SO_ATTACH_BPF, introdotta con il kernel 3.19,
% vedi http://lwn.net/Articles/625224/
% TODO aggiungere e documentare SO_INCOMING_CPU, introdotta con il kernel 3.19,
% TODO documentare SO_PEERGROUPS introdotta con il kernel 4.13, citata
% in https://lwn.net/Articles/727385/
% TODO documentare SO_ZEROCOPY, vedi https://lwn.net/Articles/726917/ (e il
% resto pure) introdotta con il kernel 4.14

La tabella elenca le costanti che identificano le singole opzioni da usare
come valore per \param{optname}; le due colonne seguenti indicano per quali
delle due funzioni (\func{getsockopt} o \func{setsockopt}) l'opzione è
disponibile, mentre la colonna successiva indica, quando si ha a che fare con
un valore di \param{optval} intero, se l'opzione è da considerare un numero o
un valore logico. Si è inoltre riportato sulla quinta colonna il tipo di dato
usato per \param{optval} ed una breve descrizione del significato delle
singole opzioni sulla sesta.

Le descrizioni delle opzioni presenti in tab.~\ref{tab:sock_opt_socklevel}
sono estremamente sommarie, è perciò necessario fornire un po' più di
informazioni. Alcune opzioni inoltre hanno una notevole rilevanza nella
gestione dei socket, e pertanto il loro utilizzo sarà approfondito
separatamente in sez.~\ref{sec:sock_options_main}. Quello che segue è quindi
soltanto un elenco più dettagliato della breve descrizione di
tab.~\ref{tab:sock_opt_socklevel} sul significato delle varie opzioni:

\begin{basedescript}{\desclabelwidth{2.2cm}\desclabelstyle{\nextlinelabel}}
\item[\constd{SO\_ACCEPTCONN}] questa opzione permette di rilevare se il socket
  su cui opera è stato posto in modalità di ricezione di eventuali connessioni
  con una chiamata a \func{listen}. L'opzione può essere usata soltanto con
  \func{getsockopt} ed utilizza per \param{optval} un intero in cui viene
  restituito 1 se il socket è in ascolto e 0 altrimenti.

\item[\constd{SO\_ATTACH\_FILTER}] questa opzione permette di agganciare ad un
  socket un filtro di selezione dei pacchetti con la stessa sintassi del BPF
  (\textit{Berkley Packet Filter}) di BSD, che consente di selezionare, fra
  tutti quelli ricevuti, verranno letti. Può essere usata solo con
  \func{setsockopt} ed utilizza per \param{optval} un puntatore ad una
  struttura \struct{sock\_fprog} (definita in
  \headfile{linux/filter.h}). Questa opzione viene usata principalmente con i
  socket di tipo \const{AF\_PACKET} (torneremo su questo in
  sez.~\ref{sec:packet_socket}) dalla libreria \texttt{libpcap} per
  implementare programmi di cattura dei pacchetti, e per questo tipo di
  applicazione è opportuno usare sempre quest'ultima.\footnote{la trattazione
    del BPF va al di là dell'argomento di questa sezione per la documentazione
    si consulti il file \texttt{networking/filter.txt} nella documentazione
    del kernel.}

\item[\constd{SO\_BINDTODEVICE}] questa opzione permette di \textsl{legare} il
  socket ad una particolare interfaccia, in modo che esso possa ricevere ed
  inviare pacchetti solo su quella. L'opzione richiede per \param{optval} il
  puntatore ad una stringa contenente il nome dell'interfaccia (ad esempio
  \texttt{eth0}); utilizzando una stringa nulla o un valore nullo per
  \param{optlen} si può rimuovere un precedente collegamento.

  Il nome della interfaccia deve essere specificato con una stringa terminata
  da uno zero e di lunghezza massima pari a \constd{IFNAMSIZ}; l'opzione è
  effettiva solo per alcuni tipi di socket, ed in particolare per quelli della
  famiglia \const{AF\_INET}; non è invece supportata per i \textit{packet
    socket} (vedi sez.~\ref{sec:socket_raw}). 

\item[\constd{SO\_BROADCAST}] questa opzione abilita il \textit{broadcast}:
  quando abilitata i socket di tipo \const{SOCK\_DGRAM} riceveranno i
  pacchetti inviati all'indirizzo di \textit{broadcast} e potranno scrivere
  pacchetti su tale indirizzo.  Prende per \param{optval} un intero usato come
  valore logico. L'opzione non ha effetti su un socket di tipo
  \const{SOCK\_STREAM}.

\item[\constd{SO\_BSDCOMPAT}] questa opzione abilita la compatibilità con il
  comportamento di BSD (in particolare ne riproduce alcuni bug).  Attualmente
  è una opzione usata solo per il protocollo UDP e ne è prevista la rimozione.
  L'opzione utilizza per \param{optval} un intero usato come valore logico.

  Quando viene abilitata gli errori riportati da messaggi ICMP per un socket
  UDP non vengono passati al programma in \textit{user space}. Con le versioni
  2.0.x del kernel erano anche abilitate altre opzioni di compatibilità per i
  socket raw (modifiche casuali agli header, perdita del flag di
  \textit{broadcast}) che sono state rimosse con il passaggio al 2.2; è
  consigliato correggere i programmi piuttosto che usare questa funzione. Dal
  kernel 2.4 viene ignorata, e dal 2.6 genera un messaggio di log del kernel.

\item[\constd{SO\_BUSY\_POLL}] questa opzione, presente dal kernel 3.11,
  imposta un tempo approssimato in microsecondi, per cui in caso di ricezione
  bloccante verrà eseguito un \itindex{busy~poll} ``\textit{busy
    poll}'',\footnote{con \textit{busy poll} si fa riferimento al
    \textit{polling} su una risorsa occupata; si continuerà cioè a tentare di
    leggere anche quando non ci sono dati senza portare il processo stato di
    \textit{sleep}, in alcuni casi, quando ci si aspetta che i dati arrivino a
    breve, questa tecnica può dare un miglioramento delle prestazioni.}  da
  indicare in \param{optval} con un valore intero. Si tratta di una opzione
  utilizzabile solo con socket che ricevono dati da un dispositivo di rete che
  la supporti, e che consente di ridurre la latenza per alcune applicazioni,
  ma che comporta un maggiore utilizzo della CPU (e quindi di energia); per
  questo il valore può essere aumentato solo da processi con i privilegi di
  amministratore (in particolare con la \textit{capability}
  \const{CAP\_NET\_ADMIN}).  Il valore di default viene controllato dal file
  \sysctlrelfile{net/core}{busy\_read} per le funzioni di lettura mentre il
  file \sysctlrelfile{net/core}{busy\_poll} controlla il ``\textit{busy
    poll}'' per \func{select} e \func{poll}.

\item[\constd{SO\_DEBUG}] questa opzione abilita il debugging delle operazioni
  dei socket; l'opzione utilizza per \param{optval} un intero usato come
  valore logico, e può essere utilizzata solo da un processo con i privilegi
  di amministratore (in particolare con la \textit{capability}
  \const{CAP\_NET\_ADMIN}).  L'opzione necessita inoltre dell'opportuno
  supporto nel kernel;\footnote{deve cioè essere definita la macro di
    preprocessore \macrod{SOCK\_DEBUGGING} nel file \file{include/net/sock.h}
    dei sorgenti del kernel, questo è sempre vero nei kernel delle serie
    superiori alla 2.3, per i kernel delle serie precedenti invece è
    necessario aggiungere a mano detta definizione; è inoltre possibile
    abilitare anche il tracciamento degli stati del TCP definendo la macro
    \macrod{STATE\_TRACE} in \file{include/net/tcp.h}.}  quando viene
  abilitata una serie di messaggi con le informazioni di debug vengono inviati
  direttamente al sistema del kernel log.\footnote{si tenga presente che il
    comportamento è diverso da quanto avviene con BSD, dove l'opzione opera
    solo sui socket TCP, causando la scrittura di tutti i pacchetti inviati
    sulla rete su un buffer circolare che viene letto da un apposito
    programma, \cmd{trpt}.}

\item[\constd{SO\_DETACH\_FILTER}] consente di distaccare un filtro
  precedentemente aggiunto ad un socket con l'opzione
  \const{SO\_ATTACH\_FILTER}, in genere non viene usata direttamente in quanto
  i filtri BPF vengono automaticamente rimossi alla chiusura del socket, il
  suo utilizzo è pertanto limitato ai rari casi in cui si vuole rimuovere un
  precedente filtro per inserirne uno diverso. Come \const{SO\_ATTACH\_FILTER}
  può essere usato solo \func{setsockopt} e prende per \param{optval} un
  intero usato come valore logico.

\item[\constd{SO\_DOMAIN}] questa opzione, presente dal kernel 2.6.32, legge
  il ``\textsl{dominio}'' (la famiglia di indirizzi) del socket.
. Funziona solo con \func{getsockopt}, ed
  utilizza per \param{optval} un intero in cui verrà restituito il valore
  numerico che lo identifica (ad esempio \texttt{AF\_INET}).

\item[\constd{SO\_DONTROUTE}] questa opzione richiede che l'invio dei
  pacchetti del socket sia eseguito soltanto verso una destinazione
  direttamente connessa, impedendo l'uso di un \textit{gateway} e saltando
  ogni processo relativo all'uso della tabella di routing del kernel. Prende
  per \param{optval} un intero usato come valore logico. È equivalente all'uso
  del flag \const{MSG\_DONTROUTE} su una \func{send} (vedi
  sez.~\ref{sec:net_send_recv}).

\item[\constd{SO\_ERROR}] questa opzione riceve un errore presente sul socket;
  può essere utilizzata soltanto con \func{getsockopt} e prende per
  \param{optval} un valore intero, nel quale viene restituito il codice di
  errore, e la condizione di errore sul socket viene cancellata. Viene
  usualmente utilizzata per ricevere il codice di errore, come accennato in
  sez.~\ref{sec:TCP_sock_select}, quando si sta osservando il socket con una
  \func{select} che ritorna a causa dello stesso.

\item[\const{SO\_KEEPALIVE}] questa opzione abilita un meccanismo di verifica
  della persistenza di una connessione associata al socket; è pertanto
  effettiva solo sui socket che supportano le connessioni, ed è usata
  principalmente con il TCP. L'opzione utilizza per \param{optval} un intero
  usato come valore logico. Maggiori dettagli sul suo funzionamento sono
  forniti in sez.~\ref{sec:sock_options_main}.

\item[\const{SO\_LINGER}] questa opzione controlla le modalità con cui viene
  chiuso un socket quando si utilizza un protocollo che supporta le
  connessioni (è pertanto usata con i socket TCP ed ignorata per UDP) e
  modifica il comportamento delle funzioni \func{close} e \func{shutdown}.
  L'opzione richiede che l'argomento \param{optval} sia una struttura di tipo
  \struct{linger}, definita in \headfile{sys/socket.h} ed illustrata in
  fig.~\ref{fig:sock_linger_struct}.  Maggiori dettagli sul suo funzionamento
  sono forniti in sez.~\ref{sec:sock_options_main}.

\item[\constd{SO\_LOCK\_FILTER}] consente di bloccare un filtro
  precedentemente aggiunto ad un socket con l'opzione
  \const{SO\_ATTACH\_FILTER}, in modo che non possa essere né rimosso né
  modificato, questo consente di impostare un filtro su un socket, bloccarlo e
  poi cedere i privilegi con la sicurezza che il filtro permarrà fino alla
  chiusura del socket. Come \const{SO\_ATTACH\_FILTER} può essere usato solo
  \func{setsockopt} e prende per \param{optval} un intero usato come valore
  logico.

\item[\constd{SO\_MARK}] questa opzione, presente dal kernel 2.6.25, imposta
  un valore di marcatura sui pacchetti del socket. Questa è una funzionalità
  specifica di Linux, ottenibile anche con l'uso del target \texttt{MARK}
  del comando \texttt{iptables} (l'argomento è trattato in sez.~3.3.5 di
  \cite{SGL}). Il valore di marcatura viene mantenuto all'interno dello stack
  di rete del kernel e può essere usato sia dal \itindex{netfilter}
  \textit{netfilter}\footnote{il \textit{netfilter} è l'infrastruttura usata
    per il filtraggio dei pacchetti del kernel, per maggiori dettagli si
    consulti il cap.~2 di \cite{FwGL}.} di Linux che per impostare politiche
  di routing avanzato. Il valore deve essere specificato in \param{optval}
  come un intero. L'opzione richiede i privilegi di amministratore con la
  \textit{capability} \const{CAP\_NET\_ADMIN}.

\item[\constd{SO\_OOBINLINE}] se questa opzione viene abilitata i dati
  \textit{out-of-band} vengono inviati direttamente nel flusso di dati del
  socket (e sono quindi letti con una normale \func{read}) invece che restare
  disponibili solo per l'accesso con l'uso del flag \const{MSG\_OOB} di
  \func{recvmsg}. L'argomento è trattato in dettaglio in
  sez.~\ref{sec:TCP_urgent_data}. L'opzione funziona soltanto con socket che
  supportino i dati \textit{out-of-band} (non ha senso per socket UDP ad
  esempio), ed utilizza per \param{optval} un intero usato come valore logico.

\item[\constd{SO\_PASSCRED}] questa opzione abilita sui socket unix-domain
  (vedi sez.~\ref{sec:unix_socket}) la ricezione dei messaggi di controllo di
  tipo \const{SCM\_CREDENTIALS}. Prende come \param{optval} un intero usato
  come valore logico.

\item[\constd{SO\_PEEK\_OFF}] questa opzione, disponibile a partire dal kernel
  3.4, imposta un ``\textit{peek offset}'' sul socket (attualmente disponibile
  solo per i socket unix-domain (vedi sez.~\ref{sec:unix_socket}). La funzione
  serve come ausilio per l'uso del flag \const{MSG\_PEEK} di \func{recv} che
  consente di ``\textsl{sbirciare}'' nei dati di un socket, cioè di leggerli
  senza rimuoverli dalla coda in cui sono mantenuti, così che vengano
  restituiti anche da una successiva lettura ordinaria.
 
  Un valore negativo (il default è $-1$) riporta alla situazione ordinaria in
  cui si ``\textsl{sbircia}'' a partire dalla testa della coda, un valore
  positivo consente di leggere a partire dalla posizione indicata nella coda e
  tutte le volte che si sbirciano dei dati il valore dell'offset viene
  automaticamente aumentato della quantità di dati sbirciati, in modo che si
  possa proseguire da dove si era arrivati. Il valore deve essere specificato
  in \param{optval} come intero.

\item[\constd{SO\_PEERCRED}] questa opzione restituisce le credenziali del
  processo remoto connesso al socket; l'opzione è disponibile solo per socket
  unix-domain di tipo \textit{stream} e anche per quelli di tipo
  \textit{datagram} quando se ne crea una coppia con \func{socketpair}, e può
  essere usata solo con \func{getsockopt}.  Utilizza per
  \param{optval} una apposita struttura \struct{ucred} (vedi
  sez.~\ref{sec:unix_socket}). 

\item[\constd{SO\_PRIORITY}] questa opzione permette di impostare le priorità
  per tutti i pacchetti che sono inviati sul socket, prende per \param{optval}
  un valore intero. Con questa opzione il kernel usa il valore per ordinare le
  priorità sulle code di rete,\footnote{questo richiede che sia abilitato il
    sistema di \textit{Quality of Service} disponibile con le opzioni di
    routing avanzato.} i pacchetti con priorità più alta vengono processati
  per primi, in modalità che dipendono dalla disciplina di gestione della
  coda. Nel caso di protocollo IP questa opzione permette anche di impostare i
  valori del campo \textit{type of service} (noto come TOS, vedi
  sez.~\ref{sec:IP_header}) per i pacchetti uscenti. Per impostare una
  priorità al di fuori dell'intervallo di valori fra 0 e 6 sono richiesti i
  privilegi di amministratore con la \textit{capability}
  \const{CAP\_NET\_ADMIN}.

\item[\constd{SO\_PROTOCOL}] questa opzione, presente dal kernel 2.6.32, legge
  il protocollo usato dal socket. Funziona solo con \func{getsockopt}, ed
  utilizza per \param{optval} un intero in cui verrà restituito il valore
  numerico che lo identifica (ad esempio \texttt{IPPROTO\_TCP}).

\item[\constd{SO\_RCVBUF}] questa opzione imposta la dimensione del buffer di
  ricezione del socket. Prende per \param{optval} un intero indicante il
  numero di byte. Il valore di default ed il valore massimo che si può
  specificare come argomento per questa opzione sono impostabili tramiti gli
  opportuni valori di \func{sysctl} (vedi sez.~\ref{sec:sock_sysctl}).

  Si tenga presente che nel caso di socket TCP il kernel alloca effettivamente
  una quantità di memoria doppia rispetto a quanto richiesto con
  \func{setsockopt} per entrambe le opzioni \const{SO\_RCVBUF} e
  \const{SO\_SNDBUF}. Questo comporta che una successiva lettura con
  \func{getsockopt} riporterà un valore diverso da quello impostato con
  \func{setsockopt}.  Questo avviene perché TCP necessita dello spazio in più
  per mantenere dati amministrativi e strutture interne, e solo una parte
  viene usata come buffer per i dati, mentre il valore letto da
  \func{getsockopt} e quello riportato nei vari parametri di
  \textit{sysctl}\footnote{cioè \sysctlrelfile{net/core}{wmem\_max} e
    \sysctlrelfile{net/core}{rmem\_max} in \texttt{/proc/sys/net/core} e
    \sysctlrelfile{net/ipv4}{tcp\_wmem} e \sysctlrelfile{net/ipv4}{tcp\_rmem}
    in \texttt{/proc/sys/net/ipv4}, vedi sez.~\ref{sec:sock_sysctl}.} indica
  la memoria effettivamente impiegata.  Si tenga presente inoltre che le
  modifiche alle dimensioni dei buffer di ricezione e trasmissione, per poter
  essere effettive, devono essere impostate prima della chiamata alle funzioni
  \func{listen} o \func{connect}.

\item[\constd{SO\_RCVBUFFORCE}] questa opzione, presente dal kernel 2.6.14, è
  identica a \const{SO\_RCVBUF} ma consente ad un processo con i privilegi di
  amministratore (per la precisione con la \textit{capability}
  \const{CAP\_NET\_ADMIN}) di impostare in valore maggiore del limite di
  \sysctlrelfile{net/core}{rmem\_max}.

\item[\constd{SO\_RCVLOWAT}] questa opzione imposta il valore che indica il
  numero minimo di byte che devono essere presenti nel buffer di ricezione
  perché il kernel passi i dati all'utente, restituendoli ad una \func{read} o
  segnalando ad una \func{select} (vedi sez.~\ref{sec:TCP_sock_select}) che ci
  sono dati in ingresso. L'opzione utilizza per \param{optval} un intero che
  specifica il numero di byte; con Linux questo valore è sempre 1 e può essere
  cambiato solo con i kernel a partire dal 2.4. Si tenga presente però che per
  i kernel prima del 2.6.28 sia \func{poll} che \func{select} non supportano
  questa funzionalità e ritornano comunque, indicando il socket come
  leggibile, non appena almeno un byte è presente, con una successiva lettura
  con \func{read} che si blocca fintanto che non diventa disponibile la
  quantità di byte indicati. Con \func{getsockopt} si può leggere questo
  valore mentre \func{setsockopt} darà un errore di \errcode{ENOPROTOOPT}
  quando il cambiamento non è supportato.

\item[\constd{SO\_RCVTIMEO}] l'opzione permette di impostare un tempo massimo
  sulle operazioni di lettura da un socket, e prende per \param{optval} una
  struttura di tipo \struct{timeval} (vedi fig.~\ref{fig:sys_timeval_struct})
  identica a quella usata con \func{select}. Con \func{getsockopt} si può
  leggere il valore attuale, mentre con \func{setsockopt} si imposta il tempo
  voluto, usando un valore nullo per \struct{timeval} il timeout viene
  rimosso. 

  Se l'opzione viene attivata tutte le volte che una delle funzioni di lettura
  (\func{read}, \func{readv}, \func{recv}, \func{recvfrom} e \func{recvmsg})
  si blocca in attesa di dati per un tempo maggiore di quello impostato, essa
  ritornerà un valore -1 e la variabile \var{errno} sarà impostata con un
  errore di \errcode{EAGAIN} e \errcode{EWOULDBLOCK}, così come sarebbe
  avvenuto se si fosse aperto il socket in modalità non bloccante.\footnote{in
    teoria, se il numero di byte presenti nel buffer di ricezione fosse
    inferiore a quello specificato da \const{SO\_RCVLOWAT}, l'effetto potrebbe
    essere semplicemente quello di provocare l'uscita delle funzioni di
    lettura restituendo il numero di byte fino ad allora ricevuti.}

  In genere questa opzione non è molto utilizzata se si ha a che fare con la
  lettura dei dati, in quanto è sempre possibile usare una \func{select} che
  consente di specificare un \textit{timeout}; l'uso di \func{select} non
  consente però di impostare il timeout per l'uso di \func{connect}, per avere
  il quale si può ricorrere a questa opzione (nel qual caso il raggiungimento
  del \textit{timeout} restituisce l'errore \errcode{EINPROGRESS}). 

\item[\const{SO\_REUSEADDR}] questa opzione permette di eseguire la funzione
  \func{bind} su indirizzi locali che siano già in uso da altri socket;
  l'opzione utilizza per \param{optval} un intero usato come valore logico.
  Questa opzione modifica il comportamento normale dell'interfaccia dei socket
  che fa fallire l'esecuzione della funzione \func{bind} con un errore di
  \errcode{EADDRINUSE} quando l'indirizzo locale (più propriamente il
  controllo viene eseguito sulla porta) è già in uso da parte di un altro
  socket.  Maggiori dettagli sul suo funzionamento sono forniti in
  sez.~\ref{sec:sock_options_main}.

\item[\const{SO\_REUSEPORT}] questa opzione, presente a partire dal kernel
  3.9, permette di far usare a più socket di tipo \const{AF\_INET} o
  \const{AF\_INET6} lo stesso indirizzo locale, e costituisce una estensione
  della precedente \const{SO\_REUSEADDR}.  Maggiori dettagli sul suo
  funzionamento sono forniti in sez.~\ref{sec:sock_options_main}.

\item[\constd{SO\_RXQ\_OVFL}] questa opzione, presente dal kernel 2.6.33,
  permette di abilitare o disabilitare sul socket la ricezione di un messaggio
  ancillare (tratteremo l'argomento in sez.~\ref{sec:net_ancillary_data})
  contenente un intero senza segno a 32 bit che indica il numero di pacchetti
  scartati sul socket fra l'ultimo pacchetto ricevuto e questo.  L'opzione
  utilizza per \param{optval} un intero usato come valore logico.

% TODO documentare SO_RXQ_OVFL, vedi https://lwn.net/Articles/626150/ capire
% che significa 'questo'

\item[\constd{SO\_SNDLOWAT}] questa opzione imposta il valore che indica il
  numero minimo di byte che devono essere presenti nel buffer di trasmissione
  perché il kernel li invii al protocollo successivo, consentendo ad una
  \func{write} di ritornare o segnalando ad una \func{select} (vedi
  sez.~\ref{sec:TCP_sock_select}) che è possibile eseguire una scrittura.
  L'opzione utilizza per \param{optval} un intero che specifica il numero di
  byte; con Linux questo valore è sempre 1 e non può essere cambiato;
  \func{getsockopt} leggerà questo valore mentre \func{setsockopt} darà un
  errore di \errcode{ENOPROTOOPT}.

\item[\constd{SO\_SNDBUF}] questa opzione imposta la dimensione del buffer di
  trasmissione del socket. Prende per \param{optval} un intero indicante il
  numero di byte. Il valore di default ed il valore massimo che si possono
  specificare come argomento per questa opzione sono impostabili
  rispettivamente tramite gli opportuni valori di \func{sysctl} (vedi
  sez.~\ref{sec:sock_sysctl}).

\item[\constd{SO\_SNDBUFFORCE}] questa opzione, presente dal kernel 2.6.14, è
  identica a \const{SO\_SNDBUF} ma consente ad un processo con i privilegi di
  amministratore (per la precisione con la \textit{capability}
  \const{CAP\_NET\_ADMIN}) di impostare in valore maggiore del limite di
  \sysctlrelfile{net/core}{wmem\_max}.

% TODO verificare il timeout con un programma di test

\item[\constd{SO\_SNDTIMEO}] l'opzione permette di impostare un tempo massimo
  sulle operazioni di scrittura su un socket, ed usa gli stessi valori di
  \const{SO\_RCVTIMEO}.  In questo caso però si avrà un errore di
  \errcode{EAGAIN} o \errcode{EWOULDBLOCK} per le funzioni di scrittura
  \func{write}, \func{writev}, \func{send}, \func{sendto} e \func{sendmsg}
  qualora queste restino bloccate per un tempo maggiore di quello specificato. 

\item[\constd{SO\_TIMESTAMP}] questa opzione, presente dal kernel 2.6.30,
  permette di abilitare o disabilitare sul socket la ricezione di un messaggio
  ancillare di tipo \const{SO\_TIMESTAMP}. Il messaggio viene inviato con
  livello \const{SOL\_SOCKET} ed nel campo \var{cmsg\_data} (per i dettagli si
  veda sez.~\ref{sec:net_ancillary_data}) viene impostata una struttura
  \struct{timeval} che indica il tempo di ricezione dell'ultimo pacchetto
  passato all'utente.  L'opzione utilizza per \param{optval} un intero usato
  come valore logico.

% TODO capire meglio la tempistica di questi messaggi ancillari
% TODO documentare SO_TIMESTAMP e le altre opzioni di timestamping dei 
% pacchetti, introdotte nel 2.6.30, vedi nei sorgenti del kernel:
% Documentation/networking/timestamping.txt

\item[\constd{SO\_TYPE}] questa opzione permette di leggere il tipo di socket
  su cui si opera; funziona solo con \func{getsockopt}, ed utilizza per
  \param{optval} un intero in cui verrà restituito il valore numerico che lo
  identifica (ad esempio \const{SOCK\_STREAM}). 


\end{basedescript}


\subsection{L'uso delle principali opzioni dei socket}
\label{sec:sock_options_main}

La descrizione sintetica del significato delle opzioni generiche dei socket,
riportata nell'elenco in sez.~\ref{sec:sock_generic_options}, è
necessariamente sintetica, alcune di queste però possono essere utilizzate
per controllare delle funzionalità che hanno una notevole rilevanza nella
programmazione dei socket.  Per questo motivo faremo in questa sezione un
approfondimento sul significato delle opzioni generiche più importanti.


\constbeg{SO\_KEEPALIVE}
\subsubsection{L'opzione \const{SO\_KEEPALIVE}}

La prima opzione da approfondire è \const{SO\_KEEPALIVE} che permette di
tenere sotto controllo lo stato di una connessione. Una connessione infatti
resta attiva anche quando non viene effettuato alcun traffico su di essa; è
allora possibile, in caso di una interruzione completa della rete, che la
caduta della connessione non venga rilevata, dato che sulla stessa non passa
comunque alcun traffico.

Se si imposta questa opzione, è invece cura del kernel inviare degli appositi
messaggi sulla rete, detti appunto \textit{keep-alive}, per verificare se la
connessione è attiva.  L'opzione funziona soltanto con i socket che supportano
le connessioni (non ha senso per socket UDP ad esempio) e si applica
principalmente ai socket TCP.

Con le impostazioni di default (che sono riprese da BSD) Linux emette un
messaggio di \textit{keep-alive} (in sostanza un segmento ACK vuoto, cui sarà
risposto con un altro segmento ACK vuoto) verso l'altro capo della connessione
se questa è rimasta senza traffico per più di due ore.  Se è tutto a posto il
messaggio viene ricevuto e verrà emesso un segmento ACK di risposta, alla cui
ricezione ripartirà un altro ciclo di attesa per altre due ore di inattività;
il tutto avviene all'interno del kernel e le applicazioni non riceveranno
nessun dato.

Qualora ci siano dei problemi di rete si possono invece verificare i due casi
di terminazione precoce del server già illustrati in
sez.~\ref{sec:TCP_conn_crash}. Il primo è quello in cui la macchina remota ha
avuto un crollo del sistema ed è stata riavviata, per cui dopo il riavvio la
connessione non esiste più.\footnote{si ricordi che un normale riavvio o il
  crollo dell'applicazione non ha questo effetto, in quanto in tal caso si
  passa sempre per la chiusura del processo, e questo, come illustrato in
  sez.~\ref{sec:file_open_close}, comporta anche la regolare chiusura del
  socket con l'invio di un segmento FIN all'altro capo della connessione.} In
questo caso all'invio del messaggio di \textit{keep-alive} si otterrà come
risposta un segmento RST che indica che l'altro capo non riconosce più
l'esistenza della connessione ed il socket verrà chiuso riportando un errore
di \errcode{ECONNRESET}.

Se invece non viene ricevuta nessuna risposta (indice che la macchina non è
più raggiungibile) l'emissione dei messaggi viene ripetuta ad intervalli di 75
secondi per un massimo di 9 volte\footnote{entrambi questi valori possono
  essere modificati a livello di sistema (cioè per tutti i socket) con gli
  opportuni parametri illustrati in sez.~\ref{sec:sock_sysctl} ed a livello di
  singolo socket con le opzioni \texttt{TCP\_KEEP*} di
  sez.~\ref{sec:sock_tcp_udp_options}.}  (per un totale di 11 minuti e 15
secondi) dopo di che, se non si è ricevuta nessuna risposta, il socket viene
chiuso dopo aver impostato un errore di \errcode{ETIMEDOUT}. Qualora la
connessione si sia ristabilita e si riceva un successivo messaggio di risposta
il ciclo riparte come se niente fosse avvenuto.  Infine se si riceve come
risposta un pacchetto ICMP di destinazione irraggiungibile (vedi
sez.~\ref{sec:ICMP_protocol}), verrà restituito l'errore corrispondente.

In generale questa opzione serve per individuare una caduta della connessione
anche quando non si sta facendo traffico su di essa.  Viene usata
principalmente sui server per evitare di mantenere impegnate le risorse che
verrebbero dedicate a trattare delle connessioni che in realtà sono già
terminate (quelle che vengono anche chiamate connessioni
\textsl{semi-aperte}); in tutti quei casi cioè in cui il server si trova in
attesa di dati in ingresso su una connessione che non arriveranno mai o perché
il client sull'altro capo non è più attivo o perché non è più in grado di
comunicare con il server via rete.

\begin{figure}[!htbp]
  \footnotesize \centering
  \begin{minipage}[c]{\codesamplewidth}
    \includecodesample{listati/TCP_echod_fourth.c}
  \end{minipage}
  \normalsize
  \caption{La sezione della nuova versione del server del servizio
    \textit{echo} che prevede l'attivazione del \textit{keepalive} sui
    socket.}
  \label{fig:echod_keepalive_code}
\end{figure}

Abilitandola dopo un certo tempo le connessioni effettivamente terminate
verranno comunque chiuse per cui, utilizzando ad esempio una \func{select}, se
ne potrà rilevare la conclusione e ricevere il relativo errore. Si tenga
presente però che non può avere la certezza assoluta che un errore di
\errcode{ETIMEDOUT} ottenuto dopo aver abilitato questa opzione corrisponda
necessariamente ad una reale conclusione della connessione, il problema
potrebbe anche essere dovuto ad un problema di routing che perduri per un
tempo maggiore di quello impiegato nei vari tentativi di ritrasmissione del
\textit{keep-alive} (anche se questa non è una condizione molto probabile).

Come esempio dell'utilizzo di questa opzione introduciamo all'interno del
nostro server per il servizio \textit{echo} la nuova opzione \texttt{-k} che
permette di attivare il \textit{keep-alive} sui socket; tralasciando la parte
relativa alla gestione di detta opzione (che si limita ad assegnare ad 1 la
variabile \var{keepalive}) tutte le modifiche al server sono riportate in
fig.~\ref{fig:echod_keepalive_code}. Al solito il codice completo è contenuto
nel file \texttt{TCP\_echod\_fourth.c} dei sorgenti allegati alla guida.

Come si può notare la variabile \var{keepalive} è preimpostata (\texttt{\small
  8}) ad un valore nullo; essa viene utilizzata sia come variabile logica per
la condizione (\texttt{\small 14}) che controlla l'attivazione del
\textit{keep-alive} che come valore dell'argomento \param{optval} della
chiamata a \func{setsockopt} (\texttt{\small 16}).  A seconda del suo valore
tutte le volte che un processo figlio viene eseguito in risposta ad una
connessione verrà pertanto eseguita o meno la sezione (\texttt{\small 14--17})
che esegue l'impostazione di \const{SO\_KEEPALIVE} sul socket connesso,
attivando il relativo comportamento.
\constend{SO\_KEEPALIVE}



\constbeg{SO\_REUSEADDR}
\subsubsection{Le opzioni \const{SO\_REUSEADDR} e \const{SO\_REUSEPORT}}

La seconda opzione da approfondire è \constd{SO\_REUSEADDR}, che consente di
eseguire \func{bind} su un socket anche quando la porta specificata è già in
uso da parte di un altro socket. Si ricordi infatti che, come accennato in
sez.~\ref{sec:TCP_func_bind}, normalmente la funzione \func{bind} fallisce con
un errore di \errcode{EADDRINUSE} se la porta scelta è già utilizzata da un
altro socket, proprio per evitare che possano essere lanciati due server sullo
stesso indirizzo e la stessa porta, che verrebbero a contendersi i pacchetti
aventi quella destinazione.

Esistono però situazioni ed esigenze particolari in cui non si vuole che
questo comportamento di salvaguardia accada, ed allora si può fare ricorso a
questa opzione.  La questione è comunque abbastanza complessa in quanto, come
sottolinea Stevens in \cite{UNP1}, si distinguono ben quattro casi diversi in
cui è prevista la possibilità di un utilizzo di questa opzione, il che la
rende una delle più difficili da capire.

Il primo caso in cui si fa ricorso a \const{SO\_REUSEADDR}, che è anche il più
comune, è quello in cui un server è terminato ma esistono ancora dei processi
figli che mantengono attiva almeno una connessione remota che utilizza
l'indirizzo locale, mantenendo occupata la porta. Quando si riesegue il server
allora questo riceve un errore sulla chiamata a \func{bind} dato che la porta
è ancora utilizzata in una connessione esistente.\footnote{questa è una delle
  domande più frequenti relative allo sviluppo, in quanto è piuttosto comune
  trovarsi in questa situazione quando si sta sviluppando un server che si
  ferma e si riavvia in continuazione dopo aver fatto modifiche.}  Inoltre se
si usa il protocollo TCP questo può avvenire anche dopo tutti i processi figli
sono terminati, dato che una connessione può restare attiva anche dopo la
chiusura del socket, mantenendosi nello stato \texttt{TIME\_WAIT} (vedi
sez.~\ref{sec:TCP_time_wait}).

\begin{figure}[!htbp]
  \footnotesize \centering
  \begin{minipage}[c]{\codesamplewidth}
    \includecodesample{listati/sockbindopt.c}
  \end{minipage}
  \normalsize
  \caption{Le sezioni della funzione \texttt{sockbindopt} modificate rispetto al
    codice della precedente \texttt{sockbind}.} 
  \label{fig:sockbindopt_code}
\end{figure}

Usando \const{SO\_REUSEADDR} fra la chiamata a \func{socket} e quella a
\func{bind} si consente a quest'ultima di avere comunque successo anche se la
connessione è attiva (o nello stato \texttt{TIME\_WAIT}). È bene però
ricordare (si riveda quanto detto in sez.~\ref{sec:TCP_time_wait}) che la
presenza dello stato \texttt{TIME\_WAIT} ha una ragione, ed infatti se si usa
questa opzione esiste sempre una probabilità, anche se estremamente
remota,\footnote{perché ciò avvenga infatti non solo devono coincidere gli
  indirizzi IP e le porte degli estremi della nuova connessione, ma anche i
  numeri di sequenza dei pacchetti, e questo è estremamente improbabile.}  che
eventuali pacchetti rimasti intrappolati in una precedente connessione possano
finire fra quelli di una nuova.

Come esempio di uso di questa connessione abbiamo predisposto una nuova
versione della funzione \texttt{sockbind} (vedi fig.~\ref{fig:sockbind_code})
che consenta l'impostazione di questa opzione. La nuova funzione è
\texttt{sockbindopt}, e le principali differenze rispetto alla precedente sono
illustrate in fig.~\ref{fig:sockbindopt_code}, dove si sono riportate le
sezioni di codice modificate rispetto alla versione precedente. Il codice
completo della funzione si trova, insieme alle altre funzioni di servizio dei
socket, all'interno del file \texttt{SockUtils.c} dei sorgenti allegati alla
guida.

In realtà tutto quello che si è fatto è stato introdurre nella nuova funzione
(\texttt{\small 1}) un nuovo argomento intero, \param{reuse}, che conterrà il
valore logico da usare nella successiva chiamata (\texttt{\small 14}) a
\func{setsockopt}. Si è poi aggiunta una sezione (\texttt{\small 13--17}) che
esegue l'impostazione dell'opzione fra la chiamata a \func{socket} e quella a
\func{bind}.

\begin{figure}[!htbp] 
  \footnotesize \centering
  \begin{minipage}[c]{\codesamplewidth}
    \includecodesample{listati/TCP_echod_fifth.c}
  \end{minipage}
  \normalsize
  \caption{Il nuovo codice per l'apertura passiva del server \textit{echo} che
    usa la nuova funzione \texttt{sockbindopt}.}
  \label{fig:TCP_echod_fifth}
\end{figure}

A questo punto basterà modificare il  server per utilizzare la nuova
funzione; in fig.~\ref{fig:TCP_echod_fifth} abbiamo riportato le sezioni
modificate rispetto alla precedente versione di
fig.~\ref{fig:TCP_echod_third}. Al solito il codice completo è coi sorgenti
allegati alla guida, nel file \texttt{TCP\_echod\_fifth.c}.

Anche in questo caso si è introdotta (\texttt{\small 8}) una nuova variabile
\var{reuse} che consente di controllare l'uso dell'opzione e che poi sarà
usata (\texttt{\small 14}) come ultimo argomento di \func{setsockopt}. Il
valore di default di questa variabile è nullo, ma usando l'opzione \texttt{-r}
nell'invocazione del server (al solito la gestione delle opzioni non è
riportata in fig.~\ref{fig:TCP_echod_fifth}) se ne potrà impostare ad 1 il
valore, per cui in tal caso la successiva chiamata (\texttt{\small 13--17}) a
\func{setsockopt} attiverà l'opzione \const{SO\_REUSEADDR}.

Il secondo caso in cui viene usata \const{SO\_REUSEADDR} è quando si ha una
macchina cui sono assegnati diversi indirizzi IP (o come suol dirsi
\textit{multi-homed}) e si vuole porre in ascolto sulla stessa porta un
programma diverso (o una istanza diversa dello stesso programma) per indirizzi
IP diversi. Si ricordi infatti che è sempre possibile indicare a \func{bind}
di collegarsi solo su di un indirizzo specifico; in tal caso se un altro
programma cerca di riutilizzare la stessa porta (anche specificando un
indirizzo diverso) otterrà un errore, a meno di non aver preventivamente
impostato \const{SO\_REUSEADDR}.

Usando questa opzione diventa anche possibile eseguire \func{bind}
sull'indirizzo generico, e questo permetterà il collegamento per tutti gli
indirizzi (di quelli presenti) per i quali la porta non risulti occupata da
una precedente chiamata più specifica. Si tenga presente infatti che con il
protocollo TCP non è in genere possibile far partire più server che eseguano
\func{bind} sullo stesso indirizzo e la stessa porta se su di esso c'è già un
socket in ascolto, cioè ottenere quello che viene chiamato un
\itindex{completely~duplicate~binding} \textit{completely duplicate binding}
(per questo è stata introdotta \const{SO\_REUSEPORT}).

Il terzo impiego è simile al precedente e prevede l'uso di \func{bind}
all'interno dello stesso programma per associare indirizzi locali diversi a
socket diversi. In genere questo viene fatto per i socket UDP quando è
necessario ottenere l'indirizzo a cui sono rivolte le richieste del client ed
il sistema non supporta l'opzione \constd{IP\_RECVDSTADDR};\footnote{nel caso
  di Linux questa opzione è stata supportata per in certo periodo nello
  sviluppo del kernel 2.1.x, ma è in seguito stata soppiantata dall'uso di
  \const{IP\_PKTINFO} (vedi sez.~\ref{sec:sock_ipv4_options}).} in tale modo
si può sapere a quale socket corrisponde un certo indirizzo.  Non ha senso
fare questa operazione per un socket TCP dato che su di essi si può sempre
invocare \func{getsockname} una volta che si è completata la connessione.

Infine il quarto caso è quello in cui si vuole effettivamente ottenere un
\textit{completely duplicate binding}, quando cioè si vuole eseguire
\func{bind} su un indirizzo ed una porta che sono già ``\textsl{legati}'' ad
un altro socket.  Come accennato questo con TCP non è possibile, ed ha anche
poco senso pensare di mettere in ascolto due server sulla stessa porta. Se
però si prende in considerazione il traffico in \textit{multicast}, diventa
assolutamente normale che i pacchetti provenienti dal traffico in
\textit{multicast} possano essere ricevuti da più applicazioni o da diverse
istanze della stessa applicazione sulla stessa porte di un indirizzo di
\textit{multicast}.

Un esempio classico è quello di uno streaming di dati (audio, video, ecc.) in
cui l'uso del \textit{multicast} consente di trasmettere un solo pacchetto da
recapitare a tutti i possibili destinatari (invece di inviarne un duplicato a
ciascuno); in questo caso è perfettamente logico aspettarsi che sulla stessa
macchina più utenti possano lanciare un programma che permetta loro di
ricevere gli stessi dati, e quindi effettuare un \textit{completely duplicate
  binding}.

In questo caso utilizzando \const{SO\_REUSEADDR} si può consentire ad una
applicazione eseguire \func{bind} sulla stessa porta ed indirizzo usata da
un'altra, così che anche essa possa ricevere gli stessi pacchetti. Come detto
la cosa non è possibile con i socket TCP (per i quali il \textit{multicast}
comunque non è applicabile), ma lo è per quelli UDP che è il protocollo
normalmente in uso da parte di queste applicazioni. La regola è che quando si
hanno più applicazioni che hanno eseguito \func{bind} sulla stessa porta, di
tutti pacchetti destinati ad un indirizzo di \textit{broadcast} o di
\textit{multicast} viene inviata una copia a ciascuna applicazione. Non è
definito invece cosa accade qualora il pacchetto sia destinato ad un indirizzo
normale (\textit{unicast}).

% TODO documentare SO_REUSEPORT, vedi https://lwn.net/Articles/542260/

\constbeg{SO\_REUSEPORT}

Esistono però dei casi, in particolare per l'uso di programmi
\textit{multithreaded}, in cui può servire un \textit{completely duplicate
  binding} anche per delle ordinarie connessioni TCP.  Per supportare queste
esigenze a partire dal kernel 3.9 è stata introdotta un'altra opzione,
\const{SO\_REUSEPORT} (già presente in altri sistemi come BSD), che consente
di eseguire il \textit{completely duplicate binding}, fintanto che essa venga
specificata per tutti i socket interessati. Come per \const{SO\_REUSEADDR}
sarà possibile usare l'opzione su un socket legato allo stesso indirizzo e
porta solo se il programma che ha eseguito il primo \func{bind} ha impostato
questa opzione.

Nel caso di \const{SO\_REUSEPORT} oltre al fatto che l'opzione deve essere
attivata sul socket prima di chiamare \func{bind} ed attiva su tutti i socket
con \textit{completely duplicate binding}, è richiesto pure che tutti i
processi che si mettono in ascolto sullo stesso indirizzo e porta abbiano lo
stesso \ids{UID} effettivo, per evitare che un altro utente possa ottenere il
relativo traffico (eseguendo quello che viene definito \textit{port hijacking}
o \textit{port stealing}). L'opzione utilizza per \param{optval} un intero
usato come valore logico.

L'opzione si usa sia per socket TCP che UDP, nel primo caso consente un uso
distribuito di \func{accept} in una applicazione \textit{multithreaded}
passando un diverso \textit{listening socket} ad ogni thread, cosa che
migliora le prestazioni rispetto all'approccio tradizionale di usare un thread
per usare \func{accept} e distribuire le connessioni agli altri o avere più
thread che competono per usare \func{accept} sul socket. Nel caso di UDP
l'opzione consente di distribuire meglio i pacchetti su più processi o thread,
rispetto all'approccio tradizionale di far competere gli stessi per l'accesso
in ricezione al socket.

% TODO documentare SO_REUSEPORT introdotta con il kernel 3.9, vedi
% http://git.kernel.org/linus/c617f398edd4db2b8567a28e899a88f8f574798d TODO:
% in cosa differisce da REUSEADDR?

\constend{SO\_REUSEPORT}
\constend{SO\_REUSEADDR}


\constbeg{SO\_LINGER}
\subsubsection{L'opzione \const{SO\_LINGER}}

La terza opzione da approfondire è \const{SO\_LINGER}; essa, come il nome
suggerisce, consente di ``\textsl{indugiare}'' nella chiusura di un socket. Il
comportamento standard sia di \func{close} che \func{shutdown} è infatti
quello di terminare immediatamente dopo la chiamata, mentre il procedimento di
chiusura della connessione (o di un lato di essa) ed il rispettivo invio sulla
rete di tutti i dati ancora presenti nei buffer, viene gestito in sottofondo
dal kernel.

\begin{figure}[!htb]
  \footnotesize \centering
  \begin{minipage}[c]{0.80\textwidth}
    \includestruct{listati/linger.h}
  \end{minipage}
  \caption{La struttura \structd{linger} richiesta come valore dell'argomento
    \param{optval} per l'impostazione dell'opzione dei socket
    \const{SO\_LINGER}.}
  \label{fig:sock_linger_struct}
\end{figure}

L'uso di \const{SO\_LINGER} con \func{setsockopt} permette di modificare (ed
eventualmente ripristinare) questo comportamento in base ai valori passati nei
campi della struttura \struct{linger}, illustrata in
fig.~\ref{fig:sock_linger_struct}.  Fintanto che il valore del campo
\var{l\_onoff} di \struct{linger} è nullo la modalità che viene impostata
(qualunque sia il valore di \var{l\_linger}) è quella standard appena
illustrata; questa combinazione viene utilizzata per riportarsi al
comportamento normale qualora esso sia stato cambiato da una precedente
chiamata.

Se si utilizza un valore di \var{l\_onoff} diverso da zero, il comportamento
alla chiusura viene a dipendere dal valore specificato per il campo
\var{l\_linger}; se quest'ultimo è nullo l'uso delle funzioni \func{close} e
\func{shutdown} provoca la terminazione immediata della connessione: nel caso
di TCP cioè non viene eseguito il procedimento di chiusura illustrato in
sez.~\ref{sec:TCP_conn_term}, ma tutti i dati ancora presenti nel buffer
vengono immediatamente scartati e sulla rete viene inviato un segmento di RST
che termina immediatamente la connessione.

Un esempio di questo comportamento si può abilitare nel nostro client del
servizio \textit{echo} utilizzando l'opzione \texttt{-r}; riportiamo in
fig.~\ref{fig:TCP_echo_sixth} la sezione di codice che permette di introdurre
questa funzionalità; al solito il codice completo è disponibile nei sorgenti
allegati.

\begin{figure}[!htbp] 
  \footnotesize \centering
  \begin{minipage}[c]{\codesamplewidth}
    \includecodesample{listati/TCP_echo_sixth.c}
  \end{minipage}
  \normalsize
  \caption{La sezione del codice del client \textit{echo} che imposta la
    terminazione immediata della connessione in caso di chiusura.}
  \label{fig:TCP_echo_sixth}
\end{figure}

La sezione indicata viene eseguita dopo aver effettuato la connessione e prima
di chiamare la funzione di gestione, cioè fra le righe (\texttt{\small 12}) e
(\texttt{\small 13}) del precedente esempio di fig.~\ref{fig:TCP_echo_fifth}.
Il codice si limita semplicemente a controllare (\texttt{\small 3}) il
valore della variabile \var{reset} che assegnata nella gestione delle opzioni
in corrispondenza all'uso di \texttt{-r} nella chiamata del client. Nel caso
questa sia diversa da zero vengono impostati (\texttt{\small 5--6}) i valori
della struttura \var{ling} che permettono una terminazione immediata della
connessione. Questa viene poi usata nella successiva (\texttt{\small 7})
chiamata a \func{setsockopt}. Al solito si controlla (\texttt{\small 7--10})
il valore di ritorno e si termina il programma in caso di errore, stampandone
il valore.

Infine l'ultima possibilità, quella in cui si utilizza effettivamente
\const{SO\_LINGER} per \textsl{indugiare} nella chiusura, è quella in cui sia
\var{l\_onoff} che \var{l\_linger} hanno un valore diverso da zero. Se si
esegue l'impostazione con questi valori sia \func{close} che \func{shutdown}
si bloccano, nel frattempo viene eseguita la normale procedura di conclusione
della connessione (quella di sez.~\ref{sec:TCP_conn_term}) ma entrambe le
funzioni non ritornano fintanto che non si sia concluso il procedimento di
chiusura della connessione, o non sia passato un numero di
secondi\footnote{questa è l'unità di misura indicata da POSIX ed adottata da
  Linux, altri kernel possono usare unità di misura diverse, oppure usare il
  campo \var{l\_linger} come valore logico (ignorandone il valore) per rendere
  (quando diverso da zero) \func{close} e \func{shutdown} bloccanti fino al
  completamento della trasmissione dei dati sul buffer.}  pari al valore
specificato in \var{l\_linger}.

\constend{SO\_LINGER}



\subsection{Le opzioni per il protocollo IPv4}
\label{sec:sock_ipv4_options}

Il secondo insieme di opzioni dei socket che tratteremo è quello relativo ai
socket che usano il protocollo IPv4; come per le precedenti opzioni generiche
una descrizione di esse è disponibile nella settima sezione delle pagine di
manuale, nel caso specifico la documentazione si può consultare con
\texttt{man 7 ip}.  

\begin{table}[!htb]
  \centering
  \footnotesize
  \begin{tabular}[c]{|l|c|c|c|l|p{4.8cm}|}
    \hline
    \textbf{Opzione}&\texttt{get}&\texttt{set}&\textbf{flag}&\textbf{Tipo}&
                    \textbf{Descrizione}\\
    \hline
    \hline
    \const{IP\_ADD\_MEMBERSHIP} &         &$\bullet$&   &\struct{ip\_mreqn}& 
      Si unisce a un gruppo di \textit{multicast}.\\
    \const{IP\_ADD\_SOURCE\_MEMBERSHIP} & &$\bullet$&   &\struct{ip\_mreq\_source}& 
      Si unisce a un gruppo di \textit{multicast} per una sorgente.\\
    \const{IP\_BLOCK\_SOURCE} & &$\bullet$&   &\struct{ip\_mreq\_source}& 
      Smette di ricevere dati di \textit{multicast} per una sorgente.\\
    \const{IP\_DROP\_MEMBERSHIP}&         &$\bullet$&   &\struct{ip\_mreqn}& 
      Si sgancia da un gruppo di \textit{multicast}.\\
    \const{IP\_DROP\_SOURCE\_MEMBERSHIP}& &$\bullet$&   &\struct{ip\_mreq\_source}& 
      Si sgancia da un gruppo di \textit{multicast} per una sorgente.\\
    \const{IP\_FREEBIND}       &$\bullet$&$\bullet$&$\bullet$&\texttt{int}& 
      Abilita il \textit{binding} a un indirizzo IP non locale ancora non assegnato.\\
    \const{IP\_HDRINCL}         &$\bullet$&$\bullet$&$\bullet$&\texttt{int}& 
      Passa l'intestazione di IP nei dati.\\
    \const{IP\_MINTTL}          &$\bullet$&$\bullet$&   &\texttt{int}& 
      Imposta il valore minimo del TTL per i pacchetti accettati.\\ 
    \const{IP\_MSFILTER}        &$\bullet$&$\bullet$&   &\struct{ip\_msfilter}& 
      Accesso completo all'interfaccia per il filtraggio delle sorgenti
      \textit{multicast}.\\  
    \const{IP\_MTU}             &$\bullet$&         &         &\texttt{int}& 
      Legge il valore attuale della MTU.\\
    \const{IP\_MTU\_DISCOVER}   &$\bullet$&$\bullet$&         &\texttt{int}& 
      Imposta il \textit{Path MTU Discovery}.\\
    \const{IP\_MULTICAST\_ALL}  &$\bullet$&$\bullet$&$\bullet$&\texttt{int}& 
      Imposta l'interfaccia locale di un socket \textit{multicast}.\\ 
    \const{IP\_MULTICAST\_IF}   &         &$\bullet$&   &\struct{ip\_mreqn}& 
      Imposta l'interfaccia locale di un socket \textit{multicast}.\\ 
    \const{IP\_MULTICAST\_LOOP} &$\bullet$&$\bullet$&$\bullet$&\texttt{int}& 
      Controlla il reinvio a se stessi dei dati di \textit{multicast}.\\ 
    \const{IP\_MULTICAST\_TTL}  &$\bullet$&$\bullet$&         &\texttt{int}& 
      Imposta il TTL per i pacchetti \textit{multicast}.\\
    \const{IP\_NODEFRAG}        &$\bullet$&$\bullet$&$\bullet$&\texttt{int}&
      Disabilita il riassemblaggio di pacchetti frammentati.\\
    \const{IP\_OPTIONS}         &$\bullet$&$\bullet$&&\texttt{void *}& %??? 
      Imposta o riceve le opzioni di IP.\\
    \const{IP\_PKTINFO}         &$\bullet$&$\bullet$&$\bullet$&\texttt{int}& 
      Abilita i messaggi di informazione.\\
    \const{IP\_RECVERR}         &$\bullet$&$\bullet$&$\bullet$&\texttt{int}& 
      Abilita i messaggi di errore affidabili.\\
    \const{IP\_RECVOPTS}        &$\bullet$&$\bullet$&$\bullet$&\texttt{int}& 
      Passa un messaggio con le opzioni IP.\\
    \const{IP\_RECVORIGSTADDR}  &$\bullet$&$\bullet$&$\bullet$&\texttt{int}& 
      Abilita i messaggi con l'indirizzo di destinazione originale.\\
    \const{IP\_RECVTOS}         &$\bullet$&$\bullet$&$\bullet$&\texttt{int}& 
      Passa un messaggio col campo TOS.\\
    \const{IP\_RECVTTL}         &$\bullet$&$\bullet$&$\bullet$&\texttt{int}& 
      Abilita i messaggi col campo TTL.\\
    \const{IP\_RETOPTS}         &$\bullet$&$\bullet$&$\bullet$&\texttt{int}& 
      Abilita i messaggi con le opzioni IP non trattate.\\
    \const{IP\_ROUTER\_ALERT}   &$\bullet$&$\bullet$&$\bullet$&\texttt{int}& 
      Imposta l'opzione \textit{IP router alert} sui pacchetti.\\
    \const{IP\_TOS}             &$\bullet$&$\bullet$&         &\texttt{int}& 
      Imposta il valore del campo TOS.\\
    \const{IP\_TRANSPARENT}     &$\bullet$&$\bullet$&$\bullet$&\texttt{int}& 
      Abilita un \textit{proxing} trasparente sul socket.\\
    \const{IP\_TTL}             &$\bullet$&$\bullet$&         &\texttt{int}& 
      Imposta il valore del campo TTL.\\
    \const{IP\_UNBLOCK\_SOURCE} & &$\bullet$&   &\struct{ip\_mreq\_source}& 
      Ricomincia a ricevere dati di \textit{multicast} per una sorgente.\\
   \hline
  \end{tabular}
  \caption{Le opzioni disponibili al livello \const{IPPROTO\_IP}.} 
  \label{tab:sock_opt_iplevel}
\end{table}

Se si vuole operare su queste opzioni generiche il livello da utilizzare è
\constd{IPPROTO\_IP} (o l'equivalente \const{SOL\_IP}); si è riportato un
elenco di queste opzioni in tab.~\ref{tab:sock_opt_iplevel}.  Le costanti
indicanti le opzioni, e tutte le altre costanti ad esse collegate, sono
definite in \headfiled{netinet/ip.h}, ed accessibili includendo detto file.

Le descrizioni riportate in tab.~\ref{tab:sock_opt_iplevel} sono estremamente
succinte, una maggiore quantità di dettagli sulle varie opzioni è fornita nel
seguente elenco:
\begin{basedescript}{\desclabelwidth{2.0cm}\desclabelstyle{\nextlinelabel}}
\item[\constd{IP\_ADD\_MEMBERSHIP}] L'opzione consente di unirsi ad gruppo di
  \textit{multicast}, e può essere usata solo con \func{setsockopt}.
  L'argomento \param{optval} in questo caso deve essere una struttura di tipo
  \struct{ip\_mreqn}, illustrata in fig.~\ref{fig:ip_mreqn_struct}, che
  permette di indicare, con il campo \var{imr\_multiaddr} l'indirizzo del
  gruppo di \textit{multicast} a cui ci si vuole unire, con il campo
  \var{imr\_address} l'indirizzo dell'interfaccia locale con cui unirsi al
  gruppo di \textit{multicast} e con \var{imr\_ifindex} l'indice
  dell'interfaccia da utilizzare (un valore nullo indica una interfaccia
  qualunque).

\begin{figure}[!htb]
  \footnotesize \centering
  \begin{minipage}[c]{0.90\textwidth}
    \includestruct{listati/ip_mreqn.h}
  \end{minipage}
  \caption{La struttura \structd{ip\_mreqn} utilizzata per unirsi a un gruppo di
    \textit{multicast}.}
  \label{fig:ip_mreqn_struct}
\end{figure}

Questa struttura è presente a partire dal kernel 2.2, per compatibilità è
possibile utilizzare anche un argomento di tipo \struct{ip\_mreq}, presente
fino dal kernel 1.2, che differisce da essa soltanto per l'assenza del campo
\var{imr\_ifindex}; il kernel riconosce il tipo di struttura in base alla
differente dimensione passata in \param{optlen}. 

\item[\constd{IP\_ADD\_SOURCE\_MEMBERSHIP}] L'opzione consente di unirsi ad
  gruppo di \textit{multicast}, ricevendo i dati solo da una sorgente
  specifica; come \const{IP\_ADD\_MEMBERSHIP} può essere usata solo con
  \func{setsockopt}.

\begin{figure}[!htb]
  \footnotesize \centering
  \begin{minipage}[c]{0.90\textwidth}
    \includestruct{listati/ip_mreq_source.h}
  \end{minipage}
  \caption{La struttura \structd{ip\_mreqn} utilizzata per unirsi a un gruppo di
    \textit{multicast} per una specifica sorgente.}
  \label{fig:ip_mreq_source_struct}
\end{figure}

L'argomento \param{optval} in questo caso è una struttura
\struct{ip\_mreq\_source} illustrata in fig.~\ref{fig:ip_mreq_source_struct},
dove il campo \var{imr\_multiaddr} è l'indirizzo del gruppo di
\textit{multicast}, il campo \var{imr\_interface} l'indirizzo dell'interfaccia
locale che deve essere usata per aggiungersi al gruppo di \textit{multicast} e
il campo \var{imr\_sourceaddr} l'indirizzo della sorgente da cui
l'applicazione vuole ricevere i dati. L'opzione può essere ripetuta più volte
per collegarsi a diverse sorgenti.

\item[\constd{IP\_BLOCK\_SOURCE}] Questa opzione, disponibile dal kernel
  2.4.22, consente di smettere di ricevere dati di \textit{multicast} dalla
  sorgente (e relativo gruppo) specificati dalla struttura
  \struct{ip\_mreq\_source} (vedi fig.~\ref{fig:ip_mreq_source_struct})
  passata come argomento \param{optval}. L'opzione è utilizzabile solo se si è
  già registrati nel gruppo di \textit{multicast} indicato con un precedente
  utilizzo di \const{IP\_ADD\_MEMBERSHIP} o
  \const{IP\_ADD\_SOURCE\_MEMBERSHIP}.

\item[\constd{IP\_DROP\_MEMBERSHIP}] Lascia un gruppo di \textit{multicast},
  prende per \param{optval} la stessa struttura \struct{ip\_mreqn} (o
  \struct{ip\_mreq}) usata per \const{IP\_ADD\_MEMBERSHIP} (vedi
  fig.~\ref{fig:ip_mreqn_struct}).

\item[\constd{IP\_DROP\_SOURCE\_MEMBERSHIP}] Lascia un gruppo di
  \textit{multicast} per una specifica sorgente, prende per \param{optval} la
  stessa struttura \struct{ip\_mreq\_source} usata per
  \const{IP\_ADD\_SOURCE\_MEMBERSHIP} (vedi
  fig.~\ref{fig:ip_mreq_source_struct}). Se ci si è registrati per più
  sorgenti nello stesso gruppo, si continuerà a ricevere dati sulle altre. Per
  smettere di ricevere dati da tutte le sorgenti occorre usare l'opzione
  \const{IP\_LEAVE\_GROUP}.

\item[\constd{IP\_FREEBIND}] Se abilitata questa opzione, disponibile dal
  kernel 2.4, consente di usare \func{bind} anche su un indirizzo IP non
  locale o che ancora non è stato assegnato. Questo permette ad una
  applicazione di mettersi in ascolto su un socket prima che l'interfaccia
  sottostante o l'indirizzo che questa deve usare sia stato configurato. È
  l'equivalente a livello di singolo socket dell'uso della \textit{sysctl}
  \texttt{ip\_nonlocal\_bind} che vedremo in
  sez.~\ref{sec:sock_ipv4_sysctl}. Prende per
  \param{optval} un intero usato come valore logico.

\item[\constd{IP\_HDRINCL}] Se viene abilitata questa opzione, presente dal
  kernel 2.0, l'utente deve fornire lui stesso l'intestazione del protocollo
  IP in testa ai propri dati. L'opzione è valida soltanto per socket di tipo
  \const{SOCK\_RAW}, e quando utilizzata eventuali valori impostati con
  \const{IP\_OPTIONS}, \const{IP\_TOS} o \const{IP\_TTL} sono ignorati. In
  ogni caso prima della spedizione alcuni campi dell'intestazione vengono
  comunque modificati dal kernel, torneremo sull'argomento in
  sez.~\ref{sec:socket_raw}.  Prende per \param{optval} un intero usato come
  valore logico.

\item[\constd{IP\_MSFILTER}] L'opzione, introdotta con il kernel 2.4.22,
  fornisce accesso completo all'interfaccia per il filtraggio delle sorgenti
  di \textit{multicast} (il cosiddetto \textit{multicast source filtering},
  definito dall'\href{http://www.ietf.org/rfc/rfc3376.txt}{RFC~3376}).

\begin{figure}[!htb]
  \footnotesize \centering
  \begin{minipage}[c]{0.90\textwidth}
    \includestruct{listati/ip_msfilter.h}
  \end{minipage}
  \caption{La struttura \structd{ip\_msfilter} utilizzata per il
    \textit{multicast source filtering}.}
  \label{fig:ip_msfilter_struct}
\end{figure}

L'argomento \param{optval} deve essere una struttura di tipo
\struct{ip\_msfilter} (illustrata in fig.~\ref{fig:ip_msfilter_struct}); il
campo \var{imsf\_multiaddr} indica l'indirizzo del gruppo di
\textit{multicast}, il campo \var{imsf\_interface} l'indirizzo
dell'interfaccia locale, il campo \var{imsf\_mode} indica la modalità di
filtraggio e con i campi \var{imsf\_numsrc} e \var{imsf\_slist}
rispettivamente la lunghezza della lista, e la lista stessa, degli indirizzi
sorgente.

Come ausilio all'uso di questa opzione sono disponibili le macro
\macro{MCAST\_INCLUDE} e \macro{MCAST\_EXCLUDE} che si possono usare per
\var{imsf\_mode}. Inoltre si può usare la macro
\macro{IP\_MSFILTER\_SIZE}\texttt{(\textsf{n})} per determinare il valore
di \param{optlen} con una struttura \struct{ip\_msfilter} contenente
\texttt{\textsf{n}} sorgenti in \var{imsf\_slist}.

\item[\constd{IP\_MINTTL}] L'opzione, introdotta con il kernel 2.6.34, imposta
  un valore minimo per il campo \textit{Time to Live} dei pacchetti associati
  al socket su cui è attivata, che se non rispettato ne causa lo scarto
  automatico. L'opzione è nata per implementare
  l'\href{http://www.ietf.org/rfc/rfc5082.txt}{RFC~5082} che la prevede come
  forma di protezione per i router che usano il protocollo BGP poiché questi,
  essendo in genere adiacenti, possono, impostando un valore di 255, scartare
  automaticamente tutti gli eventuali pacchetti falsi creati da un attacco a
  questo protocollo, senza doversi curare di verificarne la
  validità.\footnote{l'attacco viene in genere portato per causare un
    \textit{Denial of Service} aumentando il consumo di CPU del router nella
    verifica dell'autenticità di un gran numero di pacchetti falsi; questi,
    arrivando da sorgenti diverse da un router adiacente, non potrebbero più
    avere un TTL di 255 anche qualora questo fosse stato il valore di
    partenza, e l'impostazione dell'opzione consente di scartarli senza carico
    aggiuntivo sulla CPU (che altrimenti dovrebbe calcolare una checksum).}

\itindbeg{Path~MTU}
\item[\constd{IP\_MTU}] Permette di leggere il valore della \textit{Path MTU}
  del socket.  L'opzione richiede per \param{optval} un intero che conterrà il
  valore della \textit{Path MTU} in byte.  Questa è una opzione introdotta con
  i kernel della serie 2.2.x, ed è specifica di Linux.

  È tramite questa opzione che un programma può leggere, quando si è avuto un
  errore di \errval{EMSGSIZE}, il valore della MTU corrente del socket. Si
  tenga presente che per poter usare questa opzione, oltre ad avere abilitato
  la scoperta della \textit{Path MTU}, occorre che il socket sia stato
  esplicitamente connesso con \func{connect}. 

  Ad esempio con i socket UDP si può ottenere una stima iniziale della
  \textit{Path MTU} eseguendo prima una \func{connect} verso la destinazione,
  e poi usando \func{getsockopt} con questa opzione. Si può anche avviare
  esplicitamente il procedimento di scoperta inviando un pacchetto di grosse
  dimensioni (che verrà scartato) e ripetendo l'invio coi dati aggiornati. Si
  tenga infine conto che durante il procedimento i pacchetti iniziali possono
  essere perduti, ed è compito dell'applicazione gestirne una eventuale
  ritrasmissione.

\item[\constd{IP\_MTU\_DISCOVER}] Questa è una opzione introdotta con i kernel
  della serie 2.2.x, ed è specifica di Linux.  L'opzione permette di scrivere
  o leggere le impostazioni della modalità usata per la determinazione della
  \textit{Path MTU} (vedi sez.~\ref{sec:net_lim_dim}) del
  socket. L'opzione prende per \param{optval} un valore intero che indica la
  modalità usata, da specificare con una delle costanti riportate in
  tab.~\ref{tab:sock_ip_mtu_discover}.

  \begin{table}[!htb]
    \centering
    \footnotesize
    \begin{tabular}[c]{|l|r|p{7cm}|}
      \hline
      \multicolumn{2}{|c|}{\textbf{Valore}}&\textbf{Significato} \\
      \hline
      \hline
      \constd{IP\_PMTUDISC\_DONT}&0& Non effettua la ricerca dalla \textit{Path
                                     MTU}.\\
      \constd{IP\_PMTUDISC\_WANT}&1& Utilizza il valore impostato per la rotta
                                     utilizzata dai pacchetti (dal comando
                                     \texttt{route}).\\ 
      \constd{IP\_PMTUDISC\_DO}  &2& Esegue la procedura di determinazione
                                     della \textit{Path MTU} come richiesto
                                     dall'\href{http://www.ietf.org/rfc/rfc1191.txt}{RFC~1191}.\\ 
      \constd{IP\_PMTUDISC\_PROBE}&?& .\\
      \hline
    \end{tabular}
    \caption{Valori possibili per l'argomento \param{optval} di
      \const{IP\_MTU\_DISCOVER}.} 
    \label{tab:sock_ip_mtu_discover}
  \end{table}

  Il valore di default applicato ai socket di tipo \const{SOCK\_STREAM} è
  determinato dal parametro \texttt{ip\_no\_pmtu\_disc} (vedi
  sez.~\ref{sec:sock_sysctl}), mentre per tutti gli altri socket di default la
  ricerca è disabilitata ed è responsabilità del programma creare pacchetti di
  dimensioni appropriate e ritrasmettere eventuali pacchetti persi. Se
  l'opzione viene abilitata, il kernel si incaricherà di tenere traccia
  automaticamente della \textit{Path MTU} verso ciascuna destinazione, e
  rifiuterà immediatamente la trasmissione di pacchetti di dimensioni maggiori
  della MTU con un errore di \errval{EMSGSIZE}.\footnote{in caso contrario la
    trasmissione del pacchetto sarebbe effettuata, ottenendo o un fallimento
    successivo della trasmissione, o la frammentazione dello stesso.}
\itindend{Path~MTU}

\item[\constd{IP\_MULTICAST\_IF}] Imposta l'interfaccia locale per l'utilizzo
  del \textit{multicast}, ed utilizza come \param{optval} le stesse strutture
  \struct{ip\_mreqn} o \struct{ip\_mreq} delle due precedenti opzioni.

\item[\constd{IP\_MULTICAST\_LOOP}] L'opzione consente di decidere se i dati
  che si inviano su un socket usato con il \textit{multicast} vengano ricevuti
  anche sulla stessa macchina da cui li si stanno inviando.  Prende per
  \param{optval} un intero usato come valore logico.

  In generale se si vuole che eventuali client possano ricevere i dati che si
  inviano occorre che questa funzionalità sia abilitata (come avviene di
  default). Qualora però non si voglia generare traffico per dati che già sono
  disponibili in locale l'uso di questa opzione permette di disabilitare
  questo tipo di traffico.

\item[\constd{IP\_MULTICAST\_TTL}] L'opzione permette di impostare o leggere il
  valore del campo TTL per i pacchetti \textit{multicast} in uscita associati
  al socket. È importante che questo valore sia il più basso possibile, ed il
  default è 1, che significa che i pacchetti non potranno uscire dalla rete
  locale. Questa opzione consente ai programmi che lo richiedono di superare
  questo limite.  L'opzione richiede per
  \param{optval} un intero che conterrà il valore del TTL.
% TODO chiarire quale è la struttura \struct{ip\_mreq}

\item[\constd{IP\_OPTIONS}] l'opzione permette di impostare o leggere le
  opzioni del protocollo IP (si veda sez.~\ref{sec:IP_options}). L'opzione
  prende come valore dell'argomento \param{optval} un puntatore ad un buffer
  dove sono mantenute le opzioni, mentre \param{optlen} indica la dimensione
  di quest'ultimo. Quando la si usa con \func{getsockopt} vengono lette le
  opzioni IP utilizzate per la spedizione, quando la si usa con
  \func{setsockopt} vengono impostate le opzioni specificate. L'uso di questa
  opzione richiede una profonda conoscenza del funzionamento del protocollo,
  torneremo in parte sull'argomento in sez.~\ref{sec:sock_IP_options}.

\item[\constd{IP\_PKTINFO}] Quando abilitata l'opzione permette di ricevere
  insieme ai pacchetti un messaggio ancillare (vedi
  sez.~\ref{sec:net_ancillary_data}) di tipo \const{IP\_PKTINFO} contenente
  una struttura \struct{pktinfo} (vedi fig.~\ref{fig:sock_pktinfo_struct}) che
  mantiene una serie di informazioni riguardo i pacchetti in arrivo. In
  particolare è possibile conoscere l'interfaccia su cui è stato ricevuto un
  pacchetto (nel campo \var{ipi\_ifindex}),\footnote{in questo campo viene
    restituito il valore numerico dell'indice dell'interfaccia,
    sez.~\ref{sec:sock_ioctl_netdevice}.} l'indirizzo locale da esso
  utilizzato (nel campo \var{ipi\_spec\_dst}) e l'indirizzo remoto dello
  stesso (nel campo \var{ipi\_addr}).

\begin{figure}[!htb]
  \footnotesize \centering
  \begin{minipage}[c]{0.80\textwidth}
    \includestruct{listati/pktinfo.h}
  \end{minipage}
  \caption{La struttura \structd{pktinfo} usata dall'opzione
    \const{IP\_PKTINFO} per ricavare informazioni sui pacchetti di un socket
    di tipo \const{SOCK\_DGRAM}.}
  \label{fig:sock_pktinfo_struct}
\end{figure}


L'opzione è utilizzabile solo per socket di tipo \const{SOCK\_DGRAM}. Questa è
una opzione introdotta con i kernel della serie 2.2.x, ed è specifica di
Linux;\footnote{non dovrebbe pertanto essere utilizzata se si ha a cuore la
  portabilità.} essa permette di sostituire le opzioni \const{IP\_RECVDSTADDR}
e \const{IP\_RECVIF} presenti in altri Unix (la relativa informazione è quella
ottenibile rispettivamente dai campi \var{ipi\_addr} e \var{ipi\_ifindex} di
\struct{pktinfo}). 

L'opzione prende per \param{optval} un intero usato come valore logico, che
specifica soltanto se insieme al pacchetto deve anche essere inviato o
ricevuto il messaggio \const{IP\_PKTINFO} (vedi
sez.~\ref{sec:net_ancillary_data}); il messaggio stesso dovrà poi essere
letto o scritto direttamente con \func{recvmsg} e \func{sendmsg} (vedi
sez.~\ref{sec:net_sendmsg}).

\item[\constd{IP\_RECVERR}] Questa è una opzione introdotta con i kernel della
  serie 2.2.x, ed è specifica di Linux. Essa permette di usufruire di un
  meccanismo affidabile per ottenere un maggior numero di informazioni in caso
  di errori. Se l'opzione è abilitata tutti gli errori generati su un socket
  vengono memorizzati su una coda, dalla quale poi possono essere letti con
  \func{recvmsg} (vedi sez.~\ref{sec:net_sendmsg}) come messaggi ancillari
  (torneremo su questo in sez.~\ref{sec:net_ancillary_data}) di tipo
  \const{IP\_RECVERR}.  L'opzione richiede per \param{optval} un intero usato
  come valore logico e non è applicabile a socket di tipo
  \const{SOCK\_STREAM}.

\item[\constd{IP\_RECVOPTS}] Quando abilitata l'opzione permette di ricevere
  insieme ai pacchetti un messaggio ancillare (vedi
  sez.~\ref{sec:net_ancillary_data}) di tipo \const{IP\_OPTIONS}, contenente
  le opzioni IP del protocollo (vedi sez.~\ref{sec:IP_options}). Le
  intestazioni di instradamento e le altre opzioni sono già riempite con i
  dati locali. L'opzione richiede per \param{optval} un intero usato come
  valore logico.  L'opzione non è supportata per socket di tipo
  \const{SOCK\_STREAM}.

\item[\constd{IP\_RECVTOS}] Quando abilitata l'opzione permette di ricevere
  insieme ai pacchetti un messaggio ancillare (vedi
  sez.~\ref{sec:net_ancillary_data}) di tipo \const{IP\_TOS}, che contiene un
  byte con il valore del campo \textit{Type of Service} dell'intestazione IP
  del pacchetto stesso (vedi sez.~\ref{sec:IP_header}).  Prende per
  \param{optval} un intero usato come valore logico.

\item[\constd{IP\_RECVTTL}] Quando abilitata l'opzione permette di ricevere
  insieme ai pacchetti un messaggio ancillare (vedi
  sez.~\ref{sec:net_ancillary_data}) di tipo \const{IP\_RECVTTL}, contenente
  un byte con il valore del campo \textit{Time to Live} dell'intestazione IP
  (vedi sez.~\ref{sec:IP_header}).  L'opzione richiede per \param{optval} un
  intero usato come valore logico. L'opzione non è supportata per socket di
  tipo \const{SOCK\_STREAM}.

\item[\constd{IP\_RETOPTS}] Identica alla precedente \const{IP\_RECVOPTS}, ma
  in questo caso restituisce i dati grezzi delle opzioni, senza che siano
  riempiti i capi di instradamento e le marche temporali.  L'opzione richiede
  per \param{optval} un intero usato come valore logico.  L'opzione non è
  supportata per socket di tipo \const{SOCK\_STREAM}.

\item[\constd{IP\_ROUTER\_ALERT}] Questa è una opzione introdotta con i
  kernel della serie 2.2.x, ed è specifica di Linux. Prende per
  \param{optval} un intero usato come valore logico. Se abilitata
  passa tutti i pacchetti con l'opzione \textit{IP Router Alert} (vedi
  sez.~\ref{sec:IP_options}) che devono essere inoltrati al socket
  corrente. Può essere usata soltanto per socket di tipo raw.

\item[\constd{IP\_TOS}] L'opzione consente di leggere o impostare il campo
  \textit{Type of Service} dell'intestazione IP (per una trattazione più
  dettagliata, che riporta anche i valori possibili e le relative costanti di
  definizione si veda sez.~\ref{sec:IP_header}) che permette di indicare le
  priorità dei pacchetti. Se impostato il valore verrà mantenuto per tutti i
  pacchetti del socket; alcuni valori (quelli che aumentano la priorità)
  richiedono i privilegi di amministrazione con la \textit{capability}
  \const{CAP\_NET\_ADMIN}.

  Il campo TOS è di 8 bit e l'opzione richiede per \param{optval} un intero
  che ne contenga il valore. Sono definite anche alcune costanti che
  definiscono alcuni valori standardizzati per il \textit{Type of Service},
  riportate in tab.~\ref{tab:IP_TOS_values}, il valore di default usato da
  Linux è \const{IPTOS\_LOWDELAY}, ma esso può essere modificato con le
  funzionalità del cosiddetto \textit{Advanced Routing}. Si ricordi che la
  priorità dei pacchetti può essere impostata anche in maniera indipendente
  dal protocollo utilizzando l'opzione \const{SO\_PRIORITY} illustrata in
  sez.~\ref{sec:sock_generic_options}.

\item[\constd{IP\_TTL}] L'opzione consente di leggere o impostare per tutti i
  pacchetti associati al socket il campo \textit{Time to Live}
  dell'intestazione IP che indica il numero massimo di \textit{hop} (passaggi
  da un router ad un altro) restanti al paccheto (per una trattazione più
  estesa si veda sez.~\ref{sec:IP_header}).  Il campo TTL è di 8 bit e
  l'opzione richiede che \param{optval} sia un intero, che ne conterrà il
  valore.
\end{basedescript}



\subsection{Le opzioni per i protocolli TCP e UDP}
\label{sec:sock_tcp_udp_options}

In questa sezione tratteremo le varie opzioni disponibili per i socket che
usano i due principali protocolli di comunicazione del livello di trasporto;
UDP e TCP.\footnote{come per le precedenti, una descrizione di queste opzioni
  è disponibile nella settima sezione delle pagine di manuale, che si può
  consultare rispettivamente con \texttt{man 7 tcp} e \texttt{man 7 udp}; le
  pagine di manuale però, alla stesura di questa sezione (Agosto 2006) sono
  alquanto incomplete.}  Dato che questi due protocolli sono entrambi
trasportati su IP,\footnote{qui si sottintende IPv4, ma le opzioni per TCP e
  UDP sono le stesse anche quando si usa IPv6.} oltre alle opzioni generiche
di sez.~\ref{sec:sock_generic_options} saranno comunque disponibili anche le
precedenti opzioni di sez.~\ref{sec:sock_ipv4_options}.\footnote{in realtà in
  sez.~\ref{sec:sock_ipv4_options} si sono riportate le opzioni per IPv4, al
  solito, qualora si stesse utilizzando IPv6, si potrebbero utilizzare le
  opzioni di quest'ultimo.}

Il protocollo che supporta il maggior numero di opzioni è TCP; per poterle
utilizzare occorre specificare \const{SOL\_TCP} (o l'equivalente
\constd{IPPROTO\_TCP}) come valore per l'argomento \param{level}. Si sono
riportate le varie opzioni disponibili in tab.~\ref{tab:sock_opt_tcplevel},
dove sono elencate le rispettive costanti da utilizzare come valore per
l'argomento \param{optname}. Dette costanti e tutte le altre costanti e
strutture collegate all'uso delle opzioni TCP sono definite in
\headfiled{netinet/tcp.h}, ed accessibili includendo detto file.\footnote{in
  realtà questo è il file usato dalle librerie; la definizione delle opzioni
  effettivamente supportate da Linux si trova nel file
  \texttt{include/linux/tcp.h} dei sorgenti del kernel, dal quale si sono
  estratte le costanti di tab.~\ref{tab:sock_opt_tcplevel}.}

\begin{table}[!htb]
  \centering
  \footnotesize
  \begin{tabular}[c]{|l|c|c|c|l|l|}
    \hline
    \textbf{Opzione}&\texttt{get}&\texttt{set}&\textbf{flag}&\textbf{Tipo}&
                    \textbf{Descrizione}\\
    \hline
    \hline
    \const{TCP\_NODELAY}      &$\bullet$&$\bullet$&$\bullet$&\texttt{int}& 
      Spedisce immediatamente i dati in segmenti singoli.\\
    \const{TCP\_MAXSEG}       &$\bullet$&$\bullet$&         &\texttt{int}&
      Valore della MSS per i segmenti in uscita.\\  
    \const{TCP\_CORK}         &$\bullet$&$\bullet$&$\bullet$&\texttt{int}&
      Accumula i dati in un unico segmento.\\
    \const{TCP\_KEEPIDLE}     &$\bullet$&$\bullet$&         &\texttt{int}& 
      Tempo in secondi prima di inviare un \textit{keepalive}.\\
    \const{TCP\_KEEPINTVL}    &$\bullet$&$\bullet$&         &\texttt{int}&
      Tempo in secondi prima fra \textit{keepalive} successivi.\\
    \const{TCP\_KEEPCNT}      &$\bullet$&$\bullet$&         &\texttt{int}& 
      Numero massimo di \textit{keepalive} inviati.\\
    \const{TCP\_SYNCNT}       &$\bullet$&$\bullet$&         &\texttt{int}& 
      Numero massimo di ritrasmissioni di un SYN.\\
    \const{TCP\_LINGER2}      &$\bullet$&$\bullet$&         &\texttt{int}&
      Tempo di vita in stato \texttt{FIN\_WAIT2}.\\
    \const{TCP\_DEFER\_ACCEPT}&$\bullet$&$\bullet$&         &\texttt{int}&
      Ritorna da \func{accept} solo in presenza di dati.\\
    \const{TCP\_WINDOW\_CLAMP}&$\bullet$&$\bullet$&         &\texttt{int}&
      Valore della \textit{advertised window}.\\
    \const{TCP\_INFO}         &$\bullet$&        &       &\struct{tcp\_info}& 
      Restituisce informazioni sul socket.\\
    \const{TCP\_QUICKACK}     &$\bullet$&$\bullet$&$\bullet$&\texttt{int}&
      Abilita la modalità \textit{quickack}.\\
    \const{TCP\_CONGESTION}   &$\bullet$&$\bullet$&         &\texttt{char *}&
      Imposta l'algoritmo per il controllo della congestione.\\
   \hline
  \end{tabular}
  \caption{Le opzioni per i socket TCP disponibili al livello
    \const{SOL\_TCP}.}
  \label{tab:sock_opt_tcplevel}
\end{table}

Le descrizioni delle varie opzioni riportate in
tab.~\ref{tab:sock_opt_tcplevel} sono estremamente sintetiche ed indicative,
la spiegazione del funzionamento delle singole opzioni con una maggiore
quantità di dettagli è fornita nel seguente elenco:
\begin{basedescript}{\desclabelwidth{1.5cm}\desclabelstyle{\nextlinelabel}}

\item[\constd{TCP\_NODELAY}] il protocollo TCP utilizza un meccanismo di
  bufferizzazione dei dati uscenti, per evitare la trasmissione di tanti
  piccoli segmenti con un utilizzo non ottimale della banda
  disponibile.\footnote{il problema è chiamato anche
    \itindex{silly~window~syndrome} \textit{silly window syndrome}, per averne
    un'idea si pensi al risultato che si ottiene quando un programma di
    terminale invia un segmento TCP per ogni tasto premuto, 40 byte di
    intestazione di protocollo con 1 byte di dati trasmessi; per evitare
    situazioni del genere è stato introdotto \index{algoritmo~di~Nagle}
    l'\textsl{algoritmo di Nagle}.}  Questo meccanismo è controllato da un
  apposito algoritmo (detto \textsl{algoritmo di Nagle}, vedi
  sez.~\ref{sec:tcp_protocol_xxx}).  Il comportamento normale del protocollo
  prevede che i dati siano accumulati fintanto che non si raggiunge una
  quantità considerata adeguata per eseguire la trasmissione di un singolo
  segmento.

  Ci sono però delle situazioni in cui questo comportamento può non essere
  desiderabile, ad esempio quando si sa in anticipo che l'applicazione invierà
  soltanto un piccolo quantitativo di dati;\footnote{è il caso classico di una
    richiesta HTTP.} in tal caso l'attesa introdotta dall'algoritmo di
  bufferizzazione non soltanto è inutile, ma peggiora le prestazioni
  introducendo un ritardo.  Impostando questa opzione si disabilita l'uso
  dell'\textsl{algoritmo di Nagle} ed i dati vengono inviati immediatamente in
  singoli segmenti, qualunque sia la loro dimensione.  Ovviamente l'uso di
  questa opzione è dedicato a chi ha esigenze particolari come quella
  illustrata, che possono essere stabilite solo per la singola applicazione.

  Si tenga conto che questa opzione viene sovrascritta dall'eventuale
  impostazione dell'opzione \const{TCP\_CORK} (il cui scopo è sostanzialmente
  l'opposto) che blocca l'invio immediato. Tuttavia quando la si abilita viene
  sempre forzato lo scaricamento della coda di invio (con conseguente
  trasmissione di tutti i dati pendenti), anche qualora si fosse già abilitata
  \const{TCP\_CORK}.\footnote{si tenga presente però che \const{TCP\_CORK} può
    essere specificata insieme a \const{TCP\_NODELAY} soltanto a partire dal
    kernel 2.5.71.}

\item[\constd{TCP\_MAXSEG}] con questa opzione si legge o si imposta il valore
  della MSS (\textit{Maximum Segment Size}, vedi sez.~\ref{sec:net_lim_dim} e
  sez.~\ref{sec:tcp_protocol_xxx}) dei segmenti TCP uscenti. Se l'opzione è
  impostata prima di stabilire la connessione, si cambia anche il valore della
  MSS annunciata all'altro capo della connessione. Se si specificano valori
  maggiori della MTU questi verranno ignorati, inoltre TCP imporrà anche i
  suoi limiti massimo e minimo per questo valore.

\item[\constd{TCP\_CORK}] questa opzione è il complemento naturale di
  \const{TCP\_NODELAY} e serve a gestire a livello applicativo la situazione
  opposta, cioè quella in cui si sa fin dal principio che si dovranno inviare
  grosse quantità di dati. Anche in questo caso l'\textsl{algoritmo di Nagle}
  tenderà a suddividerli in dimensioni da lui ritenute
  opportune,\footnote{l'algoritmo cerca di tenere conto di queste situazioni,
    ma essendo un algoritmo generico tenderà comunque ad introdurre delle
    suddivisioni in segmenti diversi, anche quando potrebbero non essere
    necessarie, con conseguente spreco di banda.}  ma sapendo fin dall'inizio
  quale è la dimensione dei dati si potranno di nuovo ottenere delle migliori
  prestazioni disabilitandolo, e gestendo direttamente l'invio del nostro
  blocco di dati in soluzione unica.

  Quando questa opzione viene abilitata non vengono inviati segmenti di dati
  fintanto che essa non venga disabilitata; a quel punto tutti i dati rimasti
  in coda saranno inviati in un solo segmento TCP. In sostanza con questa
  opzione si può controllare il flusso dei dati mettendo una sorta di
  ``\textsl{tappo}'' (da cui il nome in inglese) al flusso di uscita, in modo
  ottimizzare a mano l'uso della banda. Si tenga presente che per l'effettivo
  funzionamento ci si deve ricordare di disattivare l'opzione al termine
  dell'invio del blocco dei dati.

  Si usa molto spesso \const{TCP\_CORK} quando si effettua il trasferimento
  diretto di un blocco di dati da un file ad un socket con \func{sendfile}
  (vedi sez.~\ref{sec:file_sendfile_splice}), per inserire una intestazione
  prima della chiamata a questa funzione; senza di essa l'intestazione
  potrebbe venire spedita in un segmento a parte, che a seconda delle
  condizioni potrebbe richiedere anche una risposta di ACK, portando ad una
  notevole penalizzazione delle prestazioni.

  Si tenga presente che l'implementazione corrente di \const{TCP\_CORK} non
  consente di bloccare l'invio dei dati per più di 200 millisecondi, passati i
  quali i dati accumulati in coda sanno inviati comunque.  Questa opzione è
  tipica di Linux\footnote{l'opzione è stata introdotta con i kernel della
    serie 2.4.x.} e non è disponibile su tutti i kernel unix-like, pertanto
  deve essere evitata se si vuole scrivere codice portabile.

\item[\constd{TCP\_KEEPIDLE}] con questa opzione si legge o si imposta
  l'intervallo di tempo, in secondi, che deve trascorrere senza traffico sul
  socket prima che vengano inviati, qualora si sia attivata su di esso
  l'opzione \const{SO\_KEEPALIVE}, i messaggi di \textit{keep-alive} (si veda
  la trattazione relativa al \textit{keep-alive} in
  sez.~\ref{sec:sock_options_main}).  Anche questa opzione non è disponibile
  su tutti i kernel unix-like e deve essere evitata se si vuole scrivere
  codice portabile.

\item[\constd{TCP\_KEEPINTVL}] con questa opzione si legge o si imposta
  l'intervallo di tempo, in secondi, fra due messaggi di \textit{keep-alive}
  successivi (si veda sempre quanto illustrato in
  sez.~\ref{sec:sock_options_main}). Come la precedente non è disponibile su
  tutti i kernel unix-like e deve essere evitata se si vuole scrivere codice
  portabile.

\item[\constd{TCP\_KEEPCNT}] con questa opzione si legge o si imposta il numero
  totale di messaggi di \textit{keep-alive} da inviare prima di concludere che
  la connessione è caduta per assenza di risposte ad un messaggio di
  \textit{keep-alive} (di nuovo vedi sez.~\ref{sec:sock_options_main}). Come
  la precedente non è disponibile su tutti i kernel unix-like e deve essere
  evitata se si vuole scrivere codice portabile.

\item[\constd{TCP\_SYNCNT}] con questa opzione si legge o si imposta il numero
  di tentativi di ritrasmissione dei segmenti SYN usati nel \textit{three way
    handshake} prima che il tentativo di connessione venga abortito (si
  ricordi quanto accennato in sez.~\ref{sec:TCP_func_connect}). Sovrascrive
  per il singolo socket il valore globale impostato con la \textit{sysctl}
  \texttt{tcp\_syn\_retries} (vedi sez.~\ref{sec:sock_ipv4_sysctl}).  Non
  vengono accettati valori maggiori di 255; anche questa opzione non è
  standard e deve essere evitata se si vuole scrivere codice portabile.

\item[\constd{TCP\_LINGER2}] con questa opzione si legge o si imposta, in
  numero di secondi, il tempo di sussistenza dei socket terminati nello stato
  \texttt{FIN\_WAIT2} (si ricordi quanto visto in
  sez.~\ref{sec:TCP_conn_term}).\footnote{si tenga ben presente che questa
    opzione non ha nulla a che fare con l'opzione \const{SO\_LINGER} che
    abbiamo visto in sez.~\ref{sec:sock_options_main}.}  Questa opzione
  consente di sovrascrivere per il singolo socket il valore globale impostato
  con la \textit{sysctl} \texttt{tcp\_fin\_timeout} (vedi
  sez.~\ref{sec:sock_ipv4_sysctl}).  Anche questa opzione è da evitare se si
  ha a cuore la portabilità del codice.

\item[\constd{TCP\_DEFER\_ACCEPT}] questa opzione consente di modificare il
  comportamento standard del protocollo TCP nello stabilirsi di una
  connessione; se ricordiamo il meccanismo del \textit{three way handshake}
  illustrato in fig.~\ref{fig:TCP_TWH} possiamo vedere che in genere un client
  inizierà ad inviare i dati ad un server solo dopo l'emissione dell'ultimo
  segmento di ACK.

  Di nuovo esistono situazioni (e la più tipica è quella di una richiesta
  HTTP) in cui sarebbe utile inviare immediatamente la richiesta all'interno
  del segmento con l'ultimo ACK del \textit{three way handshake}; si potrebbe
  così risparmiare l'invio di un segmento successivo per la richiesta e il
  ritardo sul server fra la ricezione dell'ACK e quello della richiesta.

  Se si invoca \const{TCP\_DEFER\_ACCEPT} su un socket dal lato client (cioè
  dal lato da cui si invoca \func{connect}) si istruisce il kernel a non
  inviare immediatamente l'ACK finale del \textit{three way handshake},
  attendendo per un po' di tempo la prima scrittura, in modo da inviare i dati
  di questa insieme col segmento ACK.  Chiaramente la correttezza di questo
  comportamento dipende in maniera diretta dal tipo di applicazione che usa il
  socket; con HTTP, che invia una breve richiesta, permette di risparmiare un
  segmento, con FTP, in cui invece si attende la ricezione del prompt del
  server, introduce un inutile ritardo.

  Allo stesso tempo il protocollo TCP prevede che sul lato del server la
  funzione \func{accept} ritorni dopo la ricezione dell'ACK finale, in tal
  caso quello che si fa usualmente è lanciare un nuovo processo per leggere i
  successivi dati, che si bloccherà su una \func{read} se questi non sono
  disponibili; in questo modo si saranno impiegate delle risorse (per la
  creazione del nuovo processo) che non vengono usate immediatamente.  L'uso
  di \const{TCP\_DEFER\_ACCEPT} consente di intervenire anche in questa
  situazione; quando la si invoca sul lato server (vale a dire su un socket in
  ascolto) l'opzione fa sì che \func{accept} ritorni soltanto quando sono
  presenti dei dati sul socket, e non alla ricezione dell'ACK conclusivo del
  \textit{three way handshake}.

  L'opzione prende un valore intero che indica il numero massimo di secondi
  per cui mantenere il ritardo, sia per quanto riguarda il ritorno di
  \func{accept} su un server, che per l'invio dell'ACK finale insieme ai dati
  su un client. L'opzione è specifica di Linux non deve essere utilizzata in
  codice che vuole essere portabile.\footnote{su FreeBSD è presente una
    opzione \texttt{SO\_ACCEPTFILTER} che consente di ottenere lo stesso
    comportamento di \const{TCP\_DEFER\_ACCEPT} per quanto riguarda il lato
    server.}

\item[\constd{TCP\_WINDOW\_CLAMP}] con questa opzione si legge o si imposta
  alla dimensione specificata, in byte, il valore dichiarato della
  \textit{advertised window} (vedi sez.~\ref{sec:tcp_protocol_xxx}). Il kernel
  impone comunque una dimensione minima pari a \texttt{SOCK\_MIN\_RCVBUF/2}.
  Questa opzione non deve essere utilizzata in codice che vuole essere
  portabile.

\begin{figure}[!htb]
  \footnotesize \centering
  \begin{minipage}[c]{0.80\textwidth}
    \includestruct{listati/tcp_info.h}
  \end{minipage}
  \caption{La struttura \structd{tcp\_info} contenente le informazioni sul
    socket restituita dall'opzione \const{TCP\_INFO}.}
  \label{fig:tcp_info_struct}
\end{figure}

\item[\constd{TCP\_INFO}] questa opzione, specifica di Linux, ma introdotta
  anche in altri kernel (ad esempio FreeBSD) permette di controllare lo stato
  interno di un socket TCP direttamente da un programma in \textit{user space}.
  L'opzione restituisce in una speciale struttura \struct{tcp\_info}, la cui
  definizione è riportata in fig.~\ref{fig:tcp_info_struct}, tutta una serie
  di dati che il kernel mantiene, relativi al socket.  Anche questa opzione
  deve essere evitata se si vuole scrivere codice portabile.

  Con questa opzione diventa possibile ricevere una serie di informazioni
  relative ad un socket TCP così da poter effettuare dei controlli senza dover
  passare attraverso delle operazioni di lettura. Ad esempio si può verificare
  se un socket è stato chiuso usando una funzione analoga a quella illustrata
  in fig.~\ref{fig:is_closing}, in cui si utilizza il valore del campo
  \var{tcpi\_state} di \struct{tcp\_info} per controllare lo stato del socket.

\begin{figure}[!htbp]
  \footnotesize \centering
  \begin{minipage}[c]{\codesamplewidth}
    \includecodesample{listati/is_closing.c}
  \end{minipage}
  \caption{Codice della funzione \texttt{is\_closing.c}, che controlla lo stato
    di un socket TCP per verificare se si sta chiudendo.}
  \label{fig:is_closing}
\end{figure}

%Si noti come nell'esempio si sia (


\item[\constd{TCP\_QUICKACK}] con questa opzione è possibile eseguire una forma
  di controllo sull'invio dei segmenti ACK all'interno di in flusso di dati su
  TCP. In genere questo invio viene gestito direttamente dal kernel, il
  comportamento standard, corrispondente la valore logico di vero (in genere
  1) per questa opzione, è quello di inviare immediatamente i segmenti ACK, in
  quanto normalmente questo significa che si è ricevuto un blocco di dati e si
  può passare all'elaborazione del blocco successivo.

  Qualora però la nostra applicazione sappia in anticipo che alla ricezione di
  un blocco di dati seguirà immediatamente l'invio di un altro
  blocco,\footnote{caso tipico ad esempio delle risposte alle richieste HTTP.}
  poter accorpare quest'ultimo al segmento ACK permette di risparmiare sia in
  termini di dati inviati che di velocità di risposta. Per far questo si può
  utilizzare \const{TCP\_QUICKACK} impostando un valore logico falso (cioè 0),
  in questo modo il kernel attenderà così da inviare il prossimo segmento di
  ACK insieme ai primi dati disponibili.

  Si tenga presente che l'opzione non è permanente, vale a dire che una volta
  che la si sia impostata a 0 il kernel la riporterà al valore di default dopo
  il suo primo utilizzo. Sul lato server la si può impostare anche una volta
  sola su un socket in ascolto, ed essa verrà ereditata da tutti i socket che
  si otterranno da esso al ritorno di \func{accept}.

% TODO trattare con gli esempi di apache

\item[\constd{TCP\_CONGESTION}] questa opzione permette di impostare quale
  algoritmo per il controllo della congestione\footnote{il controllo della
    congestione è un meccanismo previsto dal protocollo TCP (vedi
    sez.~\ref{sec:tcp_protocol_xxx}) per evitare di trasmettere inutilmente
    dati quando una connessione è congestionata; un buon algoritmo è
    fondamentale per il funzionamento del protocollo, dato che i pacchetti
    persi andrebbero ritrasmessi, per cui inviare un pacchetto su una linea
    congestionata potrebbe causare facilmente un peggioramento della
    situazione.} utilizzare per il singolo socket. L'opzione è stata
  introdotta con il kernel 2.6.13,\footnote{alla data di stesura di queste
    note (Set. 2006) è pure scarsamente documentata, tanto che non è neanche
    definita nelle intestazioni della \acr{glibc} per cui occorre definirla a
    mano al suo valore che è 13.} e prende come per \param{optval} il
  puntatore ad un buffer contenente il nome dell'algoritmo di controllo che
  si vuole usare. 

  L'uso di un nome anziché di un valore numerico è dovuto al fatto che gli
  algoritmi di controllo della congestione sono realizzati attraverso
  altrettanti moduli del kernel, e possono pertanto essere attivati a
  richiesta; il nome consente di caricare il rispettivo modulo e di introdurre
  moduli aggiuntivi che implementino altri meccanismi. 

  Per poter disporre di questa funzionalità occorre aver compilato il kernel
  attivando l'opzione di configurazione generale
  \texttt{TCP\_CONG\_ADVANCED},\footnote{disponibile come \textit{TCP:
      advanced congestion control} nel menù \textit{Network->Networking
      options}, che a sua volta renderà disponibile un ulteriore menù con gli
    algoritmi presenti.} e poi abilitare i singoli moduli voluti con le varie
  \texttt{TCP\_CONG\_*} presenti per i vari algoritmi disponibili; un elenco
  di quelli attualmente supportati nella versione ufficiale del kernel è
  riportato in tab.~\ref{tab:sock_tcp_congestion_algo}.\footnote{la lista è
    presa dalla versione 2.6.17.}


  Si tenga presente che prima della implementazione modulare alcuni di questi
  algoritmi erano disponibili soltanto come caratteristiche generali del
  sistema, attivabili per tutti i socket, questo è ancora possibile con la
  \textit{sysctl} \texttt{tcp\_congestion\_control} (vedi
  sez.~\ref{sec:sock_ipv4_sysctl}) che ha sostituito le precedenti
  \textit{sysctl}.\footnote{riportate anche, alla data di stesura di queste
    pagine (Set. 2006) nelle pagine di manuale, ma non più presenti.}

  \begin{table}[!htb]
    \centering
    \footnotesize
    \begin{tabular}[c]{|l|l|l|}
      \hline
      \textbf{Nome}&\textbf{Configurazione}&\textbf{Riferimento} \\
      \hline
      \hline
      reno& -- &Algoritmo tradizionale, usato in caso di assenza degli altri.\\
      \texttt{bic}     &\texttt{TCP\_CONG\_BIC}     & 
      \url{http://www.csc.ncsu.edu/faculty/rhee/export/bitcp/index.htm}.\\
      \texttt{cubic}   &\texttt{TCP\_CONG\_CUBIC}     & 
      \url{http://www.csc.ncsu.edu/faculty/rhee/export/bitcp/index.htm}.\\
      \texttt{highspeed}&\texttt{TCP\_CONG\_HSTCP}  & 
      \url{http://www.icir.org/floyd/hstcp.html}.\\
      \texttt{htcp}    &\texttt{TCP\_CONG\_HTCP}    & 
      \url{http://www.hamilton.ie/net/htcp/}.\\
      \texttt{hybla}   &\texttt{TCP\_CONG\_HYBLA}   &       
      \url{http://www.danielinux.net/projects.html}.\\
      \texttt{scalable}&\texttt{TCP\_CONG\_SCALABLE}&  
      \url{http://www.deneholme.net/tom/scalable/}.\\
      \texttt{vegas}   &\texttt{TCP\_CONG\_VEGAS}   &  
      \url{http://www.cs.arizona.edu/protocols/}.\\
      \texttt{westwood}&\texttt{TCP\_CONG\_WESTWOOD}& 
      \url{http://www.cs.ucla.edu/NRL/hpi/tcpw/}.\\
%      \texttt{}&\texttt{}& .\\
      \hline
    \end{tabular}
    \caption{Gli algoritmi per il controllo della congestione disponibili con
      Linux con le relative opzioni di configurazione da attivare.}  
    \label{tab:sock_tcp_congestion_algo}
  \end{table}

\end{basedescript}


Il protocollo UDP, anche per la sua maggiore semplicità, supporta un numero
ridotto di opzioni, riportate in tab.~\ref{tab:sock_opt_udplevel}; anche in
questo caso per poterle utilizzare occorrerà impostare l'opportuno valore per
l'argomento \param{level}, che è \const{SOL\_UDP} (o l'equivalente
\constd{IPPROTO\_UDP}).  Le costanti che identificano dette opzioni sono
definite in \headfiled{netinet/udp.h}, ed accessibili includendo detto
file.\footnote{come per TCP, la definizione delle opzioni effettivamente
  supportate dal kernel si trova in realtà nel file
  \texttt{include/linux/udp.h}, dal quale si sono estratte le costanti di
  tab.~\ref{tab:sock_opt_udplevel}.}

\begin{table}[!htb]
  \centering
  \footnotesize
  \begin{tabular}[c]{|l|c|c|c|l|l|}
    \hline
    \textbf{Opzione}&\texttt{get}&\texttt{set}&\textbf{flag}&\textbf{Tipo}&
                    \textbf{Descrizione}\\
    \hline
    \hline
    \constd{UDP\_CORK}  &$\bullet$&$\bullet$&$\bullet$&\texttt{int}& %??? 
      Accumula tutti i dati su un unico pacchetto.\\
    \constd{UDP\_ENCAP} &$\bullet$&$\bullet$&$\bullet$&\texttt{int}& %??? 
      Non documentata.\\
   \hline
  \end{tabular}
  \caption{Le opzioni per i socket UDP disponibili al livello
    \const{SOL\_UDP}.}
  \label{tab:sock_opt_udplevel}
\end{table}

% TODO documentare \const{UDP\_ENCAP} 

Ancora una volta le descrizioni contenute tab.~\ref{tab:sock_opt_udplevel}
sono un semplice riferimento, una maggiore quantità di dettagli sulle
caratteristiche delle opzioni citate è quello dell'elenco seguente:
\begin{basedescript}{\desclabelwidth{1.5cm}\desclabelstyle{\nextlinelabel}}

\item[\constd{UDP\_CORK}] questa opzione ha l'identico effetto dell'analoga
  \const{TCP\_CORK} vista in precedenza per il protocollo TCP, e quando
  abilitata consente di accumulare i dati in uscita su un solo pacchetto che
  verrà inviato una volta che la si disabiliti. L'opzione è stata introdotta
  con il kernel 2.5.44, e non deve essere utilizzata in codice che vuole
  essere portabile.

\item[\constd{UDP\_ENCAP}] Questa opzione permette di gestire l'incapsulazione
  dei dati nel protocollo UDP. L'opzione è stata introdotta con il kernel
  2.5.67, e non è documentata. Come la precedente è specifica di Linux e non
  deve essere utilizzata in codice portabile.

\end{basedescript}




\section{La gestione attraverso le funzioni di controllo}
\label{sec:sock_ctrl_func}

Benché la maggior parte delle caratteristiche dei socket sia gestibile con le
funzioni \func{setsockopt} e \func{getsockopt}, alcune proprietà possono
essere impostate attraverso le funzioni \func{fcntl} e \func{ioctl} già
trattate in sez.~\ref{sec:file_fcntl_ioctl}; in quell'occasione abbiamo
parlato di queste funzioni esclusivamente nell'ambito della loro applicazione
a file descriptor associati a dei file normali; qui tratteremo invece i
dettagli del loro utilizzo con file descriptor associati a dei socket.


\subsection{L'uso di \func{ioctl} e \func{fcntl} per i socket generici}
\label{sec:sock_ioctl}

Tratteremo in questa sezione le caratteristiche specifiche delle funzioni
\func{ioctl} e \func{fcntl} quando esse vengono utilizzate con dei socket
generici. Quanto già detto in precedenza sez.~\ref{sec:file_fcntl_ioctl}
continua a valere; quello che tratteremo qui sono le operazioni ed i comandi
che sono validi, o che hanno significati peculiari, quando queste funzioni
vengono applicate a dei socket generici.

Nell'elenco seguente si riportano i valori specifici che può assumere il
secondo argomento della funzione \func{ioctl} (\param{request}, che indica il
tipo di operazione da effettuare) quando essa viene applicata ad un socket
generico. Nell'elenco si illustrerà anche, per ciascuna operazione, il tipo di
dato usato come terzo argomento della funzione ed il significato che esso
viene ad assumere.  Dato che in caso di lettura questi dati vengono restituiti
come \textit{value result argument}, con queste operazioni il terzo argomento
deve sempre essere passato come puntatore ad una variabile (o struttura)
precedentemente allocata. Le costanti che identificano le operazioni sono le
seguenti:
\begin{basedescript}{\desclabelwidth{1.5cm}\desclabelstyle{\nextlinelabel}}
\item[\constd{SIOCGSTAMP}] restituisce il contenuto di una struttura
  \struct{timeval} con la marca temporale dell'ultimo pacchetto ricevuto sul
  socket, questa operazione può essere utilizzata per effettuare delle
  misurazioni precise del tempo di andata e ritorno\footnote{il \textit{Round
      Trip Time} cui abbiamo già accennato in sez.~\ref{sec:net_tcp}.} dei
  pacchetti sulla rete.

\item[\constd{SIOCSPGRP}] imposta il processo o il \textit{process group} a
  cui inviare i segnali \signal{SIGIO} e \signal{SIGURG} quando viene
  completata una operazione di I/O asincrono o arrivano dei dati urgenti
  (\texttt{out-of-band}). Il terzo argomento deve essere un puntatore ad una
  variabile di tipo \type{pid\_t}; un valore positivo indica direttamente il
  \ids{PID} del processo, mentre un valore negativo indica (col valore
  assoluto) il \textit{process group}. Senza privilegi di amministratore o la
  \textit{capability} \const{CAP\_KILL} si può impostare solo se stessi o il
  proprio \textit{process group}.

\item[\constd{SIOCGPGRP}] legge le impostazioni presenti sul socket
  relativamente all'eventuale processo o \textit{process group} cui devono
  essere inviati i segnali \signal{SIGIO} e \signal{SIGURG}. Come per
  \const{SIOCSPGRP} l'argomento passato deve un puntatore ad una variabile di
  tipo \type{pid\_t}, con lo stesso significato.  Qualora non sia presente
  nessuna impostazione verrà restituito un valore nullo.

\item[\const{FIOASYNC}] Abilita o disabilita la modalità di I/O asincrono sul
  socket. Questo significa (vedi sez.~\ref{sec:signal_driven_io}) che verrà
  inviato il segnale di \signal{SIGIO} (o quanto impostato con
  \const{F\_SETSIG}, vedi sez.~\ref{sec:file_fcntl_ioctl}) in caso di eventi
  di I/O sul socket.
\end{basedescript}

Nel caso dei socket generici anche \func{fcntl} prevede un paio di comandi
specifici; in questo caso il secondo argomento (\param{cmd}, che indica il
comando) può assumere i due valori \const{FIOGETOWN} e \const{FIOSETOWN},
mentre il terzo argomento dovrà essere un puntatore ad una variabile di tipo
\type{pid\_t}. Questi due comandi sono una modalità alternativa di eseguire le
stesse operazioni (lettura o impostazione del processo o del gruppo di
processo che riceve i segnali) che si effettuano chiamando \func{ioctl} con
\const{SIOCGPGRP} e \const{SIOCSPGRP}.


\subsection{L'uso di \func{ioctl} per l'accesso ai dispositivi di rete}
\label{sec:sock_ioctl_netdevice}

Benché non strettamente attinenti alla gestione dei socket, vale la pena di
trattare qui l'interfaccia di accesso a basso livello ai dispositivi di rete
che viene appunto fornita attraverso la funzione \texttt{ioctl}. Questa non è
attinente a caratteristiche specifiche di un qualche protocollo, ma si applica
a tutti i socket, indipendentemente da tipo e famiglia degli stessi, e
permette di impostare e rilevare le funzionalità delle interfacce di rete.

\begin{figure}[!htb]
  \footnotesize \centering
  \begin{minipage}[c]{0.80\textwidth}
    \includestruct{listati/ifreq.h}
  \end{minipage}
  \caption{La struttura \structd{ifreq} utilizzata dalle \func{ioctl} per le
    operazioni di controllo sui dispositivi di rete.}
  \label{fig:netdevice_ifreq_struct}
\end{figure}

Tutte le operazioni di questo tipo utilizzano come terzo argomento di
\func{ioctl} il puntatore ad una struttura \struct{ifreq}, la cui definizione
è illustrata in fig.~\ref{fig:netdevice_ifreq_struct}. Normalmente si utilizza
il primo campo della struttura, \var{ifr\_name} per specificare il nome
dell'interfaccia su cui si vuole operare (ad esempio \texttt{eth0},
\texttt{ppp0}, ecc.), e si inseriscono (o ricevono) i valori relativi alle
diversa caratteristiche e funzionalità nel secondo campo, che come si può
notare è definito come una \dirct{union} proprio in quanto il suo significato
varia a secondo dell'operazione scelta.

Si tenga inoltre presente che alcune di queste operazioni (in particolare
quelle che modificano le caratteristiche dell'interfaccia) sono privilegiate e
richiedono i privilegi di amministratore o la \textit{capability}
\const{CAP\_NET\_ADMIN}, altrimenti si otterrà un errore di \errval{EPERM}.
Le costanti che identificano le operazioni disponibili sono le seguenti:
\begin{basedescript}{\desclabelwidth{1.5cm}\desclabelstyle{\nextlinelabel}}
\item[\constd{SIOCGIFNAME}] questa è l'unica operazione che usa il campo
  \var{ifr\_name} per restituire un risultato, tutte le altre lo utilizzano
  per indicare l'interfaccia sulla quale operare. L'operazione richiede che si
  indichi nel campo \var{ifr\_ifindex} il valore numerico dell'\textsl{indice}
  dell'interfaccia, e restituisce il relativo nome in \var{ifr\_name}.

  Il kernel infatti assegna ad ogni interfaccia un numero progressivo, detto
  appunto \itindex{interface~index} \textit{interface index}, che è quello che
  effettivamente la identifica nelle operazioni a basso livello, il nome
  dell'interfaccia è soltanto una etichetta associata a detto \textsl{indice},
  che permette di rendere più comprensibile l'indicazione dell'interfaccia
  all'interno dei comandi. Una modalità per ottenere questo valore è usare il
  comando \cmd{ip link}, che fornisce un elenco delle interfacce presenti
  ordinato in base a tale valore (riportato come primo campo).
  

\item[\constd{SIOCGIFINDEX}] restituisce nel campo \var{ifr\_ifindex} il valore
  numerico dell'indice dell'interfaccia specificata con \var{ifr\_name}, è in
  sostanza l'operazione inversa di \const{SIOCGIFNAME}.

\item[\constd{SIOCGIFFLAGS}] permette di ottenere nel campo \var{ifr\_flags} il
  valore corrente dei flag dell'interfaccia specificata (con \var{ifr\_name}).
  Il valore restituito è una maschera binaria i cui bit sono identificabili
  attraverso le varie costanti di tab.~\ref{tab:netdevice_iface_flag}.

\begin{table}[htb]
  \centering
  \footnotesize
  \begin{tabular}[c]{|l|p{8cm}|}
    \hline
    \textbf{Flag} & \textbf{Significato} \\
    \hline
    \hline
    \constd{IFF\_UP}        & L'interfaccia è attiva.\\
    \constd{IFF\_BROADCAST} & L'interfaccia ha impostato un indirizzo di
                             \textit{broadcast} valido.\\
    \constd{IFF\_DEBUG}     & È attivo il flag interno di debug.\\
    \constd{IFF\_LOOPBACK}  & L'interfaccia è una interfaccia di
                             \textit{loopback}.\\ 
    \constd{IFF\_POINTOPOINT}&L'interfaccia è associata ad un collegamento
                             \textsl{punto-punto}.\\ 
    \constd{IFF\_RUNNING}   & L'interfaccia ha delle risorse allocate (non può
                             quindi essere disattivata).\\
    \constd{IFF\_NOARP}     & L'interfaccia ha il protocollo ARP disabilitato o
                             l'indirizzo del livello di rete non è impostato.\\
    \constd{IFF\_PROMISC}   & L'interfaccia è nel cosiddetto
                             \index{modo~promiscuo} \textsl{modo promiscuo},
                             riceve cioè tutti i pacchetti che vede passare,
                             compresi quelli non direttamente indirizzati a
                             lei.\\
    \constd{IFF\_NOTRAILERS}& Evita l'uso di \textit{trailer} nei pacchetti.\\
    \constd{IFF\_ALLMULTI}  & Riceve tutti i pacchetti di \textit{multicast}.\\
    \constd{IFF\_MASTER}    & L'interfaccia è il master di un bundle per il
                             bilanciamento di carico.\\
    \constd{IFF\_SLAVE}     & L'interfaccia è uno slave di un bundle per il
                             bilanciamento di carico.\\
    \constd{IFF\_MULTICAST} & L'interfaccia ha il supporto per il
                             \textit{multicast} attivo.\\
    \constd{IFF\_PORTSEL}   & L'interfaccia può impostare i suoi parametri
                             hardware (con l'uso di \struct{ifmap}).\\
    \constd{IFF\_AUTOMEDIA} & L'interfaccia è in grado di selezionare
                             automaticamente il tipo di collegamento.\\
    \constd{IFF\_DYNAMIC}   & Gli indirizzi assegnati all'interfaccia vengono
                             persi quando questa viene disattivata.\\
%    \const{IFF\_}      & .\\
    \hline
  \end{tabular}
  \caption{Le costanti che identificano i vari bit della maschera binaria
    \var{ifr\_flags} che esprime i flag di una interfaccia di rete.}
  \label{tab:netdevice_iface_flag}
\end{table}


\item[\constd{SIOCSIFFLAGS}] permette di impostare il valore dei flag
  dell'interfaccia specificata (sempre con \var{ifr\_name}, non staremo a
  ripeterlo oltre) attraverso il valore della maschera binaria da passare nel
  campo \var{ifr\_flags}, che può essere ottenuta con l'OR aritmetico delle
  costanti di tab.~\ref{tab:netdevice_iface_flag}; questa operazione è
  privilegiata.

\item[\constd{SIOCGIFMETRIC}] permette di leggere il valore della metrica del
  dispositivo associato all'interfaccia specificata nel campo
  \var{ifr\_metric}.  Attualmente non è implementato, e l'operazione
  restituisce sempre un valore nullo.

\item[\constd{SIOCSIFMETRIC}] permette di impostare il valore della metrica del
  dispositivo al valore specificato nel campo \var{ifr\_metric}, attualmente
  non ancora implementato, restituisce un errore di \errval{EOPNOTSUPP}.

\item[\constd{SIOCGIFMTU}] permette di leggere il valore della \textit{Maximum
    Transfer Unit} del dispositivo nel campo \var{ifr\_mtu}.

\item[\constd{SIOCSIFMTU}] permette di impostare il valore della
  \textit{Maximum Transfer Unit} del dispositivo al valore specificato campo
  \var{ifr\_mtu}. L'operazione è privilegiata, e si tenga presente che
  impostare un valore troppo basso può causare un blocco del kernel.

\item[\constd{SIOCGIFHWADDR}] permette di leggere il valore dell'indirizzo
  hardware del dispositivo associato all'interfaccia nel campo
  \var{ifr\_hwaddr}; questo viene restituito come struttura \struct{sockaddr}
  in cui il campo \var{sa\_family} contiene un valore \texttt{ARPHRD\_*}
  indicante il tipo di indirizzo ed il campo \var{sa\_data} il valore binario
  dell'indirizzo hardware a partire dal byte 0.

\item[\constd{SIOCSIFHWADDR}] permette di impostare il valore dell'indirizzo
  hardware del dispositivo associato all'interfaccia attraverso il valore
  della struttura \struct{sockaddr} (con lo stesso formato illustrato per
  \const{SIOCGIFHWADDR}) passata nel campo \var{ifr\_hwaddr}. L'operazione è
  privilegiata.

\item[\constd{SIOCSIFHWBROADCAST}] imposta l'indirizzo \textit{broadcast}
  hardware dell'interfaccia al valore specificato dal campo
  \var{ifr\_hwaddr}. L'operazione è privilegiata.

\item[\constd{SIOCGIFMAP}] legge alcuni parametri hardware (memoria, interrupt,
  canali di DMA) del driver dell'interfaccia specificata, restituendo i
  relativi valori nel campo \var{ifr\_map}; quest'ultimo contiene una
  struttura di tipo \struct{ifmap}, la cui definizione è illustrata in
  fig.~\ref{fig:netdevice_ifmap_struct}.

\begin{figure}[!htb]
  \footnotesize \centering
  \begin{minipage}[c]{0.80\textwidth}
    \includestruct{listati/ifmap.h}
  \end{minipage}
  \caption{La struttura \structd{ifmap} utilizzata per leggere ed impostare i
    valori dei parametri hardware di un driver di una interfaccia.}
  \label{fig:netdevice_ifmap_struct}
\end{figure}

\item[\constd{SIOCSIFMAP}] imposta i parametri hardware del driver
  dell'interfaccia specificata, restituendo i relativi valori nel campo
  \var{ifr\_map}. Come per \const{SIOCGIFMAP} questo deve essere passato come
  struttura \struct{ifmap}, secondo la definizione di
  fig.~\ref{fig:netdevice_ifmap_struct}.

\item[\constd{SIOCADDMULTI}] aggiunge un indirizzo di \textit{multicast} ai
  filtri del livello di collegamento associati dell'interfaccia. Si deve usare
  un indirizzo hardware da specificare attraverso il campo \var{ifr\_hwaddr},
  che conterrà l'opportuna struttura \struct{sockaddr}; l'operazione è
  privilegiata. Per una modalità alternativa per eseguire la stessa operazione
  si possono usare i \textit{packet socket}, vedi
  sez.~\ref{sec:packet_socket}.

\item[\constd{SIOCDELMULTI}] rimuove un indirizzo di \textit{multicast} ai
  filtri del livello di collegamento dell'interfaccia, vuole un indirizzo
  hardware specificato come per \const{SIOCADDMULTI}. Anche questa operazione
  è privilegiata e può essere eseguita in forma alternativa con i
  \textit{packet socket}.

\item[\constd{SIOCGIFTXQLEN}] permette di leggere la lunghezza della coda di
  trasmissione del dispositivo associato all'interfaccia specificata nel campo
  \var{ifr\_qlen}.

\item[\constd{SIOCSIFTXQLEN}] permette di impostare il valore della lunghezza
  della coda di trasmissione del dispositivo associato all'interfaccia, questo
  deve essere specificato nel campo \var{ifr\_qlen}. L'operazione è
  privilegiata. 

\item[\constd{SIOCSIFNAME}] consente di cambiare il nome dell'interfaccia
  indicata da \var{ifr\_name} utilizzando il nuovo nome specificato nel campo
  \var{ifr\_rename}.

\end{basedescript}


% TODO aggiunta con il kernel 3.14 SIOCGHWTSTAMP per ottenere il timestamp
% hardware senza modificarlo

Una ulteriore operazione, che consente di ricavare le caratteristiche delle
interfacce di rete, è \constd{SIOCGIFCONF}; però per ragioni di compatibilità
questa operazione è disponibile soltanto per i socket della famiglia
\const{AF\_INET} (vale ad dire per socket IPv4). In questo caso l'utente dovrà
passare come argomento una struttura \struct{ifconf}, definita in
fig.~\ref{fig:netdevice_ifconf_struct}.

\begin{figure}[!htb]
  \footnotesize \centering
  \begin{minipage}[c]{0.80\textwidth}
    \includestruct{listati/ifconf.h}
  \end{minipage}
  \caption{La struttura \structd{ifconf}.}
  \label{fig:netdevice_ifconf_struct}
\end{figure}

Per eseguire questa operazione occorrerà allocare preventivamente un buffer di
contenente un vettore di strutture \struct{ifreq}. La dimensione (in byte) di
questo buffer deve essere specificata nel campo \var{ifc\_len} di
\struct{ifconf}, mentre il suo indirizzo andrà specificato nel campo
\var{ifc\_req}. Qualora il buffer sia stato allocato come una stringa, il suo
indirizzo potrà essere fornito usando il campo \var{ifc\_buf}.\footnote{si
  noti che l'indirizzo del buffer è definito in \struct{ifconf} con una
  \dirct{union}, questo consente di utilizzare una delle due forme a piacere.}

La funzione restituisce nel buffer indicato una serie di strutture
\struct{ifreq} contenenti nel campo \var{ifr\_name} il nome dell'interfaccia e
nel campo \var{ifr\_addr} il relativo indirizzo IP.  Se lo spazio allocato nel
buffer è sufficiente il kernel scriverà una struttura \struct{ifreq} per
ciascuna interfaccia attiva, restituendo nel campo \var{ifc\_len} il totale
dei byte effettivamente scritti. Il valore di ritorno è 0 se l'operazione ha
avuto successo e negativo in caso contrario. 

Si tenga presente che il kernel non scriverà mai sul buffer di uscita dati
eccedenti numero di byte specificato col valore di \var{ifc\_len} impostato
alla chiamata della funzione, troncando il risultato se questi non dovessero
essere sufficienti. Questa condizione non viene segnalata come errore per cui
occorre controllare il valore di \var{ifc\_len} all'uscita della funzione, e
verificare che esso sia inferiore a quello di ingresso. In caso contrario si è
probabilmente\footnote{probabilmente perché si potrebbe essere nella
  condizione in cui sono stati usati esattamente quel numero di byte.} avuta
una situazione di troncamento dei dati.

\begin{figure}[!htbp]
  \footnotesize \centering
  \begin{minipage}[c]{\codesamplewidth}
    \includecodesample{listati/iflist.c}
  \end{minipage}
  \caption{Il corpo principale del programma \texttt{iflist.c}.}
  \label{fig:netdevice_iflist}
\end{figure}

Come esempio dell'uso di queste funzioni si è riportato in
fig.~\ref{fig:netdevice_iflist} il corpo principale del programma
\texttt{iflist} in cui si utilizza l'operazione \const{SIOCGIFCONF} per
ottenere una lista delle interfacce attive e dei relativi indirizzi. Al solito
il codice completo è fornito nei sorgenti allegati alla guida.

Il programma inizia (\texttt{\small 7--11}) con la creazione del socket
necessario ad eseguire l'operazione, dopo di che si inizializzano
opportunamente (\texttt{\small 13--14}) i valori della struttura
\struct{ifconf} indicando la dimensione del buffer ed il suo
indirizzo;\footnote{si noti come in questo caso si sia specificato l'indirizzo
  usando il campo \var{ifc\_buf}, mentre nel seguito del programma si accederà
  ai valori contenuti nel buffer usando \var{ifc\_req}.} si esegue poi
l'operazione invocando \func{ioctl}, controllando come sempre la corretta
esecuzione, ed uscendo in caso di errore (\texttt{\small 15--19}).

Si esegue poi un controllo sulla quantità di dati restituiti segnalando un
eventuale overflow del buffer (\texttt{\small 21--23}); se invece è tutto a
posto (\texttt{\small 24--27}) si calcola e si stampa a video il numero di
interfacce attive trovate.  L'ultima parte del programma (\texttt{\small
  28--33}) è il ciclo sul contenuto delle varie strutture \struct{ifreq}
restituite in cui si estrae (\texttt{\small 30}) l'indirizzo ad esse
assegnato\footnote{si è definito \var{access} come puntatore ad una struttura
  di tipo \struct{sockaddr\_in} per poter eseguire un \textit{casting}
  dell'indirizzo del valore restituito nei vari campi \var{ifr\_addr}, così
  poi da poterlo poi usare come argomento di \func{inet\_ntoa}.}  e lo si
stampa (\texttt{\small 31--32}) insieme al nome dell'interfaccia.



\subsection{L'uso di \func{ioctl} per i socket TCP e UDP}
\label{sec:sock_ioctl_IP}

Non esistono operazioni specifiche per i socket IP in quanto tali,\footnote{a
  parte forse \const{SIOCGIFCONF}, che però resta attinente alle proprietà
  delle interfacce di rete, per cui l'abbiamo trattata in
  sez.~\ref{sec:sock_ioctl_netdevice} insieme alle altre che comunque si
  applicano anche ai socket IP.} mentre per i pacchetti di altri protocolli
trasportati su IP, qualora li si gestisca attraverso dei socket, si dovrà fare
riferimento direttamente all'eventuale supporto presente per il tipo di socket
usato: ad esempio si possono ricevere pacchetti ICMP con socket di tipo
\texttt{raw}, nel qual caso si dovrà fare riferimento alle operazioni di
quest'ultimo.

Tuttavia la gran parte dei socket utilizzati nella programmazione di rete
utilizza proprio il protocollo IP, e quello che succede è che in realtà la
funzione \func{ioctl} consente di effettuare alcune operazioni specifiche per
i socket che usano questo protocollo, ma queste vendono eseguite, invece che a
livello di IP, al successivo livello di trasporto, vale a dire in maniera
specifica per i socket TCP e UDP.

Le operazioni di controllo disponibili per i socket TCP sono illustrate dalla
relativa pagina di manuale, accessibile con \texttt{man 7 tcp}, e prevedono
come possibile valore per il secondo argomento della funzione le costanti
illustrate nell'elenco seguente; il terzo argomento della funzione, gestito
come \textit{value result argument}, deve essere sempre il puntatore ad una
variabile di tipo \ctyp{int}:
\begin{basedescript}{\desclabelwidth{1.5cm}\desclabelstyle{\nextlinelabel}}
\item[\constd{SIOCINQ}] restituisce la quantità di dati non ancora letti
  presenti nel buffer di ricezione; il socket non deve essere in stato
  \texttt{LISTEN}, altrimenti si avrà un errore di \errval{EINVAL}.
\item[\constd{SIOCATMARK}] ritorna un intero non nullo, da intendere come
  valore logico, se il flusso di dati letti sul socket è arrivato sulla
  posizione (detta anche \textit{urgent mark}) in cui sono stati ricevuti dati
  urgenti (vedi sez.~\ref{sec:TCP_urgent_data}).  Una operazione di lettura da
  un socket non attraversa mai questa posizione, per cui è possibile
  controllare se la si è raggiunta o meno con questa operazione.

  Questo è utile quando si attiva l'opzione \const{SO\_OOBINLINE} (vedi
  sez.~\ref{sec:sock_generic_options}) per ricevere i dati urgenti all'interno
  del flusso dei dati ordinari del socket;\footnote{vedremo in
    sez.~\ref{sec:TCP_urgent_data} che in genere i dati urgenti presenti su un
    socket si leggono \textit{out-of-band} usando un opportuno flag per
    \func{recvmsg}.}  in tal caso quando \const{SIOCATMARK} restituisce un
  valore non nullo si saprà che la successiva lettura dal socket restituirà i
  dati urgenti e non il normale traffico; torneremo su questo in maggior
  dettaglio in sez.~\ref{sec:TCP_urgent_data}.

\item[\constd{SIOCOUTQ}] restituisce la quantità di dati non ancora inviati
  presenti nel buffer di spedizione; come per \const{SIOCINQ} il socket non
  deve essere in stato \texttt{LISTEN}, altrimenti si avrà un errore di
  \errval{EINVAL}.
\end{basedescript}

Le operazioni di controllo disponibili per i socket UDP, anch'esse illustrate
dalla relativa pagina di manuale accessibile con \texttt{man 7 udp}, sono
quelle indicate dalle costanti del seguente elenco; come per i socket TCP il
terzo argomento viene gestito come \textit{value result argument} e deve
essere un puntatore ad una variabile di tipo \ctyp{int}:
\begin{basedescript}{\desclabelwidth{1.5cm}\desclabelstyle{\nextlinelabel}}
\item[\const{FIONREAD}] restituisce la dimensione in byte del primo pacchetto
  in attesa di ricezione, o 0 qualora non ci sia nessun pacchetto.
\item[\const{TIOCOUTQ}] restituisce il numero di byte presenti nella coda di
  invio locale; questa opzione è supportata soltanto a partire dal kernel 2.4
\end{basedescript}



\section{La gestione con \func{sysctl} ed il filesystem \texttt{/proc}}
\label{sec:sock_sysctl_proc}

Come ultimo argomento di questo capitolo tratteremo l'uso della funzione
\func{sysctl} (che è stata introdotta nelle sue funzionalità generiche in
sez.~\ref{sec:sys_sysctl}) per quanto riguarda le sue capacità di effettuare
impostazioni relative alle proprietà dei socket.  Dato che le stesse
funzionalità sono controllabili direttamente attraverso il filesystem
\texttt{/proc}, le tratteremo attraverso i file presenti in quest'ultimo.


\subsection{L'uso di \func{sysctl} e \texttt{/proc} per le proprietà della
  rete}
\label{sec:sock_sysctl}

La differenza nell'uso di \func{sysctl} e del filesystem \texttt{/proc}
rispetto a quello delle funzioni \func{ioctl} e \func{fcntl} visto in
sez.~\ref{sec:sock_ctrl_func} o all'uso di \func{getsockopt} e
\func{setsockopt} è che queste funzioni consentono di controllare le proprietà
di un singolo socket, mentre con \func{sysctl} e con \texttt{/proc} si
impostano proprietà (o valori di default) validi a livello dell'intero
sistema, e cioè per tutti i socket.

Le opzioni disponibili per le proprietà della rete, nella gerarchia dei valori
impostabili con \func{sysctl}, sono riportate sotto il nodo \texttt{net}, o,
se acceduti tramite l'interfaccia del filesystem \texttt{/proc}, sotto
\texttt{/proc/sys/net}. In genere sotto questa directory compaiono le
sottodirectory (corrispondenti ad altrettanti sotto-nodi per \func{sysctl})
relative ai vari protocolli e tipi di interfacce su cui è possibile
intervenire per effettuare impostazioni; un contenuto tipico di questa
directory è il seguente:
\begin{verbatim}
/proc/sys/net/
|-- core
|-- ethernet
|-- ipv4
|-- ipv6
|-- irda
|-- token-ring
`-- unix
\end{verbatim}
e sono presenti varie centinaia di parametri, molti dei quali non sono neanche
documentati; nel nostro caso ci limiteremo ad illustrare quelli più
significativi.

Si tenga presente infine che se è sempre possibile utilizzare il filesystem
\texttt{/proc} come sostituto di \func{sysctl}, dato che i valori di nodi e
sotto-nodi di quest'ultima sono mappati come file e directory sotto
\texttt{/proc/sys/}, non è vero il contrario, ed in particolare Linux consente
di impostare alcuni parametri o leggere lo stato della rete a livello di
sistema sotto \texttt{/proc/net}, dove sono presenti dei file che non
corrispondono a nessun nodo di \func{sysctl}.


\subsection{I valori di controllo per i socket generici}
\label{sec:sock_gen_sysctl}

Nella directory \texttt{/proc/sys/net/core/} sono presenti i file
corrispondenti ai parametri generici di \textit{sysctl} validi per tutti i
socket.  Quelli descritti anche nella pagina di manuale, accessibile con
\texttt{man 7 socket} sono i seguenti:

\begin{basedescript}{\desclabelwidth{2.2cm}\desclabelstyle{\nextlinelabel}}
\item[\sysctlrelfiled{net/core}{rmem\_default}] imposta la dimensione
  di default del buffer di ricezione (cioè per i dati in ingresso) dei socket.
\item[\sysctlrelfiled{net/core}{rmem\_max}] imposta la dimensione
  massima che si può assegnare al buffer di ricezione dei socket attraverso
  l'uso dell'opzione \const{SO\_RCVBUF}.
\item[\sysctlrelfiled{net/core}{wmem\_default}] imposta la dimensione
  di default del buffer di trasmissione (cioè per i dati in uscita) dei
  socket.
\item[\sysctlrelfiled{net/core}{wmem\_max}] imposta la dimensione
  massima che si può assegnare al buffer di trasmissione dei socket attraverso
  l'uso dell'opzione \const{SO\_SNDBUF}.
\item[\sysctlrelfiled{net/core}{message\_cost},
  \sysctlrelfiled{net/core}{message\_burst}] contengono le impostazioni del
  \textit{bucket filter} che controlla l'emissione di
  messaggi di avviso da parte del kernel per eventi relativi a problemi sulla
  rete, imponendo un limite che consente di prevenire eventuali attacchi di
  \textit{Denial of Service} usando i log.\footnote{senza questo limite un
    attaccante potrebbe inviare ad arte un traffico che generi
    intenzionalmente messaggi di errore, per saturare il sistema dei log.}

  \itindbeg{bucket~filter} 

  Il \textit{bucket filter} è un algoritmo generico che permette di impostare
  dei limiti di flusso su una quantità\footnote{uno analogo viene usato per
    imporre dei limiti sul flusso dei pacchetti nel \textit{netfilter} di
    Linux.}  senza dovere eseguire medie temporali, che verrebbero a dipendere
  in misura non controllabile dalla dimensione dell'intervallo su cui si media
  e dalla distribuzione degli eventi;\footnote{in caso di un picco di flusso
    (il cosiddetto \textit{burst}) il flusso medio verrebbe a dipendere in
    maniera esclusiva dalla dimensione dell'intervallo di tempo su cui calcola
    la media.} in questo caso si definisce la dimensione di un
  ``\textsl{bidone}'' (il \textit{bucket}) e del flusso che da esso può
  uscire, la presenza di una dimensione iniziale consente di assorbire
  eventuali picchi di emissione, l'aver fissato un flusso di uscita garantisce
  che a regime questo sarà il valore medio del flusso ottenibile dal
  \textit{bucket}.

  \itindend{bucket~filter} 

  I due valori indicano rispettivamente il flusso a regime (non sarà inviato
  più di un messaggio per il numero di secondi specificato da
  \texttt{message\_cost}) e la dimensione iniziale per in caso di picco di
  emissione (verranno accettati inizialmente fino ad un massimo di
  \texttt{message\_cost/message\_burst} messaggi).

\item[\sysctlrelfiled{net/core}{netdev\_max\_backlog}] numero massimo
  di pacchetti che possono essere contenuti nella coda di ingresso generale.

\item[\sysctlrelfiled{net/core}{optmem\_max}] lunghezza massima dei
  dati ancillari e di controllo (vedi sez.~\ref{sec:net_ancillary_data}).
\end{basedescript}

Oltre a questi, nella directory \texttt{/proc/sys/net/core} si trovano diversi
altri file, la cui documentazione, come per gli altri, dovrebbe essere
mantenuta nei sorgenti del kernel nel file
\texttt{Documentation/networking/ip-sysctl.txt}; la maggior parte di questi
però non è documentato:
\begin{basedescript}{\desclabelwidth{2.2cm}\desclabelstyle{\nextlinelabel}}
\item[\sysctlrelfiled{net/core}{dev\_weight}] blocco di lavoro (\textit{work
    quantum}) dello \textit{scheduler} di processo dei pacchetti.

% TODO da documentare meglio

\item[\sysctlrelfiled{net/core}{lo\_cong}] valore per l'occupazione
  della coda di ricezione sotto la quale si considera di avere una bassa
  congestione.

\item[\sysctlrelfiled{net/core}{mod\_cong}] valore per l'occupazione della
  coda di ricezione sotto la quale si considera di avere una congestione
  moderata.

\item[\sysctlrelfiled{net/core}{no\_cong}] valore per l'occupazione della coda
  di ricezione sotto la quale si considera di non avere congestione.

\item[\sysctlrelfiled{net/core}{no\_cong\_thresh}] valore minimo (\textit{low
    water mark}) per il riavvio dei dispositivi congestionati.

  % \item[\sysctlrelfiled{net/core}{netdev\_fastroute}] è presente
  %   soltanto quando si è compilato il kernel con l'apposita opzione di
  %   ottimizzazione per l'uso come router.

\item[\sysctlrelfiled{net/core}{somaxconn}] imposta la dimensione massima
  utilizzabile per il \textit{backlog} della funzione \func{listen} (vedi
  sez.~\ref{sec:TCP_func_listen}), e corrisponde al valore della costante
  \constd{SOMAXCONN}; il suo valore di default è 128.

\end{basedescript}


\subsection{I valori di controllo per il protocollo IPv4}
\label{sec:sock_ipv4_sysctl}

Nella directory \texttt{/proc/sys/net/ipv4} sono presenti i file che
corrispondono ai parametri dei socket che usano il protocollo IPv4, relativi
quindi sia alle caratteristiche di IP, che a quelle degli altri protocolli che
vengono usati all'interno di quest'ultimo (come ICMP, TCP e UDP) o a fianco
dello stesso (come ARP).

I file che consentono di controllare le caratteristiche specifiche del
protocollo IP in quanto tale, che sono descritti anche nella relativa pagina
di manuale accessibile con \texttt{man 7 ip}, sono i seguenti:
\begin{basedescript}{\desclabelwidth{2.2cm}\desclabelstyle{\nextlinelabel}}

\item[\sysctlrelfiled{net/ipv4}{ip\_default\_ttl}] imposta il valore di
  default per il campo TTL (vedi sez.~\ref{sec:IP_header}) di tutti i
  pacchetti uscenti, stabilendo così il numero massimo di router che i
  pacchetti possono attraversare. Il valore può essere modificato anche per il
  singolo socket con l'opzione \const{IP\_TTL}.  Prende un valore intero, ma
  dato che il campo citato è di 8 bit hanno senso solo valori fra 0 e 255. Il
  valore di default è 64, e normalmente non c'è nessuna necessità di
  modificarlo.\footnote{l'unico motivo sarebbe per raggiungere macchine
    estremamente ``{lontane}'' in termini di \textit{hop}, ma è praticamente
    impossibile trovarne.} Aumentare il valore è una pratica poco gentile, in
  quanto in caso di problemi di routing si allunga inutilmente il numero di
  ritrasmissioni.

\item[\sysctlrelfiled{net/ipv4}{ip\_forward}] abilita l'inoltro dei
  pacchetti da una interfaccia ad un altra, e può essere impostato anche per
  la singola interfaccia. Prende un valore logico (0 disabilita, diverso da
  zero abilita), di default è disabilitato.

\item[\sysctlrelfiled{net/ipv4}{ip\_dynaddr}] abilita la riscrittura
  automatica degli indirizzi associati ad un socket quando una interfaccia
  cambia indirizzo. Viene usato per le interfacce usate nei collegamenti in
  dial-up, il cui indirizzo IP viene assegnato dinamicamente dal provider, e
  può essere modificato. Prende un valore intero, con 0 si disabilita la
  funzionalità, con 1 la si abilita, con 2 (o con qualunque altro valore
  diverso dai precedenti) la si abilità in modalità \textsl{prolissa}; di
  default la funzionalità è disabilitata.

\item[\sysctlrelfiled{net/ipv4}{ip\_autoconfig}] specifica se
  l'indirizzo IP è stato configurato automaticamente dal kernel all'avvio
  attraverso DHCP, BOOTP o RARP. Riporta un valore logico (0 falso, 1 vero)
  accessibile solo in lettura, è inutilizzato nei kernel recenti ed eliminato
  a partire dal kernel 2.6.18.

\item[\sysctlrelfiled{net/ipv4}{ip\_local\_port\_range}] imposta l'intervallo
  dei valori usati per l'assegnazione delle porte effimere, permette cioè di
  modificare i valori illustrati in fig.~\ref{fig:TCP_port_alloc}; prende due
  valori interi separati da spazi, che indicano gli estremi
  dell'intervallo. Si abbia cura di non definire un intervallo che si
  sovrappone a quello delle porte usate per il \textit{masquerading}, il
  kernel può gestire la sovrapposizione, ma si avrà una perdita di
  prestazioni. Si imposti sempre un valore iniziale maggiore di 1024 (o meglio
  ancora di 4096) per evitare conflitti con le porte usate dai servizi noti.

\item[\sysctlrelfiled{net/ipv4}{ip\_no\_pmtu\_disc}] permette di disabilitare
  per i socket \const{SOCK\_STREAM} la ricerca automatica della \textit{Path
    MTU} (vedi sez.~\ref{sec:net_lim_dim} e
  sez.~\ref{sec:sock_ipv4_options}). Prende un valore logico, e di default è
  disabilitato (cioè la ricerca viene eseguita).

  In genere si abilita questo parametro quando per qualche motivo il
  procedimento del \textit{Path MTU discovery} fallisce; dato che questo può
  avvenire a causa di router\footnote{ad esempio se si scartano tutti i
    pacchetti ICMP, il problema è affrontato anche in sez.~3.4.4 di
    \cite{SGL}.} o interfacce\footnote{ad esempio se i due capi di un
    collegamento \textit{point-to-point} non si accordano sulla stessa MTU.}
  mal configurati è opportuno correggere le configurazioni, perché
  disabilitare globalmente il procedimento con questo parametro ha pesanti
  ripercussioni in termini di prestazioni di rete.

\item[\sysctlrelfiled{net/ipv4}{ip\_always\_defrag}] fa sì che tutti i
  pacchetti IP frammentati siano riassemblati, anche in caso in successivo
  immediato inoltro.\footnote{introdotto con il kernel 2.2.13, nelle versioni
    precedenti questo comportamento poteva essere solo stabilito un volta per
    tutte in fase di compilazione del kernel con l'opzione
    \texttt{CONFIG\_IP\_ALWAYS\_DEFRAG}.} Prende un valore logico e di default
  è disabilitato. Con i kernel dalla serie 2.4 in poi la deframmentazione
  viene attivata automaticamente quando si utilizza il sistema del
  \textit{netfilter}, e questo parametro non è più presente.

\item[\sysctlrelfiled{net/ipv4}{ipfrag\_high\_thresh}] indica il limite
  massimo (espresso in numero di byte) sui pacchetti IP frammentati presenti
  in coda; quando questo valore viene raggiunto la coda viene ripulita fino al
  valore \texttt{ipfrag\_low\_thresh}. Prende un valore intero.

\item[\sysctlrelfiled{net/ipv4}{ipfrag\_low\_thresh}] indica la dimensione
  (specificata in byte) della soglia inferiore a cui viene riportata la coda
  dei pacchetti IP frammentati quando si raggiunge il valore massimo dato da
  \texttt{ipfrag\_high\_thresh}. Prende un valore intero.

\item[\sysctlrelfiled{net/ipv4}{ip\_nonlocal\_bind}] se abilitato rende
  possibile ad una applicazione eseguire \func{bind} anche su un indirizzo che
  non è presente su nessuna interfaccia locale. Prende un valore logico e di
  default è disabilitato. La funzionalità può essere abilitata a livello di
  singolo socket con l'opzione \const{IP\_FREEBIND} illustrata in
  sez.~\ref{sec:sock_ipv4_options}. 

  Questo può risultare utile per applicazioni particolari (come gli
  \textit{sniffer}) che hanno la necessità di ricevere pacchetti anche non
  diretti agli indirizzi presenti sulla macchina, ad esempio per intercettare
  il traffico per uno specifico indirizzo che si vuole tenere sotto
  controllo. Il suo uso però può creare problemi ad alcune applicazioni.

% \item[\texttt{neigh/*}] La directory contiene i valori 
% TODO trattare neigh/* nella parte su arp, da capire dove sarà.

\end{basedescript}

I file di \texttt{/proc/sys/net/ipv4} che invece fanno riferimento alle
caratteristiche specifiche del protocollo TCP, elencati anche nella rispettiva
pagina di manuale (accessibile con \texttt{man 7 tcp}), sono i seguenti:
\begin{basedescript}{\desclabelwidth{2.2cm}\desclabelstyle{\nextlinelabel}}

\item[\sysctlrelfiled{net/ipv4}{tcp\_abort\_on\_overflow}] indica al
  kernel di azzerare le connessioni quando il programma che le riceve è troppo
  lento ed incapace di accettarle. Prende un valore logico ed è disabilitato
  di default.  Questo consente di recuperare le connessioni se si è avuto un
  eccesso dovuto ad un qualche picco di traffico, ma ovviamente va a discapito
  dei client che interrogano il server. Pertanto è da abilitare soltanto
  quando si è sicuri che non è possibile ottimizzare il server in modo che sia
  in grado di accettare connessioni più rapidamente.

\item[\sysctlrelfiled{net/ipv4}{tcp\_adv\_win\_scale}] indica al kernel quale
  frazione del buffer associato ad un socket\footnote{quello impostato con
    \sysctlrelfile{net/ipv4}{tcp\_rmem}.} deve essere utilizzata per la
  finestra del protocollo TCP\footnote{in sostanza il valore che costituisce
    la \textit{advertised window} annunciata all'altro capo del socket.} e
  quale come buffer applicativo per isolare la rete dalle latenze
  dell'applicazione.  Prende un valore intero che determina la suddetta
  frazione secondo la formula
  $\texttt{buffer}/2^\texttt{tcp\_adv\_win\_scale}$ se positivo o con
  $\texttt{buffer}-\texttt{buffer}/2^\texttt{tcp\_adv\_win\_scale}$ se
  negativo.  Il default è 2 che significa che al buffer dell'applicazione
  viene riservato un quarto del totale.

\item[\sysctlrelfiled{net/ipv4}{tcp\_app\_win}] indica la frazione della
  finestra TCP che viene riservata per gestire l'overhaed dovuto alla
  bufferizzazione. Prende un valore intero che consente di calcolare la
  dimensione in byte come il massimo fra la MSS e
  $\texttt{window}/2^\texttt{tcp\_app\_win}$. Un valore nullo significa che
  non viene riservato nessuno spazio; il valore di default è 31.

% vecchi, presumibilmente usati quando gli algoritmi di congestione non erano
% modularizzabili 
% \item[\texttt{tcp\_bic}]
% \item[\texttt{tcp\_bic\_low\_window}] 
% \item[\texttt{tcp\_bic\_fast\_convergence}] 

\item[\sysctlrelfiled{net/ipv4}{tcp\_dsack}] abilita il supporto,
  definito nell'\href{http://www.ietf.org/rfc/rfc2884.txt}{RFC~2884}, per il
  cosiddetto \textit{Duplicate SACK}.\footnote{si indica con SACK
    (\textit{Selective Acknowledgement}) un'opzione TCP, definita
    nell'\href{http://www.ietf.org/rfc/rfc2018.txt}{RFC~2018}, usata per dare
    un \textit{acknowledgement} unico su blocchi di pacchetti non contigui,
    che consente di diminuire il numero di pacchetti scambiati.} Prende un
  valore logico e di default è abilitato.
% TODO documentare o descrivere che cos'è il Duplicate SACK o
% mettere riferimento nelle appendici


\item[\sysctlrelfiled{net/ipv4}{tcp\_ecn}] abilita il meccanismo della
  \textit{Explicit Congestion Notification} (in breve ECN) nelle connessioni
  TCP. Prende valore logico che di default è disabilitato. La \textit{Explicit
    Congestion Notification} \itindex{Explicit~Congestion~Notification~(ECN)}
  è un meccanismo che consente di notificare quando una rotta o una rete è
  congestionata da un eccesso di traffico,\footnote{il meccanismo è descritto
    in dettaglio nell'\href{http://www.ietf.org/rfc/rfc3168.txt}{RFC~3168}
    mentre gli effetti sulle prestazioni del suo utilizzo sono documentate
    nell'\href{http://www.ietf.org/rfc/rfc2884.txt}{RFC~2884}.} si può così
  essere avvisati e cercare rotte alternative oppure diminuire l'emissione di
  pacchetti (in modo da non aumentare la congestione).

  Si tenga presente che se si abilita questa opzione si possono avere dei
  malfunzionamenti apparentemente casuali dipendenti dalla destinazione,
  dovuti al fatto che alcuni vecchi router non supportano il meccanismo ed
  alla sua attivazione scartano i relativi pacchetti, bloccando completamente
  il traffico.
% TODO documentare o descrivere che cos'è l'Explicit Congestion Notification o
% mettere riferimento nelle appendici


\item[\sysctlrelfiled{net/ipv4}{tcp\_fack}] abilita il supporto per il
  \textit{TCP Forward Acknowledgement}, un algoritmo per il controllo della
  congestione del traffico. Prende un valore logico e di default è abilitato.

% TODO documentare o descrivere che cos'è il TCP Forward Acknowledgement o
% mettere riferimento nelle appendici

\item[\sysctlrelfiled{net/ipv4}{tcp\_fin\_timeout}] specifica il numero di
  secondi da passare in stato \texttt{FIN\_WAIT2} nell'attesa delle ricezione
  del pacchetto FIN conclusivo, passati quali il socket viene comunque chiuso
  forzatamente.  Prende un valore intero che indica i secondi e di default è
  60.\footnote{nei kernel della serie 2.2.x era il valore utilizzato era
    invece di 120 secondi.} L'uso di questa opzione realizza quella che in
  sostanza è una violazione delle specifiche del protocollo TCP, ma è utile
  per fronteggiare alcuni attacchi di \textit{Denial of Service}.

\item[\sysctlrelfiled{net/ipv4}{tcp\_frto}] abilita il supporto per
  l'algoritmo F-RTO, un algoritmo usato per la ritrasmissione dei timeout del
  protocollo TCP, che diventa molto utile per le reti wireless dove la perdita
  di pacchetti è usualmente dovuta a delle interferenze radio, piuttosto che
  alla congestione dei router. Prende un valore logico e di default è
  disabilitato.

\item[\sysctlrelfiled{net/ipv4}{tcp\_keepalive\_intvl}] indica il
  numero di secondi che deve trascorrere fra l'emissione di due successivi
  pacchetti di test quando è abilitata la funzionalità del \textit{keepalive}
  (vedi sez.~\ref{sec:sock_options_main}). Prende un valore intero che di
  default è 75.

\item[\sysctlrelfiled{net/ipv4}{tcp\_keepalive\_probes}] indica il
  massimo numero pacchetti di \textit{keepalive} (vedi
  sez.~\ref{sec:sock_options_main}) che devono essere inviati senza ricevere
  risposta prima che il kernel decida che la connessione è caduta e la
  termini. Prende un valore intero che di default è 9.

\item[\sysctlrelfiled{net/ipv4}{tcp\_keepalive\_time}] indica il numero di
  secondi che devono passare senza traffico sulla connessione prima che il
  kernel inizi ad inviare pacchetti di \textit{keepalive}.\footnote{ha effetto
    solo per i socket per cui si è impostata l'opzione \const{SO\_KEEPALIVE}
    (vedi sez.~\ref{sec:sock_options_main}.}  Prende un valore intero che di
  default è 7200, pari a due ore.

\item[\sysctlrelfiled{net/ipv4}{tcp\_low\_latency}] indica allo stack
  TCP del kernel di ottimizzare il comportamento per ottenere tempi di latenza
  più bassi a scapito di valori più alti per l'utilizzo della banda. Prende un
  valore logico che di default è disabilitato in quanto un maggior utilizzo
  della banda è preferito, ma esistono applicazioni particolari in cui la
  riduzione della latenza è più importante (ad esempio per i cluster di
  calcolo parallelo) nelle quali lo si può abilitare.

\item[\sysctlrelfiled{net/ipv4}{tcp\_max\_orphans}] indica il numero
  massimo di socket TCP ``\textsl{orfani}'' (vale a dire non associati a
  nessun file descriptor) consentito nel sistema.\footnote{trattasi in genere
    delle connessioni relative a socket chiusi che non hanno completato il
    processo di chiusura.}  Quando il limite viene ecceduto la connessione
  orfana viene resettata e viene stampato un avvertimento. Questo limite viene
  usato per contrastare alcuni elementari attacchi di \textit{denial of
    service}.  Diminuire il valore non è mai raccomandato, in certe condizioni
  di rete può essere opportuno aumentarlo, ma si deve tenere conto del fatto
  che ciascuna connessione orfana può consumare fino a 64K di memoria del
  kernel. Prende un valore intero, il valore di default viene impostato
  inizialmente al valore del parametro del kernel \texttt{NR\_FILE}, e viene
  aggiustato a seconda della memoria disponibile.

% TODO verificare la spiegazione di connessione orfana.

\item[\sysctlrelfiled{net/ipv4}{tcp\_max\_syn\_backlog}] indica la lunghezza
  della coda delle connessioni incomplete, cioè delle connessioni per le quali
  si è ricevuto un SYN di richiesta ma non l'ACK finale del \textit{three way
    handshake} (si riveda quanto illustrato in
  sez.~\ref{sec:TCP_func_listen}).

  Quando questo valore è superato il kernel scarterà immediatamente ogni
  ulteriore richiesta di connessione. Prende un valore intero; il default, che
  è 256, viene automaticamente portato a 1024 qualora nel sistema ci sia
  sufficiente memoria (se maggiore di 128Mb) e ridotto a 128 qualora la
  memoria sia poca (inferiore a 32Mb).\footnote{si raccomanda, qualora si
    voglia aumentare il valore oltre 1024, di seguire la procedura citata
    nella pagina di manuale di TCP, e modificare il valore della costante
    \texttt{TCP\_SYNQ\_HSIZE} nel file \texttt{include/net/tcp.h} dei sorgenti
    del kernel, in modo che sia $\mathtt{tcp\_max\_syn\_backlog} \ge
    \mathtt{16*TCP\_SYNQ\_HSIZE}$, per poi ricompilare il kernel.}

\item[\sysctlrelfiled{net/ipv4}{tcp\_max\_tw\_buckets}] indica il
  numero massimo di socket in stato \texttt{TIME\_WAIT} consentito nel
  sistema. Prende un valore intero di default è impostato al doppio del valore
  del parametro \texttt{NR\_FILE}, ma che viene aggiustato automaticamente a
  seconda della memoria presente. Se il valore viene superato il socket viene
  chiuso con la stampa di un avviso; l'uso di questa funzionalità consente di
  prevenire alcuni semplici attacchi di \textit{denial of service}.
  

\item[\sysctlrelfiled{net/ipv4}{tcp\_mem}] viene usato dallo stack TCP
  per gestire le modalità con cui esso utilizzerà la memoria. Prende una
  tripletta di valori interi, che indicano un numero di pagine:

  \begin{itemize*}
  \item il primo valore, chiamato \textit{low} nelle pagine di manuale, indica
    il numero di pagine allocate sotto il quale non viene usato nessun
    meccanismo di regolazione dell'uso della memoria.

  \item il secondo valore, chiamato \textit{pressure} indica il numero di
    pagine allocate passato il quale lo stack TCP inizia a moderare il suo
    consumo di memoria; si esce da questo stato di \textsl{pressione} sulla
    memoria quando il numero di pagine scende sotto il precedente valore
    \textit{low}.

  \item il terzo valore, chiamato \textit{high} indica il numero massimo di
    pagine che possono essere utilizzate dallo stack TCP/IP, e soprassiede
    ogni altro valore specificato dagli altri limiti del kernel.
  \end{itemize*}

\item[\sysctlrelfiled{net/ipv4}{tcp\_orphan\_retries}] indica il numero
  massimo di volte che si esegue un tentativo di controllo sull'altro capo di
  una connessione che è stata già chiusa dalla nostra parte. Prende un valore
  intero che di default è 8.

\item[\sysctlrelfiled{net/ipv4}{tcp\_reordering}] indica il numero
  massimo di volte che un pacchetto può essere riordinato nel flusso di dati,
  prima che lo stack TCP assuma che è andato perso e si ponga nello stato di
  \textit{slow start} (si veda sez.~\ref{sec:tcp_protocol_xxx}) viene usata
  questa metrica di riconoscimento dei riordinamenti per evitare inutili
  ritrasmissioni provocate dal riordinamento. Prende un valore intero che di
  default che è 3, e che non è opportuno modificare.

\item[\sysctlrelfiled{net/ipv4}{tcp\_retrans\_collapse}] in caso di
  pacchetti persi durante una connessione, per ottimizzare l'uso della banda
  il kernel cerca di eseguire la ritrasmissione inviando pacchetti della
  massima dimensione possibile; in sostanza dati che in precedenza erano stati
  trasmessi su pacchetti diversi possono essere ritrasmessi riuniti su un solo
  pacchetto (o su un numero minore di pacchetti di dimensione
  maggiore). Prende un valore logico e di default è abilitato.

\item[\sysctlrelfiled{net/ipv4}{tcp\_retries1}] imposta il massimo
  numero di volte che protocollo tenterà la ritrasmissione si un pacchetto su
  una connessione stabilita prima di fare ricorso ad ulteriori sforzi che
  coinvolgano anche il livello di rete. Passato questo numero di
  ritrasmissioni verrà fatto eseguire al livello di rete un tentativo di
  aggiornamento della rotta verso la destinazione prima di eseguire ogni
  successiva ritrasmissione. Prende un valore intero che di default è 3.

\item[\sysctlrelfiled{net/ipv4}{tcp\_retries2}] imposta il numero di
  tentativi di ritrasmissione di un pacchetto inviato su una connessione già
  stabilita per il quale non si sia ricevuto una risposta di ACK (si veda
  anche quanto illustrato in sez.~\ref{sec:TCP_server_crash}). Prende un
  valore intero che di default è 15, il che comporta un tempo variabile fra 13
  e 30 minuti; questo non corrisponde a quanto richiesto
  nell'\href{http://www.ietf.org/rfc/rfc1122.txt}{RFC~1122} dove è indicato un
  massimo di 100 secondi, che però è un valore considerato troppo basso.

\item[\sysctlrelfiled{net/ipv4}{tcp\_rfc1337}] indica al kernel di
  abilitare il comportamento richiesto
  nell'\href{http://www.ietf.org/rfc/rfc1337.txt}{RFC~1337}. Prende un valore
  logico e di default è disabilitato, il che significa che alla ricezione di
  un segmento RST in stato \texttt{TIME\_WAIT} il socket viene chiuso
  immediatamente senza attendere la conclusione del periodo di
  \texttt{TIME\_WAIT}.

\item[\sysctlrelfiled{net/ipv4}{tcp\_rmem}] viene usato dallo stack TCP
  per controllare dinamicamente le dimensioni dei propri buffer di ricezione,
  anche in rapporto alla memoria disponibile.  Prende una tripletta di valori
  interi separati da spazi che indicano delle dimensioni in byte:

  \begin{itemize}
  \item il primo valore, chiamato \textit{min} nelle pagine di manuale, indica
    la dimensione minima in byte del buffer di ricezione; il default è 4Kb, ma
    in sistemi con poca memoria viene automaticamente ridotto a
    \const{PAGE\_SIZE}.  Questo valore viene usato per assicurare che anche in
    situazioni di pressione sulla memoria (vedi quanto detto per
    \sysctlrelfile{net/ipv4}{tcp\_rmem}) le allocazioni al di sotto di
    questo limite abbiamo comunque successo.  Questo valore non viene comunque
    ad incidere sulla dimensione del buffer di ricezione di un singolo socket
    dichiarata con l'opzione \const{SO\_RCVBUF}.

  \item il secondo valore, denominato \textit{default} nelle pagine di
    manuale, indica la dimensione di default, in byte, del buffer di ricezione
    di un socket TCP.  Questo valore sovrascrive il default iniziale impostato
    per tutti i socket con \sysctlfile{net/core/mem\_default} che vale per
    qualunque protocollo. Il default è 87380 byte, ridotto a 43689 per sistemi
    con poca memoria. Se si desiderano dimensioni più ampie per tutti i socket
    si può aumentare questo valore, ma se si vuole che in corrispondenza
    aumentino anche le dimensioni usate per la finestra TCP si deve abilitare
    il \textit{TCP window scaling} (di default è abilitato, vedi più avanti
    \sysctlrelfile{net/ipv4}{tcp\_window\_scaling}).

  \item il terzo valore, denominato \textit{max} nelle pagine di manuale,
    indica la dimensione massima in byte del buffer di ricezione di un socket
    TCP; il default è 174760 byte, che viene ridotto automaticamente a 87380
    per sistemi con poca memoria. Il valore non può comunque eccedere il
    limite generale per tutti i socket posto con
    \sysctlfile{net/core/rmem\_max}. Questo valore non viene ad
    incidere sulla dimensione del buffer di ricezione di un singolo socket
    dichiarata con l'opzione \const{SO\_RCVBUF}.
  \end{itemize}

\item[\sysctlrelfiled{net/ipv4}{tcp\_sack}] indica al kernel di
  utilizzare il meccanismo del \textit{TCP selective acknowledgement} definito
  nell'\href{http://www.ietf.org/rfc/rfc2018.txt}{RFC~2018}. Prende un valore
  logico e di default è abilitato.

\item[\sysctlrelfiled{net/ipv4}{tcp\_stdurg}] indica al kernel di
  utilizzare l'interpretazione che viene data
  dall'\href{http://www.ietf.org/rfc/rfc1122.txt}{RFC~1122} del puntatore dei
  \textit{dati urgenti} (vedi sez.~\ref{sec:TCP_urgent_data}) in cui questo
  punta all'ultimo byte degli stessi; se disabilitato viene usata
  l'interpretazione usata da BSD per cui esso punta al primo byte successivo.
  Prende un valore logico e di default è disabilitato, perché abilitarlo può
  dar luogo a problemi di interoperabilità.

\item[\sysctlrelfiled{net/ipv4}{tcp\_synack\_retries}] indica il numero
  massimo di volte che verrà ritrasmesso il segmento SYN/ACK nella creazione di
  una connessione (vedi sez.~\ref{sec:TCP_conn_cre}). Prende un valore intero
  ed il valore di default è 5; non si deve superare il valore massimo di 255.

\item[\sysctlrelfiled{net/ipv4}{tcp\_syncookies}] abilita i \textit{TCP
    syncookies}.\footnote{per poter usare questa funzionalità è necessario
    avere abilitato l'opzione \texttt{CONFIG\_SYN\_COOKIES} nella compilazione
    del kernel.} Prende un valore logico, e di default è disabilitato. Questa
  funzionalità serve a fornire una protezione in caso di un attacco di tipo
  \textit{SYN flood}, e deve essere utilizzato come ultima risorsa dato che
  costituisce una violazione del protocollo TCP e confligge con altre
  funzionalità come le estensioni e può causare problemi per i client ed il
  reinoltro dei pacchetti.

\item[\sysctlrelfiled{net/ipv4}{tcp\_syn\_retries}] imposta il numero di
  tentativi di ritrasmissione dei pacchetti SYN di inizio connessione del
  \textit{three way handshake} (si ricordi quanto illustrato in
  sez.~\ref{sec:TCP_func_connect}). Prende un valore intero che di default è
  5; non si deve superare il valore massimo di 255.

\item[\sysctlrelfiled{net/ipv4}{tcp\_timestamps}] abilita l'uso dei
  \textit{TCP timestamps}, come definiti
  nell'\href{http://www.ietf.org/rfc/rfc1323.txt}{RFC~1323}. Prende un valore
  logico e di default è abilitato.

\item[\sysctlrelfiled{net/ipv4}{tcp\_tw\_recycle}] abilita il riutilizzo
  rapido dei socket in stato \texttt{TIME\_WAIT}. Prende un valore logico e di
  default è disabilitato. Non è opportuno abilitare questa opzione che può
  causare problemi con il NAT.\footnote{la
    \itindex{Network~Address~Translation~(NAT)} \textit{Network Address
      Translation} (abbreviato in NAT) è una tecnica, impiegata nei firewall e
    nei router, che consente di modificare al volo gli indirizzi dei pacchetti
    che transitano per una macchina, Linux la supporta con il
    \textit{netfilter}.}

\item[\sysctlrelfiled{net/ipv4}{tcp\_tw\_reuse}] abilita il riutilizzo
  dello stato \texttt{TIME\_WAIT} quando questo è sicuro dal punto di vista
  del protocollo. Prende un valore logico e di default è disabilitato. 

\item[\sysctlrelfiled{net/ipv4}{tcp\_window\_scaling}] un valore
  logico, attivo di default, che abilita la funzionalità del
  \itindex{TCP~window~scaling} \textit{TCP window scaling} definita
  dall'\href{http://www.ietf.org/rfc/rfc1323.txt}{RFC~1323}. Prende un valore
  logico e di default è abilitato. Come accennato in
  sez.~\ref{sec:TCP_TCP_opt} i 16 bit della finestra TCP comportano un limite
  massimo di dimensione di 64Kb, ma esiste una opportuna opzione del
  protocollo che permette di applicare un fattore di scale che consente di
  aumentarne le dimensioni. Questa è pienamente supportata dallo stack TCP di
  Linux, ma se lo si disabilita la negoziazione del
  \itindex{TCP~window~scaling} \textit{TCP window scaling} con l'altro capo
  della connessione non viene effettuata.

%\item[\sysctlrelfiled{net/ipv4}{tcp\_vegas\_cong\_avoid}] 
% TODO: controllare su internet

%\item[\sysctlrelfiled{net/ipv4}{tcp\_westwood}] 
% TODO: controllare su internet

\item[\sysctlrelfiled{net/ipv4}{tcp\_wmem}] viene usato dallo stack TCP
  per controllare dinamicamente le dimensioni dei propri buffer di spedizione,
  adeguandole in rapporto alla memoria disponibile.  Prende una tripletta di
  valori interi separati da spazi che indicano delle dimensioni in byte:

  \begin{itemize}
  \item il primo valore, chiamato \textit{min}, indica la dimensione minima in
    byte del buffer di spedizione; il default è 4Kb. Come per l'analogo di
    \sysctlrelfile{net/ipv4}{tcp\_rmem}) viene usato per assicurare
    che anche in situazioni di pressione sulla memoria (vedi
    \sysctlrelfile{net/ipv4}{tcp\_mem}) le allocazioni al di sotto di
    questo limite abbiamo comunque successo.  Di nuovo questo valore non viene
    ad incidere sulla dimensione del buffer di trasmissione di un singolo
    socket dichiarata con l'opzione \const{SO\_SNDBUF}.

  \item il secondo valore, denominato \textit{default}, indica la dimensione
    di default in byte del buffer di spedizione di un socket TCP.  Questo
    valore sovrascrive il default iniziale impostato per tutti i tipi di
    socket sul file \sysctlfile{net/core/wmem\_default}. Il default è 87380
    byte, ridotto a 43689 per sistemi con poca memoria. Si può aumentare
    questo valore quando si desiderano dimensioni più ampie del buffer di
    trasmissione per i socket TCP, ma come per il precedente
    \sysctlrelfile{net/ipv4}{tcp\_rmem}) se si vuole che in corrispondenza
    aumentino anche le dimensioni usate per la finestra TCP si deve abilitare
    il \textit{TCP window scaling} con
    \sysctlrelfile{net/ipv4}{tcp\_window\_scaling}.

  \item il terzo valore, denominato \textit{max}, indica la dimensione massima
    in byte del buffer di spedizione di un socket TCP; il default è 128Kb, che
    viene ridotto automaticamente a 64Kb per sistemi con poca memoria. Il
    valore non può comunque eccedere il limite generale per tutti i socket
    posto con \sysctlfile{net/core/wmem\_max}. Questo valore non viene
    ad incidere sulla dimensione del buffer di trasmissione di un singolo
    socket dichiarata con l'opzione \const{SO\_SNDBUF}.
  \end{itemize}

\end{basedescript}



% LocalWords:  socket sez dotted decimal resolver Domain Name Service cap DNS
% LocalWords:  client fig LDAP Lightweight Access Protocol NIS Information Sun
% LocalWords:  like netgroup Switch Solaris glibc libc uclib NSS tab shadow uid
% LocalWords:  username group aliases ethers MAC address hosts networks rpc RPC
% LocalWords:  protocols services dns db lib libnss org truelite it root res HS
% LocalWords:  resource init netinet resolv int void conf host LOCALDOMAIN TCP
% LocalWords:  options DEBUG debug AAONLY USEVC UDP PRIMARY IGNTC RECURSE INET
% LocalWords:  DEFNAMES search STAYOPEN DNSRCH INSECURE NOALIASES HOSTALIASES
% LocalWords:  IPv gethostbyname NOCHECKNAME KEEPTSIG TSIG BLAST RETRY retry NS
% LocalWords:  retrans query FQDN Fully Qualified const char dname class type
% LocalWords:  unsigned answer anslen CSNET Hesiod MIT CHAOS Chaosnet ANY BIND
% LocalWords:  nameser compat Berkley MF CNAME SOA MB MR NULL WKS PTR HINFO TXT
% LocalWords:  MINFO RP responsible person AFSDB AFS RT router NSAP SIG KEY PX
% LocalWords:  GPOS AAAA LOC NXT EID NIMLOC nimrod SRV ATMA ATM NAPTR naming AF
% LocalWords:  authority IXFR AXFR MAILB MAILA errno NOT FOUND RECOVERY TRY err
% LocalWords:  AGAIN herror netdb string perror error hstrerror strerror struct
% LocalWords:  hostent name addrtype length addr list sys af mygethost inet ret
% LocalWords:  ntop deep copy buf size buflen result errnop value argument len
% LocalWords:  ERANGE sethostent stayopen endhostent gethostbyaddr order pton
% LocalWords:  getipnodebyname getipnodebyaddr flags num MAPPED ALL ADDRCONFIG
% LocalWords:  freehostent ip getXXXbyname getXXXbyaddr servent getservbyname
% LocalWords:  netent getnetbyname getnetbyaddr protoent smtp udp
% LocalWords:  getprotobyname getprotobyaddr getservbyport port tcp setservent
% LocalWords:  getservent endservent setXXXent getXXXent endXXXent gethostent
% LocalWords:  setnetent getnetent endnetent setprotoent getprotoent POSIX RFC
% LocalWords:  endprotoent getaddrinfo getnameinfo nell' node service addrinfo
% LocalWords:  hints linked addrlen socklen family socktype protocol sockaddr
% LocalWords:  canonname next PF UNSPEC SOCK STREAM DGRAM bind INADDR loopback
% LocalWords:  connect sendto NUMERICHOST EAI NONAME SYSTEM BADFLAGS ADDRFAMILY
% LocalWords:  NODATA MEMORY FAIL errcode echo mygetaddr ptr casting Canonical
% LocalWords:  freeaddrinfo getservname salen hostlen serv servlen l'OR NI NUL
% LocalWords:  NOFQDN NAMEREQD NUMERICSERV MAXHOST MAXSERV sockconn SockUtil of
% LocalWords:  descriptor hint fifth sockbind setsockopt getsockopt sock level
% LocalWords:  optname optval optlen EBADF EFAULT EINVAL ENOPROTOOPT ENOTSOCK
% LocalWords:  IPPROTO Stevens ICMP ICMPV ICMPv get KEEPALIVE OOBINLINE timeval
% LocalWords:  RCVLOWAT SNDLOWAT RCVTIMEO SNDTIMEO BSDCOMPAT BSD PASSCRED ucred
% LocalWords:  PEERCRED BINDTODEVICE REUSEADDR ACCEPTCONN DONTROUTE gateway MSG
% LocalWords:  BROADCAST broadcast SNDBUF RCVBUF LINGER linger PRIORITY read IF
% LocalWords:  OOB recvmsg kernel select write readv recv recvfrom EAGAIN send
% LocalWords:  EWOULDBLOCK writev sendmsg raw domain SCM CREDENTIALS eth packet
% LocalWords:  IFNAMSIZ capabilities capability ADMIN log trpt EADDRINUSE close
% LocalWords:  listen routing sysctl shutdown Quality TOS keep alive ACK RST to
% LocalWords:  ECONNRESET ETIMEDOUT keepalive echod fourth newsgroup WAIT reuse
% LocalWords:  sockbindopt SockUtils homed completely binding RECVDSTADDR onoff
% LocalWords:  PKTINFO getsockname multicast streaming unicast REUSEPORT reset
% LocalWords:  stealing ling RECVTOS RECVTTL TTL RECVOPTS RETOPTS HDRINCL MTU
% LocalWords:  RECVERR DISCOVER Path Discovery ALERT alert ADD MEMBERSHIP mreqn
% LocalWords:  pktinfo ipi ifindex spec dst RECVIF Live IPTOS LOWDELAY Advanced
% LocalWords:  Transfer Unit PMTUDISC DONT WANT route dall' pmtu EMSGSIZE imr
% LocalWords:  multiaddr mreq fcntl ioctl request SIOCGSTAMP trip SIOCSPGRP pid
% LocalWords:  process SIGIO SIGURG KILL FIOASYNC SIOCGPGRP filesystem proc ttl
% LocalWords:  rmem wmem message cost burst bucket filter netdev backlog optmem
% LocalWords:  forward dynaddr dial autoconfig local masquerading ipfrag high
% LocalWords:  thresh low always defrag CONFIG SETSIG cmd FIOGETOWN FIOSETOWN
% LocalWords:  quest'ultime neigh dev weight cong mod somaxconn Di SIOCINQ DoS
% LocalWords:  Documentation SIOCATMARK SIOCOUTQ FIONREAD TIOCOUTQ Denial work
% LocalWords:  netfilter scheduler mark ARP DHCP BOOTP RARP nonlocal sniffer is
% LocalWords:  linux NODELAY MAXSEG CORK KEEPIDLE KEEPINTVL KEEPCNT SYNCNT INFO
% LocalWords:  DEFER ACCEPT WINDOW CLAMP QUICKACK CONGESTION ENCAP urgent MSS
% LocalWords:  Segment SYN accept advertised window info quickack Nagle ifreq
% LocalWords:  ifr ppp union EPERM SIOCGIFNAME dell' interface index IFF NOARP
% LocalWords:  SIOCGIFINDEX SIOCGIFFLAGS POINTOPOINT RUNNING PROMISC NOTRAILERS
% LocalWords:  ALLMULTI bundle PORTSEL ifmap AUTOMEDIA DYNAMIC SIOCSIFFLAGS way
% LocalWords:  SIOCGIFMETRIC SIOCSIFMETRIC SIOCGIFMTU SIOCSIFMTU SIOCGIFHWADDR
% LocalWords:  SIOCSIFHWADDR SIOCSIFHWBROADCAST SIOCGIFMAP SIOCSIFMAP sendfile
% LocalWords:  SIOCADDMULTI SIOCDELMULTI SIOCGIFTXQLEN SIOCSIFTXQLEN three syn
% LocalWords:  SIOCSIFNAME SIOCGIFCONF handshake retries MIN FreeBSD closing Mb
% LocalWords:  abort overflow adv win app bic convergence dsack ecn fack frto
% LocalWords:  intvl probes latency orphans l'ACK SYNQ HSIZE tw buckets mem rfc
% LocalWords:  orphan reordering collapse sack stdurg synack syncookies recycle
% LocalWords:  timestamps scaling vegas avoid westwood tcpi l'incapsulazione NR
% LocalWords:  metric EOPNOTSUPP mtu hwaddr ARPHRD interrupt DMA map qlen silly
% LocalWords:  rename ifconf syndrome dell'ACK FTP ACCEPTFILTER advanced reno
% LocalWords:  congestion control Networking cubic CUBIC highspeed HSTCP htcp
% LocalWords:  HTCP hybla HYBLA scalable SCALABLE ifc req iflist access ntoa Kb
% LocalWords:  hop Selective acknowledgement Explicit RTO stack firewall passwd
% LocalWords:  Notification wireless denial pressure ATTACH DETACH publickey
% LocalWords:  libpcap discovery point l'overhaed min PAGE flood NFS blast
% LocalWords:  selective COOKIES NAT Translation

%%% Local Variables: 
%%% mode: latex
%%% TeX-master: "gapil"
%%% End: 

